% Korean Noether paper 02 — mechanically assembled inherited production body.
% Origin reopened against ED0005; current scope-rebase authority: ED0011 / NOETH-DE-AUTH-v047-20260808.

% inherited: noether_paper02_ko_translation_001_20260804\targets\Noether_P02_Korean_T01_U01_UNCHECKED.tex
% sha256: CAD11967D6497AEEE5FE1C37AC016A954BD40A413400CFE3D9E33B3FE78CC78C
\begin{center}
{\Large\bfseries 2. 삼원 4차 형식의 형식계 구성에 관하여\srcfn{*)}{서론과 제I장의 짧은 발췌문이 \glqq{}Sitzungsber. der phys.-med. Sozietät Erlangen 1907\grqq{}, 176--179쪽에 발표되었다.}.}\par
\vspace{0.8em}
에를랑겐의 에미 뇌터 양.\par
\vspace{1.2em}
Journal f. d. reine u. angew. Math. 134 (1908), 23--90쪽 및 표 2개
\end{center}

\vspace{3em}

% inherited: noether_paper02_ko_translation_001_20260804\targets\Noether_P02_Korean_T01_U02_UNCHECKED_TEXREPAIR_v002_20260804.tex
% sha256: 1D39A01D8908BBEBF41BDE6D6A9E691756BE57CC493BFEE59337371FAEEB8BD7
\begin{center}
\textbf{서론.}
\end{center}

삼원 4차 형식의 형식계는 \emph{Gordan}, \emph{Maisano}, \emph{Pascal}의 연구에서 다루어졌다\srcfn{**)}{\emph{P. Gordan}, Über das volle Formensystem der ternären biquadratischen Form
\[
f=x_1^3x_2+x_2^3x_3+x_3^3x_1.
\]
(Math. Annalen, Bd. XVII (1880), S. 217--233.)
\emph{G. Maisano}, 1) \emph{Sistemi completi dei primi cinque gradi della forma ternaria biquadratica e degl' invarianti, covarianti e contravarianti di sesto grado}. (Giorn. di Battaglini XIX.)
2) \emph{Sui covarianti indipendenti di $6^\circ$ grado nei coefficienti della forma biquadratica ternaria}. (Rend. Circ. Mat. di Palermo I, 1887.)
\emph{E. Pascal}, \emph{Contributo alla teoria della forma ternaria biquadratica e delle sue varie decomposizioni in fattori}. (Memoria premiata dalla R. Accademia delle scienze fisiche e matematiche di Napoli. 1905.)}. \emph{Gordan}은 이항 영역의 형식계에 사용했던 것과 비슷한 원리를 바탕으로, 특수한 자기동형 형식 \(f=x_1^3x_2+x_2^3x_3+x_3^3x_1\)의 54개 구성식으로 이루어진 완전한 형식계를 세웠다.

\emph{Maisano}는 일반 4차 형식에 대해 계수차수 5까지의 형식들\srcfn{***)}{여기서 \glqq{}Ordnung\grqq{}은 계수에 관한 차수를, \glqq{}Grad\grqq{}는 변수에 관한 차수를 뜻한다.}과 더 높은 계수차수의 몇몇 불변식, 공변식, 반공변식을 제시하였다. 그 방법은 \emph{Gordan}이 Math. Annalen 제1권에서 삼원 3차 형식에 적용한 것이다. \emph{Pascal}은 \emph{Maisano}의 결과를 이용하여 주로 4차 형식의 인수분해 문제를 다루었다\srcfn{*)}{두 번째 짧은 논문에서 \emph{Maisano}는 계수차수 6, 변수차수 6인 세 공변식의 선형 종속성을 보이려 하였다. 완전하지 못한 이 증명과 달리 \emph{Pascal}은 계수차수 6인 세 공변식의 선형 독립성과 계수차수 9인 세 불변식의 선형 독립성을 증명했다고 여겼다. 그러나 실제로 한편의 세 공변식 사이에, 다른 한편의 세 불변식 사이에 각각 선형 관계가 존재한다는 사실은 이 논문의 공식 (13), § 11과 공식 (22), § 17 주석에서 확인되며, 그곳에 이 관계들이 명시적으로 주어져 있다. \emph{Pascal}의 통지에 따르면 이 모순은 그의 논문 46쪽에 있는 반공변식 $p$(여기서는 $\varrho$라 부른다)의 식에 생긴 수치 오류로 설명된다.}.

% inherited: noether_paper02_ko_translation_001_20260804\targets\Noether_P02_Korean_T01_U03_UNCHECKED.tex
% sha256: 1C31C069199BA2C36A146142D3AE29990D8D6BF9BB19C4F3FF6A487794D60E46
\emph{이하 연구의 목적은 일반 삼원 4차 형식의 형식계를 세우는 것이다. 이 논문에서는 그 주된 기초만을 제시하고, 이른바 \glqq{}상대 완전계\grqq{}}\srcfn{**)}{명칭에 대해서는 § 3a를 참조하라.}\emph{를 구성한다.}

이 논문은 \emph{Gordan}의 연구와 밀접하게 이어져 있지만, 그곳에서 아주 일반적으로만 제시된 원리들을 먼저 세부적으로 전개해야 했다. 수축들 사이의 관계에 관한 한 정리를 이용하고 \glqq{}형식열\grqq{}(§ 1)을 도입함으로써 환원 정리를 엄밀하게 정식화하고 완성하며(§ 3), 특수한 급수 전개를 재귀적으로 구성하는 일(§ 2)은 환원을 실제로 수행하기 위한 계산 수단을 제공한다.

삼원 형식계를 구성하는 기본 발상은 이항 영역에서와 같다. 첫 번째 상대 완전계, 곧 이항 영역에서 넘겨받은 형식들의 계에서 출발하여 일정한 법칙에 따라 모듈이 점점 더 높은 계들로 나아간다. 이 과정은 어떤 모듈의 계가 --- 그 모듈을 기본형으로 삼았을 때 --- 마침내 유한하고 알려진 것이 될 때까지, 또는 한 모듈을 불변식을 인수로 갖는 형식들로 환원할 수 있을 때까지 계속된다. 첫째 경우에는 상대 완전계를 그 모듈의 계와 전이연산하여 절대 완전계를 얻고, 둘째 경우에는 상대 완전계와 절대 완전계가 일치한다. \emph{Hilbert}가 형식계의 유한성을 일반적으로 증명했으므로, 이 절차는 반드시 유한 단계에서 끝나야 한다.

% inherited: noether_paper02_ko_translation_001_20260804\targets\Noether_P02_Korean_T02_U04_UNCHECKED.tex
% sha256: B43F6E583C9EF9E19F7C98F6CA2944C3F26C94E076BFEC672A8CFA63EB6E8274
우리의 경우 첫 번째 상대 완전계의 모듈 $(abc)$(§ 4)은 다음 모듈들로 환원된다.
\[
\Delta=(abc)^2a_x^2b_x^2c_x^2\quad\text{및}\quad \nu=(abu)^4
\]
(§ 5). 이어서 $\nu$를 법으로 하는 상대 완전계를 구한다(§ 6). 이 결과들은 일반 4차 형식에 대해 이미 \emph{Gordan}의 연구에 제시되어 있지만, 여기의 유도 방식은 더 높은 차수의 형식으로 더 쉽게 옮길 수 있다.

% inherited: noether_paper02_ko_translation_001_20260804\targets\Noether_P02_Korean_T02_U05_UNCHECKED.tex
% sha256: 51FCF8F3308E9506CFF753A98A7E81F5C474E9214AA723ED370ED26C7207F917
이제 모듈열로 다음 형식들을 택한다(§ 7).
\[
\nu,\qquad \nu(\nu)=(\nu\nu_1x)^4=s_x^4,\qquad \nu(s)=(ss'u)^4,\ldots .
\]

$s$를 법으로 하는 상대 완전계를 구성할 때, 그 계에 나타나는 두 이차형식 $u_\varrho^2$와 $t_x^2$를 모듈로 첨가한다(§ 9와 10). 따라서 $(s,\varrho,t)$를 법으로 하는 상대 완전계를 구성한다. 이어서 (§ 17) 모듈 $(ss'u)^4$가 모듈 $\varrho$와 $t$로 환원 가능함을 보인다. 이로써 그다음의 더 높은 계는 $(\varrho,t)$를 법으로 하는 상대 완전계로 넘어간다. 그런데 두 이차형식의 연립계는 유한하고 알려져 있으며, 일반적인 경우 20개의 구성식으로 이루어진다\srcfn{*)}{\emph{Clebsch-Lindemann}, Vorlesungen über Geometrie. (Bd. I, S. 288 ff.) 서로 동변인 두 형식의 계. \emph{Gordan}, Über Büschel von Kegelschnitten. (Math. Ann. Bd. XIX. S. 530). 서로 반변인 두 형식의 계.}. 그러므로 $(\varrho,t)$를 법으로 하는 상대 완전계를 세움으로써 위에서 말한 종결점에 이른다. 절대 완전계를 만들기 위하여 이 계를 $(\varrho,t)$의 계와 전이연산하는 일은 이후의 과제로 남겨 둔다.

$(\varrho,t)$를 법으로 하는 상대 완전계로 331개의 구성식을 얻으며, 첨부된 표에서는 이를 변수 $x$와 $u$에 관한 차수에 따라 배열하였다.

% inherited: noether_paper02_ko_translation_001_20260804\targets\Noether_P02_Korean_T02_U06_UNCHECKED.tex
% sha256: 91D1AEDE82C311ABCD7F5F0D6B159F340DD9270CD48F4ABB4CCA00484B97AB2B
끝으로 도입한 기호들을 한눈에 볼 수 있도록 정리한다.
\begin{align*}
f&=a_x^4,\\
\theta&=\theta_x^4u_\vartheta^2=(abu)^2a_x^2b_x^2,\\
K&=K_x^6u_k^3=(a\theta u)u_\vartheta^2a_x^3\theta_x^3,\\
j&=u_j^6=(a\theta u)^4u_\vartheta^2,\\
\Delta&=\Delta_x^6=a_\vartheta^2a_x^2\theta_x^4=(abc)^2a_x^2b_x^2c_x^2,\\
N&=N_x^8u_\eta=(a\Delta u)a_x^3\Delta_x^5,\\
\nu&=u_\nu^4=(abu)^4,\\
H&=H_x^2u_\eta^4=(\nu\nu_1x)^2u_\nu^2u_{\nu_1}^2,\\
L&=L_x^3u_l^6=(\nu\eta x)H_x^2u_\nu^3u_\eta^3,\\
g&=g_x^6=(\nu\eta x)^4H_x^2,\\
\sigma&=u_\sigma^6=H_\nu^2u_\nu^2u_\eta^4,\\
s&=s_x^4=(\nu\nu_1x)^4,\\
Z&=Z_x^4u_\zeta^2=(ss'u)^2s_x^2s_x'^2,\\
\varrho&=u_\varrho^2=\theta_\nu^4u_\vartheta^2,\\
t&=t_x^2=a_\eta^4H_x^2,\\
i&=a_\nu^4,\\
J&=s_\nu^4.
\end{align*}

% inherited: noether_paper02_ko_translation_001_20260804\targets\Noether_P02_Korean_T03_U07_UNCHECKED.tex
% sha256: 6DC2898B793F958C194AAC22D5FAF23B00CAFA6E47359C3FF45878CB6233F4CD
\begin{center}
{\Large\bfseries 제I장. \quad 삼원 형식에 관한 일반 정리.}
\end{center}

\begin{center}
\textbf{§ 1.}\quad \emph{수축 과정. 형식열.}
\end{center}

형식을 생성하는 기본 과정은 \emph{수축 과정}이며, 삼원 영역에서는 다음과 같이 정의할 수 있다. 기호곱
\[
s_x^m t_x^n u_\sigma^\mu u_\tau^\nu
\]
이 주어졌다고 하자. 인수쌍
\[
s_xt_x\qquad u_\sigma u_\tau\qquad s_xu_\sigma\ \text{또는}\ s_xu_\tau
\qquad t_xu_\tau\ \text{또는}\ t_xu_\sigma
\]
을 각각
\[
(stu)\qquad (\sigma\tau x)\qquad s_\sigma\ \text{또는}\ s_\tau
\qquad t_\tau\ \text{또는}\ t_\sigma
\]
로 바꾸는 것을 수축
\[
\text{I}\qquad\text{II}\qquad\text{III}\qquad\text{IV}
\]
라 하면, 이렇게 생기는 형식들은 원래 형식에서 수축으로 얻어진 것이다\srcfn{*)}{\emph{Gordan}, Math. Annalen Bd. XVII S. 219.}.

% inherited: noether_paper02_ko_translation_001_20260804\targets\Noether_P02_Korean_T03_U08_UNCHECKED.tex
% sha256: 03E42551EC74FCA965680A680CAC0D45B2B482E4C78D6658A41E98753FA20722
이때 식 $s_x^m t_x^n u_\sigma^\mu u_\tau^\nu$는 두 형식의 실제 곱 $S\cdot T$로 볼 수도 있고, 여러 기호열을 가진 하나의 형식으로 볼 수도 있다.
\[
s_x^m t_x^n u_\sigma^\mu u_\tau^\nu=A_x^{m+n}u_\lambda^{\mu+\nu}.
\]
첫째 경우에는 형식 $S$와 $T$의 수축이라 하고, 둘째 경우에는 \glqq{}형식 $A$의 자기 수축\grqq{}이라 한다. 간결함을 위해 이하에서는 언제나 다음을 가정한다. 이는 실제로 뒤에서 살필 모든 형식에 성립한다.
\srcnumdisplay{(a)}{\begin{gathered}
s_\sigma^\lambda=0,\\
t_\tau^\varkappa=0
\end{gathered}}
여기서 0이 아닌 임의의 $\lambda$와 $\varkappa$를 취하고, 이 식들을 임의의 다른 수축들과 결합한다. 이는 형식 $s_x^m t_y^n u_\sigma^\mu u_\tau^\nu$가 이른바 \glqq{}정규형\grqq{}임을 뜻한다. 즉 변수 $x,u$에 관해서든 변수 $y,v$에 관해서든 $\Omega$-과정을 적용하면 소멸한다.

따라서 조건 (a)는 아무런 제한도 주지 않으며 언제나 달성할 수 있다\srcfn{*)}{\emph{Gordan}, Math. Annalen Bd. V S. 104 참조.}.

개별 수축들 사이의 관계에 대해서는 이제 다음 정리가 성립한다.

\emph{정리 I. 수축 I과 II는 기본 수축이며, 합성 순서와 무관하게 이들로부터 수축 III과 IV를 합성할 수 있다. 다시 말해, 주어진 식에서 수축으로 생기는 모든 형식을 만들기 위해서는 수축 I과 II만 적용하면 된다.}

% inherited: noether_paper02_ko_translation_001_20260804\targets\Noether_P02_Korean_T03_U09_UNCHECKED.tex
% sha256: D0E8FC772D8D352F3477D857EDDCAE7678975D8328E2268EB08AD1D2D793E34E
\textbf{증명:} 항등식 정리와 행렬의 곱정리를 각각 사용하고 (a)를 고려하면 다음이 따른다.
\[
(\widehat{st}\sigma x)=-t_\sigma,\qquad
(\widehat{\sigma\tau}su)=-s_\tau,
\]
\[
(\widehat{st}\tau x)=s_\tau,\qquad
(\widehat{\sigma\tau}tu)=t_\sigma,
\]
\[
(stu)(\sigma\tau x)=s_\tau+t_\sigma-u_x\cdot s_\tau t_\sigma .
\]
따라서 수축 I과 II를 합성하면 수축 III과 IV가 생긴다. 이는 수축 I, 곧 인수 $(stu)u_\sigma u_\tau$에 수축 II를 적용할 때에도, 수축 II, 곧 인수 $(\sigma\tau x)s_xt_x$에 수축 I을 적용할 때에도 마찬가지이다.

% inherited: noether_paper02_ko_translation_001_20260804\targets\P02_T03_U10_r2.tex
% sha256: 79DFF71BBB82EAB95B85783C28760B777C69B01EE8A4CB0DDF921DD0961DDA06
\emph{형식열}\srcfn{**)}{\emph{Clebsch}, Über eine Fundamentalaufgabe der Invariantentheorie (Abh. der Gött. Ges. d. Wiss. Bd. XVII), § 17 및 § 18 끝 참조. \glqq{}형식열\grqq{}은 더 높은 형식에 따른 배열 원리에 의해 그곳에서 도입한 \glqq{}본래적으로 환원된 동치계\grqq{}와 구별된다.}\setcounter{footnote}{9}을 하나의 시작형식과 이 시작형식의 자기 수축으로 생긴 모든 형식으로 이루어진 것으로 정의하고, 그 시작형식으로 형식열을 나타낸다. 여기서 자기 수축이 더 많은 형식을 \glqq{}더 높은 형식\grqq{}이라 하며, 서로 동등한 수축 $s_\tau$와 $t_\sigma$ 가운데 하나를 더 높은 것으로 정규화해야 한다. 형식열의 구성 법칙에서 다음이 따른다.

1) 형식열의 각 형식은 시작형식의 기호들에 관해 선형이다.

2) 서로 반변인 기호열을 각각 두 개씩 가진 시작형식에서는 형식열이 유일하게 정해진다. 더 많은 기호열을 가진 시작형식에서는 항등식 정리로 연결되는 수축들 가운데 특정 수축들을 우선시함으로써 형식열을 유일하게 정규화해야 한다.

% inherited: noether_paper02_ko_translation_001_20260804\targets\P02_T03_U11_r2.tex
% sha256: 7D37668528A5337CC743E073A62AED0C82EE81D55595A21C49C5F882EB4B3333
정리 I에 따라 형식열 $s_x^m t_x^n u_\sigma^\mu u_\tau^\nu$ ($m>n;\ \mu>\nu$)는 다음 직사각형 도식으로 배열할 수 있다.
\begingroup\small
\[
\begin{array}{c|c|c|c|c}
s_x^m t_x^n u_\sigma^\mu u_\tau^\nu
& (stu)
& (stu)^2
& \cdots
& \bigstar\\[0.3em]
(\sigma\tau x)
& t_\sigma,\ s_\tau
& t_\sigma(stu),\ s_\tau(stu)
& \cdots
& \vdots\\[0.3em]
(\sigma\tau x)^2
& t_\sigma(\sigma\tau x),\ s_\tau(\sigma\tau x)
& t_\sigma^2,\ t_\sigma s_\tau,\ s_\tau^2
& \cdots
& \vdots\\[0.3em]
\vdots&\vdots&\vdots&\ddots&\vdots\\[0.3em]
(\sigma\tau x)^\nu
& t_\sigma(\sigma\tau x)^{\nu-1},\ s_\tau(\sigma\tau x)^{\nu-1}
& t_\sigma^2(\sigma\tau x)^{\nu-2},\ t_\sigma s_\tau(\sigma\tau x)^{\nu-2},\ s_\tau^2(\sigma\tau x)^{\nu-2}
& \cdots
& \vdots\\[0.3em]
\vdots
& t_\sigma(\sigma\tau x)^\nu
& t_\sigma^2(\sigma\tau x)^{\nu-1},\ t_\sigma s_\tau(\sigma\tau x)^{\nu-1}
& \cdots
& \vdots\\[0.3em]
\vdots
& \vdots
& t_\sigma^2(\sigma\tau x)^\nu
& \cdots
& \vdots
\end{array}
\]
\endgroup
\begingroup\small
\[
\begin{array}{r@{\quad}c}
\bigstar\srcfnmark{*)}
&
\begin{array}{c|c|c}
(stu)^n & \cdots & \cdots\\[0.3em]
t_\sigma(stu)^{n-1},\ s_\tau(stu)^{n-1}
& s_\tau(stu)^n & \cdots\\[0.3em]
t_\sigma^2(stu)^{n-2},\ t_\sigma s_\tau(stu)^{n-2},\ s_\tau^2(stu)^{n-2}
& t_\sigma s_\tau(stu)^{n-1},\ s_\tau^2(stu)^{n-1}
& s_\tau^2(stu)^n .
\end{array}
\end{array}
\]
\endgroup
\srcfntext{*)}{여기서 그리고 이하의 비슷한 곳에서, 별표가 붙은 아래쪽 부분은 도식의 오른쪽 위에 같은 별표로 표시한 자리에 이어 붙여야 한다.}

% inherited: noether_paper02_ko_translation_001_20260804\targets\Noether_P02_Korean_T03_U12_UNCHECKED.tex
% sha256: C137EFA41218012899DF7DB19C3F9B5368A39573DE35D6E70A1050F44F0BED1B
이 도식에서

1) 한 열씩 오른쪽으로 나아가면 바로 왼쪽의 형식들에 수축 I을 적용하여 생기는 모든 형식을 얻고,

2) 한 행씩 아래로 나아가면 바로 위의 형식들에 수축 II를 적용하여 생기는 모든 형식을 얻는다.

따라서 수축으로 생기는 모든 형식을 얻는다.

전체 형식열의 임의의 형식 $t_\sigma^\kappa s_\tau^\lambda(stu)^\mu$ 또는 $t_\sigma^\kappa s_\tau^\lambda(\sigma\tau x)^\nu$는 새로운 형식열의 시작을 이룬다. 이 새 형식열은 직사각형 도식에서 해당 형식 $t_\sigma^\kappa s_\tau^\lambda(stu)^\mu$ 또는 $t_\sigma^\kappa s_\tau^\lambda(\sigma\tau x)^\nu$를 왼쪽 위 꼭짓점으로 삼고, 인수 $t_\sigma^\kappa s_\tau^\lambda$를 갖는 모든 형식으로 이루어진다.

% inherited: noether_paper02_ko_translation_001_20260804\targets\Noether_P02_Korean_T03_U13_UNCHECKED.tex
% sha256: 9E66EEE7BD1412E5C251919D73D35EFCA5BE2411C794E53EC13769256EFD1D8C
후기. 수 $n$과 $\mu$(각각 $m$과 $\nu$)가 0이 되는 경우에는 정리 I에 예외가 생긴다. 이때 수축 I, II 및 IV는 더 이상 존재하지 않지만 수축 III(각각 IV)는 여전히 존재한다. 이 경우에는 다음 도식의 대각 성분을 형식열이라 한다.
\[
\begin{array}{c@{\qquad}c@{\qquad}c}
\begin{array}{ccccc}
s_x^m u_\tau^\nu &&&&\\
&s_\tau&&&\\
&&s_\tau^2&&\\
&&&\ddots&\\
&&&&s_\tau^\nu
\end{array}
&\text{각각}&
\begin{array}{ccccc}
t_x^n u_\sigma^\mu &&&&\\
&t_\sigma&&&\\
&&t_\sigma^2&&\\
&&&\ddots&\\
&&&&t_\sigma^\mu
\end{array}
\end{array}
\]
수축 I과 II가 존재하지 않는 것과 일치하게, 이 경우 형식열의 편극에 따른 급수 전개는 언제나 시작 성분으로 환원된다. 그러나 환원자에 관한 정리들은 성립한다.

% inherited: noether_paper02_ko_translation_001_20260804\targets\Noether_P02_Korean_T04_U14_UNCHECKED.tex
% sha256: E4EA45CC6E9CF43F71B22B84EBB75D545679D588D158008BF2EA1CB99B9D4CCF
\begin{center}
\textbf{§ 2.}\quad \emph{형식열의 편극에 따른 급수 전개}\srcfn{*)}{Study, Methoden zur Theorie der ternären Formen (Teubner 1889) II. § 7 참조. 그곳의 유도와 달리 우리의 유도는 수치 계수 $C_{pr\varrho}$를 빠르게 계산하여 정하는 방법을 제공한다. 특수화 a), a')에서 Clebsch--Study 전개는 다음과 같을 것이다.
\[
(stu)^\mu u_\tau^{\nu-\varkappa}v_\tau^\varkappa s_x^m t_x^{\,n-\lambda}(tuv)^\lambda
=\sum_{p,r,\varrho} C'_{pr\varrho}
\bigl[s_\tau^\varrho(stu)^{\mu+r-\varrho}\bigr]_{\widehat{uv}^{\lambda-r+\varrho}v^{\varkappa-\varrho},\,w^p,\,(w=\widehat{xw})}.
\]
따라서 인수 $u_x$를 가진 모든 항을 생략해야 한다는 사실을 알고 있을 때 우리의 전개에서는 계산 도중에 이미 나타나는 단순화를 그 전개에서는 할 수 없다.}.
\end{center}

$n$개 변수를 가진 형식에 대해서는 두 종류의 급수 전개가 명시적으로 주어져 있다\srcfn{**)}{Gordan, Über Kombinanten, § 2와 5 (Math. Ann. Bd. V) 참조.}.

1) 서로 반변인 변수열을 하나씩 가진 형식 $s_x^m u_\tau^\nu$에 대해서는 $u_x$의 거듭제곱에 따라 진행하는 급수이며, 그 계수들은 \glqq{}정규형\grqq{}이다.

2) 서로 동변인 변수열 두 개를 가진 형식 $s_x^m t_x^n$에 대해서는 기본 공변식의 편극(형식열 $s_x^m t_x^n$의 편극)에 따른 전개이며, 삼원 영역에서는 다음과 같다.
\[
\tag{I}
s_x^m t_x^{n-\lambda}t_y^\lambda
=\sum_{\varrho}
\frac{\binom{m}{\varrho}\binom{\lambda}{\varrho}}
 {\binom{m+n-\varrho+1}{\varrho}}\cdot
\bigl[(st\widehat{xy})^\varrho s_x^{m-\varrho}t_x^{n-\varrho}\bigr]_{y^{\lambda-\varrho}}.
\]

% inherited: noether_paper02_ko_translation_001_20260804\targets\Noether_P02_Korean_T04_U15_UNCHECKED.tex
% sha256: DD4A438DAFC2D3779DE165B8C86F2157C43A0C72C0BF0EC21DFE0589A8488251
서로 반변인 변수열을 두 개씩 가진 형식
\[
s_x^m t_x^{n-\lambda}t_y^\lambda u_\sigma^\mu u_\tau^{\nu-\varkappa}v_\tau^\varkappa
\]
에 대해 두 전개를 결합한다. 즉 $u_x$의 거듭제곱에 따른 전개를 세우되, 그 계수는 변수 $x$와 $u$에 관해 취한 형식열 $s_x^m t_x^n u_\sigma^\mu u_\tau^\nu$의 합성 편극이 되게 한다.

서로 반변인 기호열(각각 변수열)을 두 개씩 가진 경우의 급수만 세우면 충분하다. 기호열(변수열)을 두 개씩 하나의 새로운 열로 묶으면, 더 많은 기호열(변수열)을 가진 형식들을 이 경우로 환원할 수 있기 때문이다.

형식열의 모든 형식이 기호열과 변수열을 각각 두 개만 포함하게 하기 위해 다음과 같이 특수화한다.
\[
\begin{array}{ll}
\text{a) } \sigma=\widehat{st}, & \text{a') } y=\widehat{uv},\\[0.3em]
\text{b) } s=\widehat{\sigma\tau}, & \text{b') } v=\widehat{xy},
\end{array}
\]
그리고 특수화 a), a')에 대해 전개를 수행한다.

% inherited: noether_paper02_ko_translation_001_20260804\targets\Noether_P02_Korean_T04_U16_UNCHECKED.tex
% sha256: 3373680A739CF0DAA4F577DEF0F2EB9C3D258A331AC365AD68EB548D27E89267
전개 I을 직접 옮기면 형식 $(stu)^\mu s_x^m t_x^{n-\lambda}t_y^\lambda$에 대해서는 합성 편극을 얻는다. 그러나 형식 $(stu)^\mu u_\tau^{\nu-\varkappa}v_\tau^\varkappa s_x^m t_x^{n-\lambda}t_y^\lambda$에 대해서는 합성 편극 대신 그러한 편극들의 항들이 생긴다. 이 항들을 계산하려면 계속 새로운 보조변수를 도입하여 마지막에야 서로 같게 놓아야 하므로 계산이 불필요하게 복잡해진다.

따라서 우리는 간접적인 길을 택한다. 형식열의 모든 형식의 합성 편극, 곧 식
\[
\bigl[s_\tau^\varrho(stu)^{\mu-\varrho+r}\bigr]_{y^\lambda}^{v^\varkappa}
\]
을 이 편극들의 시작항들로 나타내고, 그 결과 생기는 재귀 방정식계를 역으로 풀어 원하는 편극에 따른 전개를 얻는다.

% inherited: noether_paper02_ko_translation_001_20260804\targets\Noether_P02_Korean_T04_U17_UNCHECKED.tex
% sha256: A984E0FCC41245E5B92E68E65FD45920862D90E858A0FAB2F773372B641E7DE4
전이 원리에 따라 형식 $s_x^m t_x^n$의 $\lambda$번째 편극 $[s_x^m t_x^n]_{y^\lambda}$ ($\lambda\le n$)에는 다음 전개가 성립한다\srcfn{*)}{Gordan, Die Resultante binärer Formen (Kap. I, § 5). Rend. Circ. Mat. di Palermo XXII. 1906.}. $\lambda>n$인 경우에는 $[s_y^m t_y^n]_{x^{m+n-\lambda}}$의 전개를 통하여 첫째 경우로 돌아간다.
\[
\bigl[s_x^m t_x^n\bigr]_{y^\lambda}
=\sum_r (-1)^r\,
\frac{\binom{m}{r}\binom{\lambda}{r}}{\binom{m+n}{r}}\cdot
(st\widehat{xy})^r s_x^{m-r}t_x^{n-\lambda}t_y^{\lambda-r}.
\]
이에 대응하여 $u_\sigma^\mu u_\tau^\nu$의 $\varkappa$번째 편극 $[u_\sigma^\mu u_\tau^\nu]^{v^\varkappa}$에는 다음이 성립한다.
\[
\bigl[u_\sigma^\mu u_\tau^\nu\bigr]^{v^\varkappa}
=\sum_\varrho (-1)^\varrho\,
\frac{\binom{\mu}{\varrho}\binom{\varkappa}{\varrho}}{\binom{\mu+\nu}{\varrho}}\,
(\sigma\tau\widehat{uv})^\varrho u_\sigma^{\mu-\varrho}u_\tau^{\nu-\varkappa}v_\tau^{\varkappa-\varrho}.
\]

% inherited: noether_paper02_ko_translation_001_20260804\targets\Noether_P02_Korean_T04_U18_UNCHECKED.tex
% sha256: AC9E075C49B82565858C88CB2D75A170344A43B64B5FE5C8EB6A537F0BE27149
두 식을 곱하고 특수화 a)와 a')를 도입하면 다음을 얻는다.
\[
\begin{aligned}
\bigl[s_x^m t_x^n(stu)^\mu u_\tau^\nu\bigr]_{\widehat{uv}^{\lambda}}^{v^\varkappa}
={}&
\sum_{r,\varrho}c_{r\varrho}\cdot
(stx\widehat{uv})^r(\widehat{st}\tau\widehat{uv})^\varrho
s_x^{m-r}t_x^{n-\lambda}(tuv)^{\lambda-r}(stu)^{\mu-\varrho}
u_\tau^{\nu-\varkappa}v_\tau^{\varkappa-\varrho}.
\end{aligned}
\]
여기서 $u_x$의 거듭제곱에 따라 전개하면서 첫째 항을 따로 떼어 낸다. $\sum'_{r,\varrho}$는 항 $r=0$, $\varrho=0$을 생략한다는 뜻이다. 또 26쪽의 (a)를 고려하면 다음을 얻는다.
\[
\tag{II}
\begin{aligned}
&s_x^m t_x^{n-\lambda}(tuv)^\lambda(stu)^\mu u_\tau^{\nu-\varkappa}v_\tau^\varkappa\\
&\quad =
\bigl[s_x^m t_x^n(stu)^\mu u_\tau^\nu\bigr]_{\widehat{uv}^{\lambda}}^{v^\varkappa}\\
&\qquad
-\sum_{r,\varrho}' c_{r\varrho}\cdot s_\tau^\varrho(stu)^{\mu-\varrho+r}
u_\tau^{\nu-\varkappa}v_\tau^{\varkappa-\varrho}
s_x^{m-r}t_x^{n-\lambda}(tuv)^{\lambda-r+\varrho}\cdot v_x^r\\
&\qquad
+u_x\cdot\sum_{\varrho}\cdot\sum_{r=1}^{\lambda} c_{r\varrho}\cdot
s_\tau^\varrho(stu)^{\mu-\varrho+r-1}(stv)
u_\tau^{\nu-\varkappa}v_\tau^{\varkappa-\varrho}
s_x^{m-r}t_x^{n-\lambda}(tuv)^{\lambda-r+\varrho}\cdot v_x^{r-1}\\
&\qquad
\vdots\\
&\qquad
-(-1)^\lambda u_x^\lambda\cdot\sum_\varrho c_{\lambda\varrho}\cdot
s_\tau^\varrho(stu)^{\mu-\varrho}(stv)^\lambda
u_\tau^{\nu-\varkappa}v_\tau^{\varkappa-\varrho}
s_x^{m-\lambda}t_x^{n-\lambda}(tuv)^\varrho,
\end{aligned}
\]
여기서
\[
c_{r\varrho}=(-1)^{r+\varrho}\cdot
\frac{\binom{m}{r}\binom{\lambda}{r}}{\binom{m+n}{r}}\cdot
\frac{\binom{\mu}{\varrho}\binom{\varkappa}{\varrho}}{\binom{\mu+\nu}{\varrho}}
\qquad
\begin{array}{l}
\lambda\le n,\\
\varkappa\le\nu
\end{array}
\]
이고, 수 $\lambda,n;\varkappa,\nu$의 나머지 배열들에 대해서도 이에 대응하는 식을 얻는다.

% inherited: noether_paper02_ko_translation_001_20260804\targets\Noether_P02_Korean_T04_U19_UNCHECKED.tex
% sha256: E5074116B8D4CE64D5683D15E0B0F12F149F3E2303F6F4925D4F4BC855639542
따라서 식
\[
s_x^m t_x^{n-\lambda}(tuv)^\lambda(stu)^\mu u_\tau^{\nu-\varkappa}v_\tau^\varkappa
\]
은 하나의 합성 편극으로 환원된다. 또한 항등식
\[
(stv)u_\tau=(stu)v_\tau-s_\tau(tuv)
\]
을 적용하면, 형식열의 더 높은 형식들에 대응하는 같은 방식으로 구성된 식들의 합으로 환원된다.

이를 반복하면 다음 전개에 이른다.
\[
\tag{III}
s_x^m t_x^{n-\lambda}(tuv)^\lambda(stu)^\mu u_\tau^{\nu-\varkappa}v_\tau^\varkappa
=
\sum_{p,r,\varrho}u_x^p\cdot C_{pr\varrho}\cdot
\bigl[s_\tau^\varrho(stu)^{\mu-\varrho+r}\bigr]_{\widehat{uv}^{\lambda-r+\varrho}v^{\varkappa-\varrho+p},\,v_x^{r-p}}.
\]
수치 계수 $C_{pr\varrho}$는 (II.)를 반복하여 생기는 방정식계의 계수 $c_{r\varrho}$와 $c'_{r\varrho}$로부터 유일하고 재귀적으로 계산된다(42쪽의 예 참조). 나머지 특수화들에서도 이와 같은 전개를 얻는다.

% inherited: noether_paper02_ko_translation_001_20260804\targets\P02_T05_U20_r2.tex
% sha256: C67FDDF1FE3A6C9CEC5AED0EA8C66E8B82587E6BDA8B89534B3CA821B00273A4
\begin{center}
\textbf{§ 3.}\quad \emph{환원 정리.}
\end{center}

형식과 형식계의 환원에 관한 잘 알려진 정리들을 이제 § 1과 § 2에 연결하고자 한다. 이를 위해 이항 형식론에서 가져온 몇몇 정의가 필요하다.

a) 주어진 형식열을 법으로 하는 \emph{상대 완전계}란 다음 성질을 지닌 형식계를 말한다. 이 형식들을 임의로 곱한 것을 수축하여 생기는 모든 구성식은 그 계의 형식들로 된 다항식으로 나타낼 수 있으며, 다만 계의 형식들을 법의 계와 수축하여 생긴 식들로 이루어진 가법적 항은 제외한다\srcfn{*)}{Gordan-Kerschensteiner, Vorlesungen über Invariantentheorie. Bd. II. S. 227.}.

b) 어떤 형식이 불변식을 인수로 갖는 형식들 또는 \glqq{}더 높은 형식들\grqq{}로 나타낼 수 있을 때 그 형식을 \emph{환원 가능}하다고 한다. 다시 말해, 해당 형식이 속한 전체 형식열의 더 높은 형식들, 또는 법의 기호를 더 높은 차수로 포함하는 형식들로 나타낼 수 있을 때이다.

이제 환원자에 관한 정리들\srcfn{**)}{Gordan, Math. Annalen Bd. XVII. S. 222.}은 가장 일반적인 형태로 다음과 같이 진술할 수 있다.

\textbf{정의:} \emph{환원자란 환원 가능한 형식열이다.}

% inherited: noether_paper02_ko_translation_001_20260804\targets\Noether_P02_Korean_T05_U21_UNCHECKED.tex
% sha256: 3AF1EFF132DBEC3A4573DF6C44F5FFEB0D5E11DC6BFDABA7F92F1AC2E2AF9765
\textbf{정리 II.} \emph{형식열의 시작형식이 그 항들 가운데 하나가 환원자와 수축하여 생겼기 때문에(환원자를 인수로 갖기 때문에) 환원 가능하고, 그 형식열의 끝형식도 바로 이 항에서 수축하여 생겼다면, 전체 형식열은 환원 가능하다}\srcfn{*)}{정리 II가 가져오는 단순화는 다음 예에서 볼 수 있다. 특수 형식 $f=x_1^3x_2+x_2^3x_3+x_3^3x_1$에 대해서는 $a_y^2a_x^2u_y^2$가 환원자이다. 따라서 정리 II에 의해 형식열
\[
\theta_\nu(\vartheta\nu x)\theta_x^3u_\vartheta u_\nu^2
=
a_\nu^2(abu)-\underbrace{a_\nu b_\nu(aba)}_{0}
\]
의 환원이 따른다. 이는 $\theta_\nu^4=a_\nu^2b_\nu^2(abu)^2$가 항 $a_\nu^2(abu)$에서 수축하여 생겼기 때문이다. 반면 Gordan은 앞서 인용한 문헌 231쪽에서 형식들
\[
\theta_\nu(\vartheta\nu x),\quad \theta_\nu(\vartheta\nu x)^2,\quad
\theta_\nu^2,\quad \theta_\nu^2(\vartheta\nu x),\quad
\theta_\nu^2(\vartheta\nu x)^2
\]
을 하나씩 환원한다.}.

\textbf{증명:} 형식열은 시작항을 그 자체 안에서 수축하여 생긴다. 따라서 한편으로는 환원자와 더 높게 수축하여, 다른 한편으로는 환원자를 그 자체 안에서 수축하여 생긴다. 두 경우 모두 § 2에 따라 형식열의 모든 형식은 환원자의 형식열에 대한 편극(전이연산)으로 나타낼 수 있다.

% inherited: noether_paper02_ko_translation_001_20260804\targets\P02_T05_U22_r2.tex
% sha256: 3A9D7658957064014D86A6F951B85B06541EF8B3C24838EE48638B48D92FA14C
다른 환원 방법들에서는 시작형식에서 전체 형식열로 나아가는 추론이 더 이상 성립하지 않는다. 그 방법들은 다음과 같다.

1) 이른바 \emph{\glqq{}이중 환원\grqq{}}.

이는 서로 독립적인 두 환원자를 인수로 갖는 식을 두 가지 방식으로 환원하는 것을 뜻하며, 이로써 환원의 도착점이 된 더 높은 형식들 사이의 관계가 생긴다.

한편 \glqq{}이중 환원\grqq{}을 체계적으로 적용하면 모든 관계를 얻어야 한다\srcfn{**)}{예를 들어 \glqq{}이중 환원\grqq{}으로 다음을 찾았다. 형식 $a_\eta^3$(§ 10)의 환원인데, 이 형식은 Maisano가 불필요하게 계에 포함시킨 유일한 형식이다. 이어서 서론의 주석에서 언급한 3개의 공변식과 3개의 불변식 사이의 관계도 찾았다.}. 다른 한편, 서로 다른 식들에 적용한 \glqq{}이중 환원\grqq{}이 새로운 관계를 내놓지 않을 때에는 더 높은 형식들의 독립성에 대한 판정 기준과 계산 검산을 동시에 얻게 된다.

% inherited: noether_paper02_ko_translation_001_20260804\targets\Noether_P02_Korean_T05_U23_UNCHECKED.tex
% sha256: 71ADA33BCEFFB608ED678FA88E239A7D6B853E5C1CF967DDBB76130910017092
2) \emph{분해 가능 형식과의 수축.}

여기서 \glqq{}분해 가능 형식\grqq{}이란 --- 통상적인 개념을 약간 일반화하여 --- 더 낮은 차수의 형식들의 곱과 \glqq{}더 높은\grqq{} 형식들로 나타낼 수 있는 형식을 뜻한다. 형식들의 곱은 한 번 수축하면 다시 곱으로 옮겨 간다. 마찬가지로 \emph{비선형} 형식들의 곱은 특수화 a')와 b')(§ 2)에서 곱 정리에 따라 두 번 수축하면 다시 곱으로 옮겨 간다.
\[
a_yb_y a_xb_x=\frac12a_y^2\cdot b_x^2+\frac12b_y^2\cdot a_x^2-\frac12(ab\widehat{yx})^2.
\]
따라서 분해 가능 형식의 1차 편극(전이연산)과 몇몇 2차 편극은 다시 분해 가능 형식이다. 그러므로 다음이 따른다.

분해 가능 형식에 대한 1차(2차) 전이연산의 한 \emph{항}이 분해 가능 형식과 한 번(두 번) 수축하여 생겼다면, 그 항 자체도 분해 가능 형식이 된다.

a) 분해 가능 형식의 형식열에서 바로 다음의 더 높은 형식들이 분해되거나 환원 가능해지는 경우, 또는 한 번 수축할 때 특수화 $y=\widehat{uv}$나 $v=\widehat{xy}$가 일어나는 경우,

b) 나머지 한 번의 수축들이 더 높은 형식으로 이어지는 경우(§ 12, B. $H_k$와 $H_k(k\eta x)$ 참조).

그러한 전이연산의 항은 다음 경우에도 분해될 수 있다.

c) 나머지 한 번의 수축들이 더 낮은 형식으로 이어지고, 이 낮은 형식들이 분해 가능 형식과의 수축과는 독립적으로 분해되는 경우이다. 즉 곱들과 환원하려는 형식으로 나타낼 수 있는 경우인데, 이때에는 해당 항의 수치 계수가 항등적으로 소멸하지 않는다는 사실을 매번 계산으로 먼저 확인해야 한다. 이는 낮은 형식의 \glqq{}이중 환원\grqq{}에 지나지 않는다(§ 23, C. $H_q(\theta H u)(\theta s u)$ 참조).

% inherited: noether_paper02_ko_translation_001_20260804\targets\P02_T05_U24_r2.tex
% sha256: 12814795C05A512A8922A8683E5DE25484BF88653A8B9B65B9963E37B24A0303
분해 가능 형식은 야코비 행렬식과의 모든 한 번의 수축으로 생기며\srcfn{*)}{Gordan, 앞서 인용한 문헌 § 3 참조.}, 그 밖에도 몇몇 더 높은 수축으로 생긴다. 다음이 성립한다.
\[
(st\widehat{yx})s_y=\frac12\,t\cdot s_y^2-\frac12\,s\cdot t_y^2+\frac12(st\widehat{yx})^2,
\]
\[
(st\widehat{yx})^2s_y^2=\frac13\,t\cdot s_y^3-\frac13\,s\cdot t_y^3+s_y(st\widehat{yx})^2-\frac13(st\widehat{yx})^3.
\]
유사한 공식이 쌍대 형식들
\[
(\sigma\tau\widehat{vu})v_\sigma
\quad\text{및}\quad
(\sigma\tau\widehat{vu})v_\sigma^2
\]
에 대해서도 성립한다.

% inherited: noether_paper02_ko_translation_001_20260804\targets\Noether_P02_Korean_T05_U25_UNCHECKED.tex
% sha256: 2038228E0DF3834E2AD3D2B804C7BC8DDF7197295220E355C3B3E20204D37E78
이 절의 정리들에 따라 $S$와 $T$의 수축에서 생기는 모든 환원 불가능한 형식을 세우기 위한 규칙은 다음과 같다.

가장 낮은 수축부터 시작하여, 미리 임의로 정해 둔 경로(예를 들어 행별로 또는 대각 성분에 대칭이 되도록)를 따라 형식열 $ST$를 구성해 나가다가 환원자를 인수로 갖는 형식에 이를 때까지 계속한다. 이 형식이 정하는 형식열은 끝형식이 정리 II에서 세운 조건을 만족하는 즉시 제외한다. 환원 가능한 형식열을 모두 제거한 뒤에는 이중 환원을 적용하여 나머지 환원 가능한 형식들을 모두 제거하고, 분해 가능 형식들도 모두 제거한다. 서로 독립적으로 환원 가능한 형식열들이 겹쳐서 두 번 덮는 전체 형식열의 위치에서는 더 높은 형식들을 환원할 계기가 생긴다(§ 11, D 참조).

% inherited: noether_paper02_ko_translation_001_20260804\targets\Noether_P02_Korean_T06_U26_UNCHECKED.tex
% sha256: FC4C9DF3C01A72D68EF72D523DEF0222EF114B13328F2629A250D597FF64BF47
\begin{center}
\textbf{제II장. \quad 상대 완전 기본계.}

\textbf{§ 4.}\quad\emph{첫 번째 상대 완전계.}
\end{center}

이하에서는 첫 번째 정리를 제외하고 4차 형식에 한정한다. 그러나 모듈들을 적절히 선택하면 §§ 5와 6의 결과를 더 높은 차수의 형식으로 쉽게 옮길 수 있으며, 이 결과들은 3차 형식에도 마찬가지로 성립한다\srcfn{*)}{3차 형식의 경우 이와 유사하게 유도한 대응 공식은 다음과 같다:
\[
(abc)a_x^2b_y^2c_z^2=\frac13(sxy)(sxz)(syz)+\frac23\Delta_x\Delta_y\Delta_z\cdot(xyz),\qquad (\mathrm{IV}.)
\]
\[
\begin{gathered}
s=(abc)(abu)(acu)(bcu),\qquad \Delta=(abc)^2a_xb_xc_x,\\
t=\Delta_\vartheta(\Delta\theta u)^2u_\vartheta=\theta_s^2u_su_\vartheta^2,
\qquad T=a_t^3,
\end{gathered}
\]
\[
\left.
\begin{aligned}
(a\theta u)^2&=\frac13 u_xs,\\
a_\vartheta&=\frac13u_x\cdot\Delta,
& a_\vartheta(a\theta u)&=0,
& a_\vartheta(a\theta u)^2&=s,\\
a_\vartheta^2&=\Delta,
& a_\vartheta^2(a\theta u)&=0.
\end{aligned}
\right\}\tag{1.}
\]
}.

% inherited: noether_paper02_ko_translation_001_20260804\targets\Noether_P02_Korean_T06_U27_UNCHECKED.tex
% sha256: D95501425A53AA5DFA59B178F7AAA1E1523F8DFCD88DD42EDE036484AC123B40
임의 차수의 형식에 대한 첫 번째 상대 완전계는 잘 알려진 다음 정리로 얻어진다.

\clearpage
\noindent\textbf{정리 III.}
\emph{이항 영역에서 넘겨받은 형식들, 즉 대응하는 이항 형식의 형식계에 속한 형식들로부터 Clebsch의 전이 원리에 따른 경계 첨가로 생기는 형식들은 $(abc)$를 법으로 하는 상대 완전계를 이룬다.}

\noindent\textbf{증명:}
형식 $Q'_1,\ldots,Q'_n$이 하나의 이항 기본형의 완전계를 이룬다는 다음 이항 관계
\[
Q'-F(Q')=0=\sum_\lambda\bigl\{(ab)c_x+(bc)a_x+(ca)b_x\bigr\}\varphi_\lambda\srcfnmark{*)}
\]
\srcfntext{*)}{Gordan--Kerschensteiner, 앞의 책 134쪽.}
에는 경계 첨가를 통하여 다음 삼원 관계가 대응한다.
\[
Q-F(Q_i)=\sum_\lambda u_x\cdot(abc)\varphi_\lambda .
\]
그러나 이것은 바로 $(abc)$를 법으로 하는 상대 완전계의 정의 방정식이다. 4차 형식에서는 넘겨받은 형식들의 계가 다음 형식들로 이루어진다.

% inherited: noether_paper02_ko_translation_001_20260804\targets\Noether_P02_Korean_T06_U28_UNCHECKED.tex
% sha256: 00E6662910F2859DF937132D037EA36894B6B4B62BD00279E88E50DBCD5863CA
\[
\begin{gathered}
f=a_x^4,\qquad
\theta=\theta_x^4u_\vartheta^2=(abu)^2a_x^2b_x^2,
\qquad
K=K_x^6u_k^3=(a\theta u)u_\vartheta^2a_x^3\theta_x^3,\\
j=u_j^6=(a\theta u)^4u_\vartheta^2,\qquad
\nu=u_\nu^4=(abu)^4,
\end{gathered}
\]
이들 가운데 처음 네 형식은 $((abc),\nu)$를 법으로 하는 상대 완전계를 이룬다.

% inherited: noether_paper02_ko_translation_001_20260804\targets\Noether_P02_Korean_T07_U29_UNCHECKED_TEXREPAIR_v002_20260804.tex
% sha256: B380060629B2733D4460DBB583DD2866DCD0B017F3CAB84A07EE855A9A9241DA
\begin{center}
\textbf{§ 5.}\quad\emph{모듈 $(abc)$를 모듈 $\Delta$와 $\nu$로 환원.}
\end{center}

괄호 인자 하나로만 주어진 모듈 $(abc)$는 정의(§ 3, a)에 따라 계의 형식들로 환원해야 한다. 달리 말하면 형식열 $(abc)$를 유일하게 정규화해야 한다(§ 1 참조).

\[
a_s^2a_xu_s=\frac13a_s^3\cdot u_x=\frac13S\cdot u_x. \tag{2.}
\]
\[
\left.
\begin{aligned}
(a\Delta u)^2&=a_s-\frac{S}{6}u_x^2,\\
(a\Delta u)^3&=0.
\end{aligned}
\right\} \tag{3.}
\]
\[
\left.
\begin{aligned}
\theta(s)=(ss_1x)^2&=2a_t+\frac13S\cdot\theta-\frac13Tu_x^2,\\
\theta(t)=(tt_1x)^2&=-\frac13S\cdot a_t+\frac23T\cdot a_s+\frac1{18}S^2\cdot\theta-\frac1{18}u_x^2\cdot S\cdot T.
\end{aligned}
\right\} \tag{4.}
\]

% inherited: noether_paper02_ko_translation_001_20260804\targets\Noether_P02_Korean_T07_U30_UNCHECKED.tex
% sha256: 85A15F0D8625FB426EE485672B91A72BBC2B35A39B0A24B20CDC895EE01EE3FD
모듈의 열: $(abc)$; $\Delta$; $s$; $t$ (모듈 $t$의 계는 유일한 형식 $t$로 이루어진다).

다음 형식들이
\[
\Delta=(abc)^2a_x^2b_x^2c_x^2,\qquad \nu=(abu)^4
\]
모듈로서 충분하다는 것, 즉 형식들
\[
\begin{gathered}
\Delta=(abc)^2a_x^2b_x^2c_x^2,
\qquad a_\nu=(abc)(bcu)^3a_x^3,\\
a_\nu^2=(abc)^2(bcu)^2a_x^2,
\qquad i=a_\nu^4=(abc)^4
\end{gathered}
\]
이 형식열 $(abc)$를 정한다는 것을 보일 것이다. 형식열 $a_\nu$에는 다음 항등식에 따라 형식 $a_\nu^3$이 빠진다.
\[
a_\nu^3=(abc)^3a_x(bcu)=\frac13u_x\cdot(abc)^4=\frac13i\cdot u_x.
\]

% inherited: noether_paper02_ko_translation_001_20260804\targets\Noether_P02_Korean_T07_U31_UNCHECKED.tex
% sha256: B857852737512EA7A42779E9286A8D7376E0376597B3B41DCFFAF926F5347FB4
모듈 하나를 더 높은 모듈로 환원하려면 정의에 따라, 그 모듈과의 가장 일반적인 수축 또는 고려 대상인 모든 특수 수축을 더 높은 모듈과의 수축으로 환원해야 한다.

우리의 경우 가장 일반적인 수축은 다음 식들에서 생긴다.
\[
(abc)a_x^3b_y^3c_z^3,
\qquad
(abc)a_x^2b_y^2c_z^2a_xb_xc_x
\]
(3차 형식들에서 넘겨받음),
\[
(abc)a_xb_yc_za_x^2b_x^2c_x^2
\]
(2차 형식들에서 넘겨받음), 그리고 이 식들을 편극하여 생긴다.

이 식들을 $\Delta$와 형식열 $a_\nu$의 편극, 또는 각각 그 시작항들로 계산하기 위해 § 2에서와 같이 간접적인 길을 택한다. 편극항들
\[
(abc)^2a_x^2b_y^2c_z^2,
\qquad
 a_\nu(\nu xy)(\nu yz)(\nu zx)a_xa_ya_z=(abc)(bc\widehat{xy})(bc\widehat{yz})(bc\widehat{zx})a_xa_ya_z\quad\text{등}
\]
을 기호 인자 $(abc)$를 갖는 식들의 선형 결합으로 전개하고, 방정식계를 역으로 풀어 원하는 관계들과 서로 다른 변수 조합에 따라 취한 편극들 사이의 한 관계를 얻는다. 계산에는 항등식 정리와 행렬식의 곱 정리를 이용한다.

% inherited: noether_paper02_ko_translation_001_20260804\targets\Noether_P02_Korean_T07_U32_UNCHECKED.tex
% sha256: C2882E0A1E3103D403F5012CCFD8930D52E4B4E6ED180A23054CD6FAF4DD42EB
$(abc)a_x^3b_y^3c_z^3$에 대한 방정식계는 다음과 같다.
\begingroup\small
\[
\begin{aligned}
1)\quad
\sum_3a_\nu(\nu yz)^3a_x^3
&=6(abc)a_x^3b_y^3c_z^3-6\sum_3(abc)a_x^3b_y^2b_zc_z^2c_y\srcfnmark{*)},\\[0.4em]
2)\quad
3(abc)^2a_x^2b_y^2c_z^2\cdot(xyz)-\frac12\sum_3a_\nu^2(\nu yz)^2a_x^2\cdot(xyz)
&=2\sum_3(abc)a_x^3b_y^2b_zc_z^2c_y\\
&\quad-4\sum_3(abc)a_xa_ya_zb_y^2b_xc_z^2c_x,\\[0.4em]
3)\quad
a_\nu(\nu xy)(\nu yz)(\nu zx)a_xa_ya_z
&=2\sum_3(abc)a_xa_ya_zb_y^2b_xc_z^2c_x,\\[0.4em]
4)\quad
\sum_3a_\nu(\nu yz)^3a_x^3-\frac32\sum_3a_\nu^2(\nu yz)^2a_x^2\cdot(xyz)+\frac16i\cdot(xyz)^3
&=6\sum_3(abc)a_xa_ya_zb_y^2b_xc_z^2c_x.
\end{aligned}
\]
\endgroup
\srcfntext{*)}{$\sum_3$은 시작항에서 변수들을 순환 치환하여 생기는 세 항의 합을 뜻한다. 방정식들은 그 합의 각 항에 대해서도 개별적으로 성립한다.}

% inherited: noether_paper02_ko_translation_001_20260804\targets\Noether_P02_Korean_T07_U33_UNCHECKED.tex
% sha256: DE012BABCCA5B58E66C2B65FA26DC0AAEFDF3F8BACA6E7CC037A1A3B3E10CAB0
이로부터 $(abc)a_x^3b_y^3c_z^3$에 대해 다음을 얻고, $(abc)a_x^2b_y^2c_z^2a_xb_xc_x$와 $(abc)a_xb_yc_za_x^2b_x^2c_x^2$에 대해서도 유사한 계산으로 다음을 얻는다.
\[
\begin{aligned}
(abc)a_x^3b_y^3c_z^3
&=\frac32(abc)^2a_x^2b_y^2c_z^2\cdot(xyz)
+\frac32a_\nu(\nu xy)(\nu yz)(\nu zx)a_xa_ya_z
-\frac1{12}i\cdot(xyz)^3,\\[0.4em]
(\mathrm{IV}.)\quad
(abc)a_x^2b_y^2c_z^2a_xb_xc_x
&=\frac23(abc)^2a_x^2b_yb_xc_zc_x\cdot(xyz)
+\frac13a_\nu(\nu xy)(\nu yz)(\nu zx)a_x^3\\
&\quad-\frac16a_\nu^2(\nu xy)(\nu zx)a_x^2\cdot(xyz),\\[0.4em]
(abc)a_xb_yc_za_x^2b_x^2c_x^2
&=\frac16(abc)^2a_x^2b_x^2c_x^2\cdot(xyz).
\end{aligned}
\]

이로써 네 형식 $f,\theta,K,j$로 이루어진, $(\Delta,\nu)$를 법으로 하는 상대 완전계를 얻었다.

따라서 이항 영역에서 넘겨받지 않은 $f$와 $\theta$의 모든 전이연산, 그리고 이항 영역에서 $(ab)^4$로 환원 가능한 전이연산 $(a\theta u)^2$는 모두 기호 $\Delta$와 $\nu$로 나타낼 수 있다.

% inherited: noether_paper02_ko_translation_001_20260804\targets\Noether_P02_Korean_T07_U34_UNCHECKED.tex
% sha256: 3D941420042813E2F604847FBEE1EE4DEC34955F2AD7CC27B3187428D82BD24A
이하에서 기본이 되는 환원 공식들을 전개 IV의 특수화로, 또는 더 짧게는 직접 계산(기호의 교환)으로 얻는다. $a_\vartheta(a\theta u)^2$, $a_\vartheta^2(a\theta u)^2$, $(a\theta u)^2$에 대해서는 다음과 같이 둔다.
\[
\begin{aligned}
a_\vartheta(a\theta u)^2
&=(abc)(bcu)(abu)^2-\frac13(abc)(bcu)\{(bcu)a_x-u_x\cdot(abc)\}^2,\\
a_\vartheta^2(a\theta u)^2
&=(abc)^2(abu)^2-\frac13(abc)^2\{(bcu)a_x-u_x\cdot(abc)\}^2.
\end{aligned}
\]
$(a\theta u)^2$에 대해서는\srcfn{*)}{Gordan--Kerschensteiner, 앞서 인용한 문헌 92쪽.} 다음을 쓴다.
\[
\begin{aligned}
(abu)a_y^3b_x^3&=\frac32u_\vartheta(\vartheta yx)\theta_y^2+\frac14(\nu yx)^3,\\
(abu)(acu)^3b_x^3&=\frac32u_\vartheta(\vartheta\widehat{au}x)(\theta au)^2+\frac14(\nu\widehat{au}x)^3.
\end{aligned}
\]

% inherited: noether_paper02_ko_translation_001_20260804\targets\Noether_P02_Korean_T07_U35_UNCHECKED.tex
% sha256: 9D6D205AB9AEF2984C462CBC4472A36A9249A49A1D51445682597BDC1F529401
\begingroup\small
\[
\begin{array}{r@{\quad}|l@{\qquad}c|}
& \underline{(a\theta u)^2}=\frac16\nu\cdot f-\frac23u_x\cdot a_\nu+\frac23u_x^2\cdot a_\nu^2-\frac{i}{18}u_x^4
& \bigstar\\[0.55em]
a_\vartheta=\frac13u_x\cdot\Delta
& \underline{a_\vartheta(a\theta u)}=0
& \\[0.55em]
& \underline{a_\vartheta^2}=\Delta
&
\end{array}
\]
\vspace{-0.5em}
\[
\begin{array}{r@{\quad}c@{\quad}|l@{\qquad}|l}
(1.) & \bigstar
& \underline{(a\theta u)^3}=0
& \underline{(a\theta u)^4}=j\\[0.55em]
&
& \underline{a_\vartheta(a\theta u)^2}=-\frac56a_\nu+\frac76u_x\cdot a_\nu^2-\frac{i}{9}u_x^3
& \underline{a_\vartheta(a\theta u)^3}=0\\[0.55em]
&
& \underline{a_\vartheta^2(a\theta u)}=0
& \underline{a_\vartheta^2(a\theta u)^2}=\frac23a_\nu^2-\frac{i}{9}u_x^2
\end{array}
\]
\endgroup

% inherited: noether_paper02_ko_translation_001_20260804\targets\Noether_P02_Korean_T07_U36_UNCHECKED.tex
% sha256: 8191C0DB71924B3957431FBCBEF22A18AB705682B227746962200AA394DE63F5
여기에 역시 모듈 $(abc)$의 환원에서 생긴 다음 공식도 취한다.
\[
\begin{array}{r@{\qquad}l}
(2.) & a_\nu^3a_xu_\nu=\frac13i\cdot u_x.
\end{array}
\]

\emph{공식 (1.)과 (2.), 그리고 기본형 $\nu$에 대해 이와 유사하게 만든 공식들로부터, 분해 가능 형식들에 관한 관계들을 제외한 이후의 모든 환원 공식이 편극을 통해 유도된다.} 공식 (1.)에서 다음을 알 수 있다.

형식열은 다음과 같이 환원된다.
\[
\begin{array}{rcl}
a_\vartheta(a\theta u)^2 &\text{는}& \text{기호 }\nu\text{만으로 환원되고},\\[0.2em]
a_\vartheta &\text{는}& \text{기호 }\Delta\text{와 }\nu\text{로 환원되고},\\[0.2em]
(a\theta u)^2 &\text{는}& \text{기호 }j,\Delta\text{ 및 }\nu\text{로 환원되고},\\[0.2em]
(a\theta u)^2\text{에 수축 I을 적용하면} &\text{}& \text{기호 }j\text{와 }\nu\text{로 환원되고},\\[0.2em]
(a\theta u)^2\text{에 수축 II를 적용하면} &\text{}& \text{기호 }\Delta\text{와 }\nu\text{로 환원된다}.
\end{array}
\]

% inherited: noether_paper02_ko_translation_001_20260804\targets\Noether_P02_Korean_T07_U37_UNCHECKED.tex
% sha256: FE94279496767C5DE99EAE9B18D5C311494D61C4AEF4D2ADCACBF343FAD6E7A8
따라서 다음과 같이 바꾸어 쓸 수 있다.
\begingroup
\renewcommand{\labelenumi}{\arabic{enumi})}
\begin{enumerate}
\item $(\Delta,\nu)$를 법으로 하여: $j$와의 수축을 $(a\theta u)^2$ 또는 $(a\theta u)^3$과의 대응하는 수축으로 바꾼다.
\item $(\nu)$를 법으로 하여: $\Delta$와의 수축을 $a_\vartheta$ 또는 $a_\vartheta(a\theta u)$와의 대응하는 수축으로 바꾸거나, $(a\theta u)^2$과의 특수한 수축으로도 바꾼다(후자는 한 번의 수축 I로만 이어진다).
\end{enumerate}
\endgroup
특히 다음이 성립한다.
\[
\begin{aligned}
(a\theta u)^2a_y^2\theta_x^2&=\frac13(jyx)^2\modu{\nu},&
(a\theta u)^2a_y\theta_y&=-\frac16(jyx)^2\modu{\nu},\\[0.4em]
(a\theta u)^2v_\vartheta^2&=\frac13(\Delta uv)^2\modu{\nu},&
(a\theta u)(a\theta v)u_\vartheta v_\vartheta&=-\frac16(\Delta uv)^2\modu{\nu},\\[0.4em]
(a\theta u)^2a_yv_\vartheta^2&=\frac13(\Delta uv)^2\Delta_y\modu{\nu}.
\end{aligned}
\]

% inherited: noether_paper02_ko_translation_001_20260804\targets\Noether_P02_Korean_T08_U38_UNCHECKED.tex
% sha256: B17A0666B0A7B97AF26322FB0DCE8DBF44743C89D123EA87A546F230F11991B5
\begin{center}
\textbf{§ 6.}\quad\emph{모듈 $(\Delta,\nu)$를 모듈 $(\nu)$로 환원.}
\end{center}

앞 절의 환원 공식들은 $\nu$를 법으로 하는 상대 완전계를 직접 세울 수단을 제공한다. 넘겨받은 형식들의 계에는 형식 $\Delta$와 $(a\Delta u)=N$만 덧붙이면 된다는 것을 보일 것이다(3차 형식의 경우와 유사하며, 사실 일반적으로 성립한다).

모듈 $\Delta$를 환원하기 위해 고려해야 할 모든 특수 수축, 즉 $(\Delta,\nu)$를 법으로 하는 상대 완전계와 $\Delta$의 계 사이의 수축을 살펴본다. 계수에 관한 차수가 가장 낮은 형식들, 곧 $f$와 $\Delta$의 수축에서 출발한다.

형식 $(a\Delta u)^2$, $(a\Delta u)^3$, $(a\Delta u)^4$를 환원하기 위해 § 5에 따라 기호 $\Delta$를 형식열의 더 낮은 형식으로 바꾼다. 즉 다음 식들을
\[
a_\vartheta b_\vartheta(b\theta u)^2,
\qquad a_\vartheta(a\theta u)b_\vartheta(b\theta u)^2,
\qquad a_\vartheta(a\theta u)b_\vartheta(b\theta u)^3
\]
1) 형식열 $a_\vartheta$의 편극(기호 $\Delta$와 $\nu$)에 따라,

2) 형식열 $b_\vartheta(b\theta u)^2$의 편극(기호 $\nu$)에 따라

전개하여 다음 환원 공식들을 얻는다(계산은 아래 참조).

% inherited: noether_paper02_ko_translation_001_20260804\targets\Noether_P02_Korean_T08_U39_UNCHECKED.tex
% sha256: 429786E189AB1BAEB8AD980E8E865F12DB71160E273CBACF2B7CECBF0C63A67C
\[
\begin{aligned}
\text{a)}\quad \underline{(a\Delta u)^2}
&=\frac12(\vartheta \nu x)^2-\frac25f\cdot a_\nu^2
-\frac45u_x\cdot\theta_\nu(\vartheta \nu x)^2
+u_x^2\left\{\frac16i f-\frac1{40}s\right\},\\
(3.)\quad\text{b)}\quad \underline{(a\Delta u)^3}
&=-\frac3{10}\theta_\nu(\vartheta \nu x),\\
\text{c)}\quad \underline{(a\Delta u)^4}
&=\frac{21}{10}\theta_\nu^2-\frac3{10}H
-\frac65u_x\cdot\theta_\nu^3+\frac1{10}u_x^2\cdot\varrho.
\end{aligned}
\]

($a_\vartheta^2$를 소거한 뒤 식 $a_\vartheta^2(b\theta u)^2$, $a_\vartheta^2(b\theta u)^3$, $a_\vartheta^2(a\theta u)^2(b\theta u)^2$, $a_\vartheta^2(a\theta u)(b\theta u)^3$을 이중 환원해도 같은 결과를 얻는다.)
공식 (3.)으로부터 $(a\Delta u)^2$가 모듈 $\nu$에 관한 환원자임이 따른다. 따라서 다음을 얻는다.

1) $\Delta$의 계의 환원: 형식열 $(\Delta\Delta' u)^2\equiv a_\vartheta^2(a\Delta u)^2$은 환원자 $(a\Delta u)^2$를 인수로 갖는다.

2) 모듈 $\Delta$와의 수축으로 생긴 계의 환원:

형식열 $(\theta\Delta u)^2$은 환원자 $(a\Delta u)^2$를 인수로 갖는다.

형식 $(\theta\Delta u)$와 $\Delta_\vartheta$에 대해서는 다음을 얻는다.
\[
(a\Delta u)^2(abu)=(\theta\Delta u)-u_x\cdot\Delta_\vartheta(\theta\Delta u),
\]
\[
(a\Delta u)(a\Delta b)=-\frac12\Delta_\vartheta+\frac12u_x\cdot\Delta_\vartheta^2.
\]

그러나 환원자 $(\theta\Delta u)$로부터 모듈 $\Delta$와의 수축으로 생기는 계의 환원이 따른다.

% inherited: noether_paper02_ko_translation_001_20260804\targets\Noether_P02_Korean_T08_U40_UNCHECKED.tex
% sha256: 702145461A4C04C16EE75608D71B5F57121F371AF1ABC30807F589D2AFFB4750
\bigskip

따라서 $\nu$를 법으로 하는 상대 완전계는 다음 여섯 형식으로 이루어진다.
\[
f,\ \theta,\ K,\ j,\ \A,\ N.
\]

§ 2의 급수 전개에 대한 예로 공식 (3.)a의 유도를 자세히 제시한다. 이후의 계산에서는 최종 공식 III을 직접 세울 것이다. (소멸하는 형식의 편극은 계산 도중에 이미 생략하였다. 첫째 줄은 전개 II에 따라 대입한 값을, 둘째 줄은 연립계의 마지막 방정식부터 시작하여 재귀적으로 구한 최종값을 나타낸다.)

% inherited: noether_paper02_ko_translation_001_20260804\targets\Noether_P02_Korean_T08_U41_UNCHECKED.tex
% sha256: 624C64DE3ACB6BBC24447E67F27D5989664E67BD4837B4298935196EC54945E2
전개 II에 따라 다음이 성립한다.
\[
\begin{aligned}
\text{1)}\quad&m=3.\quad n=4.\quad \lambda=2.\quad \mu=0.\quad \nu=\varkappa=1,\\[0.25em]
\underline{a_\vartheta b_\vartheta(b\theta u)^2}
&=[a_\vartheta]_{\widehat{ub^2}}b+\frac67\{a_\vartheta b_\vartheta(a\theta u)(b\theta u)-u_x\cdot a_\vartheta b_\vartheta(a\theta b)(b\theta u)\}\\
&\quad-\frac17\{a_\vartheta(a\theta u)^2b_\vartheta-2u_x\cdot a_\vartheta(a\theta u)(a\theta b)b_\vartheta+u_x^2\cdot a_\vartheta(a\theta b)^2b_\vartheta\}\\
&=\frac23(a\A u)^2+\frac15[a_\vartheta(a\theta u)^2]b+\frac1{30}f\cdot[a_\vartheta^2(a\theta u)^2]-\frac25u_x\cdot[a_\vartheta(a\theta u)^2]b^2\\
&\quad-\frac1{15}u_x\cdot[a_\vartheta^2(a\theta u)^2]b+\frac15u_x^2\cdot[a_\vartheta(a\theta u)^2]b^3+\frac1{30}u_x^2\cdot[a_\vartheta^2(a\theta u)^2]b^2;\\[0.7em]
\text{2)}\quad&m=2.\quad n=3.\quad \lambda=1.\quad \mu=1.\quad \nu=\varkappa=1,\\[0.25em]
\underline{a_\vartheta(a\theta u)b_\vartheta(b\theta u)}
&=\frac25\{a_\vartheta(a\theta u)^2b_\vartheta-u_x\cdot a_\vartheta(a\theta u)(a\theta b)b_\vartheta\}+\frac12a_\vartheta^2(b\theta u)^2\\
&\quad-\frac15\{a_\vartheta^2(b\theta u)(a\theta u)-u_x\cdot a_\vartheta^2(a\theta b)(b\theta u)\}\\
&=\frac12(a\A u)^2+\frac25[a_\vartheta(a\theta u)^2]b+\frac1{15}[a_\vartheta^2(a\theta u)^2]\cdot f\\
&\quad-\frac25u_x\cdot[a_\vartheta(a\theta u)^2]b^2-\frac1{10}u_x\cdot[a_\vartheta^2(a\theta u)^2]b+\frac1{30}u_x^2\cdot[a_\vartheta^2(a\theta u)^2]b^2;\\[0.7em]
\text{3)}\quad&m=2.\quad n=3.\quad \lambda=1.\quad \mu=1.\quad \nu=1.\quad \varkappa=2,
\end{aligned}
\]

% inherited: noether_paper02_ko_translation_001_20260804\targets\Noether_P02_Korean_T08_U42_UNCHECKED.tex
% sha256: E2551348279B7D5AEA85606D1E0B55DD1CFC5A62C55465E16DF6509CADD6AA43
\[
\begin{aligned}
\underline{a_\vartheta(a\theta b)b_\vartheta(b\theta u)}
&=\frac25\{a_\vartheta(a\theta b)(a\theta u)b_\vartheta-u_x\cdot a_\vartheta(a\theta b)^2b_\vartheta\}\\
&=\frac25[a_\vartheta(a\theta u)^2]b^2+\frac1{30}[a_\vartheta^2(a\theta u)^2]b-\frac25u_x\cdot[a_\vartheta(a\theta u)^2]b^3\\
&\quad-\frac1{30}u_x\cdot[a_\vartheta^2(a\theta u)^2]b^2;\\[0.7em]
\text{4)}\quad&m=1.\quad n=2.\quad \lambda=0.\quad \mu=2.\quad \nu=\varkappa=1,\\[0.25em]
\underline{a_\vartheta(a\theta u)^2b_\vartheta}
&=[a_\vartheta(a\theta u)^2]b+\frac23a_\vartheta^2(a\theta u)(b\theta u)\\
&=[a_\vartheta(a\theta u)^2]b+\frac16[a_\vartheta^2(a\theta u)^2]\cdot f-\frac16u_x\cdot[a_\vartheta^2(a\theta u)^2]b;\\[0.7em]
\text{5)}\quad&m=1.\quad n=2.\quad \lambda=0.\quad \mu=2.\quad \nu=1.\quad \varkappa=2\quad(\varkappa>\nu),\\[0.25em]
\underline{a_\vartheta(a\theta u)(a\theta b)b_\vartheta}
&=[a_\vartheta(a\theta u)^2]b^2+\frac13a_\vartheta^2(a\theta b)(b\theta u)\\
&=[a_\vartheta(a\theta u)^2]b^2+\frac1{12}[a_\vartheta^2(a\theta u)^2]b-\frac1{12}u_x\cdot[a_\vartheta^2(a\theta u)^2]b^2;\\[0.7em]
\text{6)}\quad&m=1.\quad n=2.\quad \lambda=0.\quad \mu=2.\quad \nu=1.\quad \varkappa=3,\\[0.25em]
\underline{a_\vartheta(a\theta b)^2b_\vartheta}&=[a_\vartheta(a\theta u)^2]b^3;\\[0.7em]
\text{7)}\quad&m=2.\quad n=4.\quad \lambda=2.\quad \mu=\nu=\varkappa=0.\quad(\text{전개 I에 따라}),\\[0.25em]
\underline{a_\vartheta^2(b\theta u)^2}&=[a_\vartheta^2]_{\widehat{ub^2}}+\frac1{10}\{[a_\vartheta^2(a\theta u)^2]\cdot f-2u_x\cdot[a_\vartheta^2(a\theta u)^2]b+u_x^2\cdot[a_\vartheta^2(a\theta u)^2]b^2\};\\[0.7em]
\text{8)}\quad&m=1.\quad n=3.\quad \lambda=1.\quad \mu=1.\quad \nu=\varkappa=0.\quad(\text{전개 I}),\\[0.25em]
\underline{a_\vartheta^2(a\theta u)(b\theta u)}&=\frac14[a_\vartheta^2(a\theta u)^2]\cdot f-\frac14u_x\cdot[a_\vartheta^2(a\theta u)^2]b;\\[0.7em]
\text{9)}\quad&m=1.\quad n=3.\quad \lambda=1.\quad \mu=\varkappa=1.\quad \nu=0.\quad(\text{전개 I}),\\[0.25em]
\underline{a_\vartheta^2(a\theta b)(b\theta u)}&=\frac14[a_\vartheta^2(a\theta u)^2]b-\frac14u_x\cdot[a_\vartheta^2(a\theta u)^2]b^2.
\end{aligned}
\]

% inherited: noether_paper02_ko_translation_001_20260804\targets\Noether_P02_Korean_T08_U43_UNCHECKED.tex
% sha256: 3BE177A12A9AF5D88A0735F173E4584E70656243EB633D8C47472C15379B6421
1)과 4)로부터 다음이 따른다.
\[
\begin{aligned}
(a\A u)^2
&=\frac65[a_\vartheta(a\theta u)^2]b+\frac15f\cdot[a_\vartheta^2(a\theta u)^2]+\frac35u_x\cdot[a_\vartheta(a\theta u)^2]b^2-\frac3{20}u_x\cdot[a_\vartheta^2(a\theta u)^2]b\\
&\quad-\frac3{10}u_x^2\cdot[a_\vartheta(a\theta u)^2]b^3-\frac1{20}u_x^2\cdot[a_\vartheta^2(a\theta u)^2]b^2,\\[0.4em]
(a\A u)^2&=-a_\nu b_\nu+\frac35f\cdot a_\nu^2+\frac45u_x\cdot a_\nu^2b_\nu-\frac3{20}u_x^2\cdot a_\nu^2b_\nu^2-\frac1{12}u_x^2\cdot i\cdot f.
\end{aligned}
\]

($\theta$의 기호로 환산하는 데 대해서는 §§ 8과 9를 참조하라.)
공식 (3.)a는 보조변수 $w=x\widehat{u}b$를 도입하여 전개 I을 직접 전이하면 더 짧게 계산할 수도 있다. 반면 공식 (3.)b와 (3.)c에서는 우리가 택한 방법이 이미 더 간단해진다. (3.)b에서는 § 9에 따라 인수 $u_x$를 가진 모든 항을 생략해야 한다는 것을 미리 알고 있다. (3.)b와 (3.)c의 계산에는 다음을 얻는다.

% inherited: noether_paper02_ko_translation_001_20260804\targets\Noether_P02_Korean_T08_U44_UNCHECKED.tex
% sha256: 2ADF9E02425CD74FAB3315A820A54CF5E58D1BF4296ED20492FC1A167CDAD372
\[
\begin{aligned}
(3.)\mathrm{b)}\quad 1)&\quad a_\vartheta(a\theta u)b_\vartheta(b\theta u)^2=\frac12(a\A u)^3+\frac45[a_\vartheta(a\theta u)^2]_{\widehat{ub}}b+\frac{77}{360}[a_\vartheta^2(a\theta u)^2]_{\widehat{ub}},\\
2)&\quad a_\vartheta(a\theta u)b_\vartheta(b\theta u)^2=[a_\vartheta(a\theta u)^2]_{\widehat{ub}}b+\frac{13}{36}[a_\vartheta^2(a\theta u)^2]_{\widehat{ub}};\\[0.5em]
(3.)\mathrm{c)}\quad 1)&\quad a_\vartheta(a\theta u)b_\vartheta(b\theta u)^3=\frac12(a\A u)^4+\frac65[a_\vartheta(a\theta u)^2]_{\widehat{ub^2}}b+\frac{53}{120}[a_\vartheta^2(a\theta u)^2]_{\widehat{ub^2}}\\
&\qquad-\frac65u_x\cdot[a_\vartheta(a\theta u)^2]_{\widehat{ub^2}}b^2-\frac35u_x\cdot[a_\vartheta^2(a\theta u)^2]_{\widehat{ub^2}}b+\frac{19}{120}u_x^2\cdot[a_\vartheta^2(a\theta u)^2]_{\widehat{ub^3}}b^2,\\
2)&\quad a_\vartheta(a\theta u)b_\vartheta(b\theta u)^3=\frac34[a_\vartheta^2(a\theta u)^2]_{\widehat{ub^2}}.
\end{aligned}
\]

% inherited: noether_paper02_ko_translation_001_20260804\targets\Noether_P02_Korean_T08_U45_UNCHECKED.tex
% sha256: 0576200DDC934F6063F499CC7E29EE4CB3FD05300E1282470B8CBC758859620D
부기: § 5에 따라 환원 가능한 $f$와 $j$의 전이연산들도 같은 방식으로 계산된다. 다음을 얻는다.
\[
\tag{4.}
\begin{aligned}
a_j&=-\theta_\nu+u_x\cdot\left\{\frac92\theta_\nu^2-\frac12H\right\}-3u_x^2\cdot\theta_\nu^3+\frac12u_x^3\cdot\varrho,\\
a_j^2&=\frac{21}{10}\theta_\nu^2-\frac3{10}H-2u_x\cdot\theta_\nu^3+\frac12u_x^2\cdot\varrho,\\
a_j^3&=-\frac7{10}\theta_\nu^3+\frac12u_x\cdot\varrho,\\
a_j^4&=\frac35\varrho.
\end{aligned}
\]

(3.)과 (4.) 및 이항 영역으로부터의 전이에 의해 다음이 따른다.
\[
(\theta\theta'u)^2=\frac13 j\cdot f\modu{\nu}\srcfn{*)}{Gordan--Kerschensteiner, 앞의 책 181쪽.}.
\]
형식열 $(\theta\theta'u)^2$의 나머지 형식들과 형식 $\theta'_\vartheta$는 기호 $\nu$로 환원 가능하다. 따라서 다음과 같이 바꿀 수 있다.
\[
\modu{\nu}:\quad \text{$f\cdot j$와의 수축을 $(\theta\theta'u)^2$와의 대응 수축으로.}
\]
또한 다음을 얻는다.
\[
(\vartheta\vartheta'x)^2=\frac43 f\cdot\A,\qquad
\theta_{\vartheta'}(\vartheta\vartheta'x)=-N.
\]

% inherited: noether_paper02_ko_translation_001_20260804\targets\Noether_P02_Korean_T09_U46_UNCHECKED.tex
% sha256: F2A520ABAD615F3FAA3D95075B4D0976BCA24E01D1C97437E273C18AE5766C24
\begin{center}
{\Large\bfseries 제III장.\quad $(s,\varrho,t)$를 법으로 하는 상대 완전계.}
\end{center}

\begin{center}
\textbf{\S\ 7.}\quad\emph{계 구성의 개요.}
\end{center}

$\nu$를 법으로 하는 상대 완전계에서 그다음으로 높은 모듈을 법으로 하는 상대 완전계로 넘어가려면, 정의 \S\ 3에 따라 이 더 높은 모듈에 관하여 $\nu$의 계를 구성하고 이를 $\nu$를 법으로 하는 상대 완전계의 형식들과 전이연산해야 한다. 그런데 $\nu=u_\nu^4$는 형식 $f$의 4차 반공변식으로서 $f$와 쌍대 관계에 있다. 따라서 $\nu$를 기본형으로 보면 $\nu(\nu)=(\nu\nu_1x)^4=s_x^4$를 법으로 하는 여섯 형식의 상대 완전계를 얻는다.

그러므로 $s$를 법으로 하는 상대 완전계(계 III)를 구성하려면 다음 두 계를 전이연산해야 한다.
\[
\begin{array}{ll}
\text{계 I:} & f,\ \theta,\ K,\ j,\ \A,\ N\quad \text{와}\vphantom{\dfrac11}\\[0.3em]
\text{계 II:} & \nu,\ H=(\nu\nu_1x)^2u_\nu^2u_{\nu_1}^2,\ L=(\nu\eta x)H_x^2u_\nu^3u_\eta^3,\ g=(\nu\eta x)^4H_x^2,\\
& \sigma=H_\eta^2u_\nu^2u_\eta^4,\ (\nu\sigma x).
\end{array}
\]
형식 $g$는 환원 가능하며, 따라서 계 II는 다섯 형식으로만 이루어진다는 것이 드러날 것이다(공식 (13.)).

% inherited: noether_paper02_ko_translation_001_20260804\targets\Noether_P02_Korean_T09_U47_UNCHECKED.tex
% sha256: 5598B17A075FBA6E17536A346C98853364B30198CC9C1F29E4D74B9F39CC97C6
계산 과정에서 이차형식들
\[
\varrho=u_\varrho^2=\theta_\nu^4u_\vartheta^2,\qquad t=t_x^2=a_\eta^4H_x^2
\]
을 모듈로 첨가하여 $(s,\varrho,t)$를 법으로 하는 상대 완전계를 얻는다. 이때 모듈 $(\varrho,t)$는 모듈 $s$보다 더 높은 모듈로 간주한다.

개별 형식들은 계수차수에 따라 배열한다.

수축
\[
\theta_\eta(K_\eta,\theta_\iota,\ldots)
\]
을 수축
\[
H_\vartheta(H_\kappa,L_\vartheta,\ldots).
\]
보다 더 높은 것으로 정규화한다. 이로써 \S\ 1과 \S\ 3b에 따라 같은 차수의 개별 형식들의 순서가 유일하게 정해진다.

% inherited: noether_paper02_ko_translation_001_20260804\targets\Noether_P02_Korean_T09_U48_UNCHECKED.tex
% sha256: A0D636A3F0E29E08C2AC05AB991CCCE58BB7BDC889F61651B40409AD2AD8DD23
같은 차수의 형식들은 표를 이용하여 구성한다.

I. 특정 계수차수를 얻기 위해 계 I과 계 II의 어떤 형식들을 서로 수축해야 하는지 제시한다.

II. 형식열의 도식에 따라 환원 불가능한 형식들을 제시한다.

III. 환원 가능하거나 분해 가능한 형식열 또는 형식을 환원한다. 즉 전체 형식열이 끊기는 위치를 밝히고, 이로써 제시한 것들이 모든 환원 불가능한 구성식을 소진한다는 것을 증명한다.

IV. 결론으로서 1) 새로 생긴 환원자들을 제시하고, 2) 같은 계의 형식들과 곱을 이루어 다른 계의 형식들과 수축하는 데 들어가지 않는 형식들을 제시한다.

% inherited: noether_paper02_ko_translation_001_20260804\targets\Noether_P02_Korean_T09_U49_UNCHECKED.tex
% sha256: 92D04D78360B67EA6A37CF65DD6E7898771CED8A311D991E41952131CBB90C88
다음 두 경우만 나타난다는 것이 드러날 것이다.

1) 계 I의 형식 하나와 계 II의 형식 하나를 각각 수축하는 경우.

2) $\theta$ 또는 $H$의 거듭제곱을 각각 수축하거나, 한 계의 두 형식을 다른 계의 형식 하나와 각각 수축하는 경우.

다음과 같은 수축 도식을 얻게 될 것이다.
\[
\begin{array}{r@{\ }c@{\ }l@{\qquad}c@{\qquad}r@{\ }c@{\ }l}
\multicolumn{3}{c}{\text{계 I}} & \text{와 수축} & \multicolumn{3}{c}{\text{계 II}}\\[0.25em]
1. & \text{차수} & f & & 2. & \text{차수} & \nu\\
2. & ,, & \theta & & 4. & ,, & H\\
3. & ,, & K,j,\A & & 6. & ,, & L,\sigma\\
4. & ,, & N,\theta^2,f\cdot j & & 8. & ,, & (\nu\sigma x),H^2\\
5. & ,, & \theta\cdot K & & 10. & ,, & H\cdot L\\
6. & ,, & \theta^3,\A\cdot j & & 12. & ,, & H^3.
\end{array}
\]

% inherited: noether_paper02_ko_translation_001_20260804\targets\Noether_P02_Korean_T09_U50_UNCHECKED.tex
% sha256: 26EEC7BE91A0A0D1D7B8DF97F606D64E9710C995E8CABA3388DFB86FC242EF6D
계산을 검산하기 위해 \glqq{}이중 환원\grqq{} 외에도 다음 두 특수 형식을 이용한다.
\[
\begin{array}{ll}
\text{I.}&\begin{array}{rlrlrl}
f&=\displaystyle\sum_3x_1^3x_2,& a_\nu^2&=\dfrac{i}{6}u_x^2,& s&=\dfrac{i}{3}f,\\
a_\eta^2&=-\dfrac19i\theta,& \A_\nu^2&=\dfrac{i}{6}\theta,& \sigma&=-\dfrac{i}{6}j,\\
g&=-\dfrac23i\A,& \varrho&=0,& t&=0,
\end{array}\\[1.2em]
\text{II.}&\begin{array}{rlrlrl}
f&=\displaystyle\sum_3x_1^4,& s&=\dfrac43if,& a_\eta^2&=\dfrac29i\theta,\\
\A_\nu^2&=\dfrac1{15}i\theta,& \sigma&=\dfrac43ij,& g&=\dfrac43i\A,\\
\varrho&=0,& t&=0.
\end{array}
\end{array}
\]

% inherited: noether_paper02_ko_translation_001_20260804\targets\Noether_P02_Korean_T09_U51_UNCHECKED.tex
% sha256: 772B236A3CE082FA7DA2F6FABCE7BD3724A784CF4AA687FC01585C725A39D22A
후기: 공식 (1.), (2.), (3.), (4.)에 대응하여 기본형 $\nu$에 대해서는 다음 공식들이 성립한다.
\[
\tag{5.}
\begin{array}{@{}c|c|c@{}}
\cdot & \underline{H_\nu=\dfrac13u_x\cdot\sigma} & \cdot\\[0.45em]
\underline{(\nu\eta x)^2}=A & H_\nu(\nu\eta x)=0 & \underline{H_\nu^2}=\sigma\\[0.45em]
\underline{(\nu\eta x)^3}=0 & \underline{H_\nu(\nu\eta x)^2}=B & \underline{H_\nu^2(\nu\eta x)}=0\\[0.45em]
\underline{(\nu\eta x)^4}=g & H_\nu(\nu\eta x)^3=0 & \underline{H_\nu^2(\nu\eta x)^2}=C,
\end{array}
\]
\[
\begin{aligned}
A&=\frac16s\cdot\nu-\frac23u_x\cdot s_\nu+\frac23u_x^2\cdot s_\nu^2-\frac{J}{18}u_x^4,\\
B&=-\frac56s_\nu+\frac76u_x\cdot s_\nu^2-\frac{J}{9}u_x^3,\\
C&=\frac23s_\nu^2-\frac{J}{9}u_x^2.
\end{aligned}
\]
\[
\tag{6.}
s_\nu^3u_\nu s_x=\frac13u_xs_\nu^4=\frac13J\cdot u_x.
\]

% inherited: noether_paper02_ko_translation_001_20260804\targets\Noether_P02_Korean_T09_U52_UNCHECKED.tex
% sha256: 29240A4BCAA7B5E2597685ED0E6A44CD43293D4D1FE3C7E0982CF4CB942070A5
\[
\tag{7.}
\begin{aligned}
\text{a)}\quad (\nu\sigma x)^2
&=\frac12(Hsu)^2-\frac25\nu\cdot s_\nu^2-\frac45u_x\cdot s_\eta(Hsu)^2+u_x^2\left\{\frac16J\cdot\nu-\frac1{40}(ss'u)^4\right\},\\
\text{b)}\quad (\nu\sigma x)^3&=-\frac3{10}s_\eta(Hsu),\\
\text{c)}\quad (\nu\sigma x)^4&=\frac{21}{10}s_\eta^2-\frac3{10}Z-\frac65u_x\cdot s_\eta^3+\frac1{10}u_x^2\cdot s_\eta^4.
\end{aligned}
\]

% inherited: noether_paper02_ko_translation_001_20260804\targets\Noether_P02_Korean_T09_U53_UNCHECKED.tex
% sha256: DA6F9A4E3DFB6CF7AA7729A82882254DE5B0EF897D4FDEC3B761CA80B6980C5D
\[
\tag{8.}
\begin{aligned}
\text{a)}\quad g'_\nu&=-s_\eta+u_x\left\{\frac92s_\eta^2-\frac12Z\right\}-3u_x^2\cdot s_\eta^3+\frac12u_x^3\cdot s_\eta^4,\\
\text{b)}\quad g_\nu^2&=\frac{21}{10}s_\eta^2-\frac3{10}Z-2u_x\cdot s_\eta^3+\frac12u_x^2\cdot s_\eta^4,\\
\text{c)}\quad g_\nu^3&=-\frac7{10}s_\eta^3+\frac12u_x\cdot s_\eta^4,\\
\text{d)}\quad g_\nu^4&=\frac35s_\eta^4.
\end{aligned}
\]

% inherited: noether_paper02_ko_translation_001_20260804\targets\Noether_P02_Korean_T10_U54_UNCHECKED.tex
% sha256: E49592E9335EFE7F42BFAED9A4E0853D6DFA48EB84F32A38DE0D9199CD369657
\begin{center}
\textbf{\S\ 8.}\quad\emph{계수차수 3의 형식들(계 III).}
\end{center}

\noindent\textbf{$f$와 $\nu$의 수축.}

\noindent\emph{환원 불가능 형식들:}
\[
\begin{array}{cccc}
a_\nu&\cdot&\cdot&\cdot\\
\cdot&a_\nu^2&\cdot&\cdot\\
\cdot&\cdot&\cdot&\cdot\\
\cdot&\cdot&\cdot&a_\nu^4=i.
\end{array}
\]

\noindent\emph{환원:}
\[
a_\nu^3=\frac13iu_x\qquad\text{(공식 (2.)).}
\]
\noindent\emph{결론:} 1) $a_\nu^3$은 환원자이다. 2) $f$는 반변 형식 $(j)$과의 곱 안에서만 $\nu$와 수축된다. $f$에 대한 $\nu$의 경우에도 마찬가지이다.

% inherited: noether_paper02_ko_translation_001_20260804\targets\Noether_P02_Korean_T11_U55_UNCHECKED.tex
% sha256: 926200615B8D59EC2A155146E0BC3E10A5333570FB08912ACBD85BE5B2BA28D3
\begin{center}
\textbf{\S\ 9.}\quad\emph{계수차수 4의 형식들(계 III).}
\end{center}

\noindent\textbf{$\theta$와 $\nu$의 수축.}

\noindent\emph{환원 불가능 형식들:}
\[
\begin{array}{cccc}
(\vartheta\nu x)&\theta_\nu&\cdot&\cdot\\
(\vartheta\nu x)^2&\theta_\nu(\vartheta\nu x)&\theta_\nu^2&\cdot\\
\cdot&\theta_\nu(\vartheta\nu x)^2&\cdot&\theta_\nu^3.
\end{array}
\]

\noindent\emph{환원:} 환원자 $a_\nu^3$으로부터 식 $a_\nu^3(abu)$, $a_\nu^3b_\nu$, $a_\nu^3b_\nu(abu)$을 다음과 같이 두고 이중 환원한다.

% inherited: noether_paper02_ko_translation_001_20260804\targets\Noether_P02_Korean_T11_U56_UNCHECKED.tex
% sha256: 1AA8A05A411DAD85219BF47DFC33AD164571951042CC8A5D83727BDFFE43DF84
\[
\begin{aligned}
\theta_\nu^2(\vartheta\nu x)
&=a_\nu^2(\widehat{ab}\nu x)(abu)-\frac13(\nu\nu_1x)^3
=a_\nu b_\nu(\widehat{ab}\nu x)(abu)+\frac16(\nu\nu_1x)^3\\
&=a_\nu^3(abu)-a_\nu^2b_\nu(abu)=2a_\nu^2b_\nu(abu)=\frac23a_\nu^3(abu),\\[.4em]
\theta_\nu^2(\vartheta\nu x)^2
&=a_\nu^2(\widehat{ab}\nu x)^2-\frac13(\nu\nu_1x)^4
=a_\nu b_\nu(\widehat{ab}\nu x)^2+\frac16(\nu\nu_1x)^4\\
&=i f-2a_\nu^3b_\nu+a_\nu^2b_\nu^2-\frac13s
=2a_\nu^3b_\nu-2a_\nu^2b_\nu^2+\frac16s
=\frac23i f-\frac23a_\nu^3b_\nu-\frac16s,\\[.4em]
\theta_\nu^3(\vartheta\nu x)&=a_\nu^2b_\nu(\widehat{ab}\nu x)(abu)=a_\nu^3b_\nu(abu).
\end{aligned}
\]

다음 공식들이 생긴다.
\[
\tag{9}
\begin{array}{c|c|c}
\underline{\theta_\nu^2(\vartheta\nu x)}=0
&\underline{\theta_\nu^3(\vartheta\nu x)}=0
&\underline{\theta_\nu^4u_\vartheta^2}=u_\varrho^2=\varrho\quad(\text{첨가 모듈, \S\ 7}),\\[0.6em]
\underline{\theta_\nu^2(\vartheta\nu x)^2}=\dfrac49 i f-\dfrac16s
&&
\end{array}
\]

\noindent\emph{결론:} 1) $\theta_\nu^2(\vartheta\nu x)^2$은 환원자이다.
2) $\nu$는 반변 형식 $(g)$과의 곱 안에서만 $\theta$와 수축된다.

% inherited: noether_paper02_ko_translation_001_20260804\targets\Noether_P02_Korean_T12_U57_UNCHECKED.tex
% sha256: 234E53B8B9199BC14BA4077D2365DEB91FD5F5BE1068AF0CDF80AF66418FE89F
\begin{center}
\textbf{\S\ 10.}\quad\emph{계수차수 5의 형식들(계 III).}
\end{center}

\[
\text{다음의 수축: }\left\{\begin{array}{l}
K,j,\A \text{와 }\nu,\\
f\text{와 }H.
\end{array}\right.
\]

\noindent\emph{환원 불가능 형식들:}
\[
\begin{array}{c@{\qquad}c@{\qquad}c@{\qquad}c}
\text{A)} & \begin{array}{cccc}
\cdot&K_\nu&\cdot&\cdot\\
(k\nu x)^2&\cdot&K_\nu^2&\cdot\\
(k\nu x)^3&K_\nu(k\nu x)^2&\cdot&\cdot\\
\cdot&K_\nu(k\nu x)^3&\cdot&\cdot
\end{array}
&
\text{B)}\ \begin{array}{c}(j\nu x)\\(j\nu x)^2\\(j\nu x)^3\\ \cdot\end{array}
&
\text{C)}\ \begin{array}{c}\A_\nu\\\cdot\\\cdot\\\cdot\end{array}
\end{array}
\]
\[
\text{D)}\qquad
\begin{array}{cccc}
(aHu)&(aHu)^2&\cdot&\cdot\\
a_\eta&a_\eta(aHu)&a_\eta(aHu)^2&\cdot\\
\cdot&a_\eta^2&\cdot&\cdot
\end{array}
\]

% inherited: noether_paper02_ko_translation_001_20260804\targets\Noether_P02_Korean_T12_U58_UNCHECKED.tex
% sha256: 00492270E77478BD6ED80CD05058C9C5DDF32F50F54766A1262E3E10B68F8ED8
\noindent\emph{환원:}

\noindent A) 형식열 $K_\nu^3$: 환원자 $a_\nu^3$.

형식 $K_\nu^2(k\nu x)^2$: 환원자 $\theta_\nu^2(\vartheta\nu x)^2$.
\[
(k\nu x),\quad K_\nu(k\nu x),\quad K_\nu^2(k\nu x)
\]
은 \S\ 3에 따라 분해 가능 형식이다.

% inherited: noether_paper02_ko_translation_001_20260804\targets\Noether_P02_Korean_T12_U59_UNCHECKED.tex
% sha256: 7ED11389ACA5A89B2A8423FD03CEF0ED58E8C62987B00080DB90289A24ACB61C
\noindent B) 형식
\[
\begin{aligned}
(j\nu x)^4&=(\widehat{a\theta}\nu x)^4
=a_\nu^4\theta+\theta_\nu^4 f-4a_\nu^3\theta_\nu
-4a_\nu\{a_\nu\theta_x-(\widehat{a\theta}\nu x)\}^3\\
&\quad+6a_\nu^2\{a_\nu\theta_x-(\widehat{a\theta}\nu x)\}^2
=\R f+\frac53 i\theta-6a_\nu^2(\widehat{a\theta}\nu x)^2
+4a_\nu(\widehat{a\theta}\nu x)^3,
\end{aligned}
\]
이고, 따라서
\[
(j\nu x)^4=\R f+\frac53 i\theta\modu{(\A,\nu)},\qquad (\text{\S\ 5 끝부분 참조}).
\]

% inherited: noether_paper02_ko_translation_001_20260804\targets\Noether_P02_Korean_T12_U60_UNCHECKED.tex
% sha256: 16790D69C41F065C371648F9346A0798D53D355DAC445BD4A1B9A34FDCFDD892
\noindent C) 형식 $\A_\nu^4$: 환원자 $\theta_\nu^4$.

형식 $\A_\nu^2$와 $\A_\nu^3$: 형식 $\A$를 형식열의 더 낮은 형식들로 바꾸어 얻는 환원자 $a_\nu^3$ 또는 $\theta_\nu^4$(\S\ 5 끝부분), 즉 식들
\[
a_\vartheta a_\nu^3,\qquad
a_\vartheta a_\nu^3\theta_\nu,\qquad
a_\vartheta\theta_\nu^4.
\]
의 이중 환원으로 얻는다. $\A_\nu^3$을 두 가지로 계산하면 더 높은 형식 $a_\eta^3$을 환원할 계기가 생긴다.

% inherited: noether_paper02_ko_translation_001_20260804\targets\Noether_P02_Korean_T12_U61_UNCHECKED.tex
% sha256: A456657D7EDAC80843A431E3E457169C6DD92B0AB0DCCD3A32F97BACE4F7672F
\noindent\emph{환원 공식:}
\[
\tag{10}
\begin{aligned}
\text{a)}\quad \A_\nu^2
&=\frac35a_\eta^2-\frac3{10}(asu)^2+\frac13 i\theta
+\frac15u_x^2(2a_\R^2-t),\\
\text{b)}\quad \A_\nu^3
&=-\frac3{10}a_\R+\frac35u_x\left(\frac{13}{10}a_\R^2-\frac25t\right),\\
\text{c)}\quad \A_\nu^4&=a_\R^2-\frac25t.
\end{aligned}
\]

% inherited: noether_paper02_ko_translation_001_20260804\targets\Noether_P02_Korean_T12_U62_UNCHECKED.tex
% sha256: A1035AAEBF0934EF6BAA484E61BC99C606A097CCA136C4738A431E2903A74FF2
\noindent\emph{공식의 유도:}
\[
\begin{aligned}
\text{a)}\quad
 a_\vartheta a_\nu^3
&=\frac13 i\theta
=[a_\vartheta]_{\nu^3}+\frac67[a_\vartheta^2]_{\nu^2}
+\frac65[a_\vartheta(a\theta u)^2]_{\nu^2}\,\widehat{\nu}x^2
+\frac{11}{30}[a_\vartheta^2(a\theta u)^2]_{\nu}\,\widehat{\nu}x^2\\
&\quad-\frac67u_x[a_\vartheta^2]_{\nu^3}
-\frac16u_x[a_\vartheta^2(a\theta u)^2]_{\nu^2}\,\widehat{\nu}x^2,\\
\frac13 i\theta
&=\A_\nu^2-\frac23u_x\A_\nu^3-a_\nu a_{\nu_1}(\nu\nu_1x)^2
+\frac25a_\nu^2(\nu\nu_1x)^2
+\frac15u_xa_\nu^2a_{\nu_1}(\nu\nu_1x)^2;
\end{aligned}
\]

% inherited: noether_paper02_ko_translation_001_20260804\targets\Noether_P02_Korean_T12_U63_UNCHECKED.tex
% sha256: A159E06DE772751B48DF1B464E7799C98227FFD5B3B2BB827BB9B2A20C9FF7CF
\[
\begin{aligned}
\text{b)}\quad
 a_\vartheta a_\nu^3\theta_\nu
&=0=[a_\vartheta]_{\nu^4}+\frac9{14}[a_\vartheta^2]_{\nu^3}
+\frac35[a_\vartheta(a\theta u)^2]_{\nu^2}\,\widehat{\nu}x^2
+\frac{17}{120}[a_\vartheta^2(a\theta u)^2]_{\nu}\,\widehat{\nu}x^2\\
&\quad-\frac9{14}u_x[a_\vartheta^2]_{\nu^4}
-\frac1{24}u_x[a_\vartheta^2(a\theta u)^2]_{\nu^2}\,\widehat{\nu}x^2,\\
0&=\frac56\A_\nu^3-\frac12u_x\A_\nu^4-\frac14a_\eta^3+\frac1{20}u_xa_\eta^4;
\end{aligned}
\]

% inherited: noether_paper02_ko_translation_001_20260804\targets\Noether_P02_Korean_T12_U64_UNCHECKED.tex
% sha256: DEF1101C0B34B1C6AE7FD82373DA45FB5FBB8AF1ABAFDAD9FD59BA9656CB27D5
\[
\begin{aligned}
\text{b')}\quad
 a_\vartheta\theta_\nu^4
&=a_\R=a_\vartheta\theta_\nu\{a_\nu-(a\theta\widehat{\nu}x)\}^3\\
&=-\frac32[a_\vartheta^2]_{\nu^3}
+\frac35[a_\vartheta(a\theta u)^2]_{\nu^2}\,\widehat{\nu}x^2
-\frac{37}{120}[a_\vartheta^2(a\theta u)^2]_{\nu}\,\widehat{\nu}x^2\\
&\quad+\frac32u_x[a_\vartheta^2]_{\nu^4}
+\frac{49}{120}u_x[a_\vartheta^2(a\theta u)^2]_{\nu^2}\,\widehat{\nu}x^2,\\
a_\R&=-\frac32\A_\nu^3+\frac32u_x\A_\nu^4-\frac{11}{20}a_\eta^3+\frac7{20}u_xa_\eta^4;
\end{aligned}
\]

% inherited: noether_paper02_ko_translation_001_20260804\targets\Noether_P02_Korean_T12_U65_UNCHECKED.tex
% sha256: FDE71068EC508A8FDAE2706B4BE0231C59ED1B4EE74BD3FFB8152E364B92F2D0
\[
\begin{aligned}
\text{c)}\quad
 a_\vartheta^2\theta_\nu^4
&=a_\R^2=[a_\vartheta^2]_{\nu^4}+\frac35[a_\vartheta^2(a\theta u)^2]_{\nu^2}\,\widehat{\nu}x^2,\\
a_\R^2&=\A_\nu^4+\frac25a_\eta^4.
\end{aligned}
\]

\noindent D) C)의 $\A_\nu^3$을 이중 환원하여 $a_\eta^3$을 환원한다.

나머지 환원들은 \S\ 9의 계산과 쌍대인 계산으로 얻는다.

% inherited: noether_paper02_ko_translation_001_20260804\targets\Noether_P02_Korean_T12_U66_UNCHECKED.tex
% sha256: 6DDEBA55C921A90381E8DF040AFDCABBF121CC5BB7216533CDF3B6636E3ABBFF
\noindent\emph{환원 공식:}
\[
\tag{11}
\begin{array}{rl|rl}
\underline{a_\eta^2(aHu)}&=-\dfrac16(asu)^3,&
\underline{a_\eta^2(aHu)^2}&=\dfrac49 i \nu-\dfrac16(asu)^4,\\[0.7em]
\underline{a_\eta^3}&=-a_\R+\dfrac35u_xa_\R^2+\dfrac15u_xt,
&\underline{a_\eta^3(aHu)}&=0,\\[0.7em]
&&\underline{a_\eta^4}&=t\quad(\text{모듈로 첨가됨}).
\end{array}
\]

\noindent\emph{결론:}

1) $(j\nu x)^4$, $\A_\nu^2$, $a_\eta^2(aHu)$는 환원자이다.

2) $\nu$는 반변 형식들과의 곱 안에서만 $K,j,\A$와 수축된다. $f$가 $H$와 수축되는 경우와 $j$가 $\nu$와 수축되는 경우에도 마찬가지다. 왜냐하면 $\theta_\nu(\vartheta\cdot j,\nu,x)^3$이 \S\ 11 B에 따라 환원 가능해지기 때문이다.

% inherited: noether_paper02_ko_translation_001_20260804\targets\Noether_P02_Korean_T13_U67_UNCHECKED.tex
% sha256: EFCBF90B4FF0C62FCB379C44A036B365DAE65CC17CA2E47538702DA07D65C0FB
\begin{center}
\textbf{\S\ 11.\quad 계수차수 6의 형식들(계 III).}
\end{center}
\[
\text{다음의 수축: }\left\{\begin{array}{l}
\theta^2,\ f\cdot j,\ N \text{과 }\nu,\\
\theta\text{와 }H.
\end{array}\right.
\]


\noindent\emph{환원 불가능 형식들:}
\[
\renewcommand{\arraystretch}{1.25}
\begin{array}{@{}r@{\quad}c@{\qquad\qquad}r@{\quad}c@{\qquad\qquad}r@{\quad}c@{}}
\text{A)}&(\vartheta^2\nu x)^3&\text{B)}&a_\nu(j\nu x)&\text{C)}&N_\nu,\\[0.45em]
&\cdot&&\cdot&&\\[-0.1em]
&\theta_\nu(\vartheta^2\nu x)^3&&a_\nu^2(j\nu x)&&
\end{array}
\]
\[
\renewcommand{\arraystretch}{1.25}
\begin{array}{@{}r@{\quad}c|c|c|c@{}}
&\cdot&(\theta Hu)&(\theta Hu)^2&\cdot\\
\text{D)}&(\vartheta\eta x)&\theta_\eta&H_\vartheta(\theta Hu),\ \theta_\eta(\theta Hu)&\theta_\eta(\theta Hu)^2\\
&(\vartheta\eta x)^2&H_\vartheta(\vartheta\eta x),\ \theta_\eta(\vartheta\eta x)&\theta_\eta^2&\theta_\eta^2(\theta Hu)\\
&\cdot&\theta_\eta(\vartheta\eta x)^2&\cdot&\cdot
\end{array}
\]

% inherited: noether_paper02_ko_translation_001_20260804\targets\P02_T13_U68_r2.tex
% sha256: A6B87E38855EDA0656D6BDE1FB1B9212DC3842465209527073C2C42D5377EDEC
\noindent\emph{환원:}

\noindent A) $(\vartheta^2\nu x)^4$는 공식 (3.)a에 따라 분해된다.
\[
(a\A\widehat{\vartheta}x)^2=\frac12(\vartheta^2\nu x)^4-\frac35 f\cdot a_\nu^2(\nu\vartheta x)^2
=\frac13\A^2+f\A_\vartheta^2.
\]

\noindent B) $a_\nu(j\nu x)^2$은 공식 (9.)a에 따라 분해된다. 이때 \S\ 6 끝부분에 따라 $f\cdot j$를 $(\theta\theta'u)^2$으로 바꾸고 $\theta_\vartheta'$의 환원 가능성을 이용한다.
\[
\frac13a_\nu(j\nu x)^2=(\widehat{\theta\theta'}\nu x)^2\theta_\nu
=\theta_\nu^3\theta-\theta_\nu^2\theta_\nu'
=\theta_\nu^3\theta+\theta_\nu^2(\widehat{\vartheta}\nu\theta'u)
=\theta_\nu^3\theta-\frac23\theta_\nu^3\theta.
\]
형식 $a_\nu(j\nu x)^3$과 $a_\nu^2(j\nu x)^2$: 환원자 $a_j$와 $(j\nu x)^4$.
\[
a_\nu(j\nu x)^3=a_j(j\nu x)^3-(j\nu x)^3(j\nu\widehat{a}u),
\]
\[
a_\nu^2(j\nu x)^2=a_j^2(j\nu x)^2-2a_j(j\nu x)^2(j\nu\widehat{a}u)+(j\nu x)^2(j\nu\widehat{a}u)^2.
\]

\noindent C) 형식열 $N_\nu^2$: 환원자 $\A_\nu^2$.
$(n\nu x)$와 $N_\nu(n\nu x)$는 \S\ 3에 따라 야코비 행렬식에 대한 전이연산으로서 분해된다.

% inherited: noether_paper02_ko_translation_001_20260804\targets\Noether_P02_Korean_T13_U69_UNCHECKED.tex
% sha256: FDE0C22AF2CBF8A0BD0DCF8217EF8DCCFC34DC66EC3CB006DDA5A200770F9875
\noindent D) \emph{환원의 개요:}

a) 형식열 $\theta_\eta^2H_\vartheta$: 환원자 $a_\eta^2(aHu)$. (더 높은 기호를 갖는 형식들로 이어진다.)

b) 형식열 $\theta_\eta H_\vartheta$: 이중 환원을 통한 환원자 $a_\eta^2(aHu)$.
(더 높은 기호를 갖는 형식들과 전체 형식열의 더 높은 형식들, 즉 형식열 $\theta_\eta^2$로 이어진다.) $\theta_\eta H_\vartheta$ 대신 형식열
\[
\theta_\eta(\vartheta\eta x)(\theta Hu)=\theta_\eta H_\vartheta+\theta_\eta^2
\]
을 고찰한다. 여기서 기호 $\theta$를 더 낮은 형식 $(abu)$로 바꾸면 형식열
\[
a_\eta^2(aHu)(abu)=\theta_\eta H_\vartheta+\frac32\theta_\eta^2\modu{s}
\]
의 이중 환원을 거쳐 끝형식
\[
a_\eta^3b_\eta(bHu)(abu)=\theta_\eta^3H_\vartheta+\frac12\theta_\eta^4
\]
에 이른다.

% inherited: noether_paper02_ko_translation_001_20260804\targets\Noether_P02_Korean_T13_U70_UNCHECKED.tex
% sha256: B357589DD8346506285B40B027755BC4FDF87AA1C96151BD210AB3C0A6FC1098
c) 형식열 $\theta_\eta^3$: 형식열 $\theta_\eta H_\vartheta$와 $\theta_\eta^2H_\vartheta$에 동시에 속하는 형식들, 즉 형식열 $\theta_\eta^2H_\vartheta$의 형식들을 이중 환원하여 얻는다. 한편으로는 더 높은 기호를 갖는 형식들로 환원하고, 다른 한편으로는 더 높은 기호를 갖는 형식들과 더 높은 형식열 $\theta_\eta^3$으로 환원한다.

d) 형식 $\theta_\eta^2(\vartheta\eta x)$, $\theta_\eta^2(\vartheta\eta x)^2$, $\theta_\eta^3(\vartheta\eta x)$: \S\ 9의 계산과 마찬가지로 환원자 $a$를 이용한 이중 환원으로 얻는다.

e) 형식 $\theta_\eta^2(\theta Hu)^2$와 $\theta_\eta^2(\vartheta\eta x)^2$: 환원자 $\A_\nu^2$와 $(a\A u)^4$를 이용하여 식 $\A_\nu^2(a\A u)^4$와 $(a\A\widehat{\nu}x)^4$을 이중 환원함으로써 얻는다.

f) 형식 $H_\vartheta$와 $H_\vartheta^2$: 분해 가능 형식이다. $H_\vartheta^2$의 경우 이는 식 $\A_\nu^2(a\A u)^2$의 이중 환원으로, 또는 곱 $(a_\nu^2)^2$, $a_\nu^2\cdot b_{\nu_1}$을 계의 형식들로 직접 계산함으로써 얻는다. ($(a_\nu)^2$을 이와 같이 계산하면 분해 가능 형식들 사이의 관계가 생긴다.)

g) 형식열 $\theta\cdot H$의 형식들이 두 가지 환원 방법을 허용하기 때문에, 즉 두 환원 가능한 형식열에 속하기 때문에 다음과 같이 더 높은 형식들이 환원된다.
\[
\begin{array}{rl}
g&\text{는 }\theta_\eta^2(\vartheta\eta x)^2\text{의 이중 환원으로: 환원 d)와 e)가 가능하다;}\\[0.3em]
s_\vartheta^2(\theta s u)&\text{는 }\theta_\eta^3(\vartheta\eta x)\text{의 이중 환원으로: 환원 c)와 d)가 가능하다;}\\[0.3em]
s_\vartheta^2(\theta s u)^2&\text{은 }\theta_\eta^2H_\vartheta^2\text{의 이중 환원으로: 환원 a)가 가능하며}\\
&\text{환원자 }\theta_\nu^2(\vartheta\nu x)^2\text{을 인수로 갖는다;}\\[0.3em]
s_\nu^2&\text{은 }\theta_\eta^3H_\vartheta\text{의 이중 환원으로: 환원 a)와 c)가 가능하다.}
\end{array}
\]
나머지 형식들의 독립성은 형식열 $\A_\nu^2(a\A u)^2\equiv\theta_\eta^2$의 이중 환원으로 증명된다.

% inherited: noether_paper02_ko_translation_001_20260804\targets\Noether_P02_Korean_T13_U71_UNCHECKED.tex
% sha256: 45F69FB3A2136067C8C7386315C0768F9A7083AD22E10B25605D07E897699E3B
\noindent\emph{환원의 실행:}

제시한 구성 법칙에 따르면 이중 환원에서 생기는 관계들이 서로 독립이므로, 환원 a), b), c), d)의 존재는 선험적으로 명백하다. 환원 e), f), g)에서는 항등식이 생기지 않는다는 것을 계산으로 증명해야 한다.

형식열 $\theta\cdot H$의 대각 성분에 대칭을 이루도록 형식 $\theta_\eta$까지 나아가면서, 환원 a), b), c), d)로 얻은 공식들도 일부 계산의 개요와 함께 제시한다. 이 공식들은 일부는 환원 e), f), g)에, 일부는 뒤에서 사용되기 때문이다.

\noindent 1) f)에 따른 $H_\vartheta$와 $H_\vartheta^2$. 다음과 같이 둔다\srcfn{*)}{등호 대신 쓴 기호 $\sim$은 수축 III과 IV에서 생기는 형식들보다 더 높은 형식에 $u_x$를 곱한 곱들을 생략했다는 뜻이다.}.
\[
\begin{aligned}
a_\nu\cdot b_{\nu_1}^2
&=\nu\cdot a_\nu b_\nu^2-(aHu)(bHu)b_\eta
\sim \nu\cdot a_\nu b_\nu^2-f\cdot a_\eta(aHu)^2-(abu)(aHu)b_\eta+(abu)^2b_\eta,\\
a_\nu^2\cdot b_{\nu_1}^2
&=b_{\nu_1}^2\{a_{\nu_1}u_\nu-(a_\nu\widehat{\nu}_1u)\}^2
\sim \nu\cdot a_\nu^2b_\nu^2-2a_\eta b_\eta(aHu)(bHu)-\frac16(asu)^2(bsu)^2\\
&\sim \nu\cdot a_\nu^2b_\nu^2-2a_\nu^2(aHu)(abu)-\frac16(asu)^2(bsu)^2+(\vartheta\eta x)^2(\theta Hu)^2.
\end{aligned}
\]
$\theta_\eta H_\vartheta$의 값을 대입하면 다음 공식들이 생긴다.
\[
\tag{12}
\begin{aligned}
\text{a)}\quad H_\vartheta&\sim 2a_\nu\cdot a_\nu^2+2\nu\cdot\theta_\nu(\vartheta\nu x)^2+2f\cdot a_\eta(aHu)^2-2\theta_\eta,\\
\text{b)}\quad H_\vartheta^2&\sim (a_\nu^2)^2+2\theta_\eta^2+\frac23(\theta su)^2+\frac1{12}s\cdot\nu-\frac19 i f\cdot\nu-\frac16f\cdot(asu)^4.
\end{aligned}
\]

% inherited: noether_paper02_ko_translation_001_20260804\targets\Noether_P02_Korean_T13_U72_UNCHECKED.tex
% sha256: C833173AA31351BA95E54BBFEE3CC47925D19AA9323DEEE9B2D6FD8D11E8D933
\noindent 2) b)에 따른 $\theta_\eta H_\vartheta$:
\[
\begin{aligned}
a_\eta^2(aHu)(abu)
&\sim [a_\eta^2(aHu)]_{\widehat{bu}}+\frac12 f\cdot[a_\eta^2(aHu)^2]
=a_\eta b_\eta(aHu)(abu)\\
&\quad+a_\eta(\widehat{ab}\eta x)(abu)(aHu)-\theta_\eta(\vartheta\eta x)(\theta Hu)+\frac12\theta_\eta^2+\frac14(\nu\eta x)^2.
\end{aligned}
\]
공식 (5.)와 (11.)에 따라 다음이 성립한다.
\[
\theta_\eta H_\vartheta\sim-\frac32\theta_\eta^2-\frac14(\theta su)^2-\frac1{12}s\cdot\nu+\frac29 i f\cdot\nu.
\]

\noindent 3) $\theta_\eta H_\vartheta(\theta Hu)$는 b)에 따라 $\theta_\eta H_\vartheta$와 같은 방식으로 계산된다.
\[
\theta_\eta H_\vartheta(\theta Hu)\sim-\theta_\eta^2(\theta Hu)-\frac16(\theta su)^3.
\]

\noindent 4) b)에 따른 $\theta_\eta H_\vartheta(\vartheta\eta x)$; d)에 따른 $\theta_\eta^2(\vartheta\eta x)$(\S\ 9 참조):
\[
\theta_\eta H_\vartheta(\vartheta\eta x)\sim-\theta_\eta^2(\vartheta\eta x)-\frac16s_\vartheta(\theta su)
\sim\frac13(\vartheta\R x)-\frac16s_\vartheta(\theta su),
\]
\[
\theta_\eta^2(\vartheta\eta x)\sim\frac23a_\eta^3(abu)\sim-\frac13(\vartheta\R x).
\]

% inherited: noether_paper02_ko_translation_001_20260804\targets\Noether_P02_Korean_T13_U73_UNCHECKED.tex
% sha256: 23AA3A299BE004F69952280566F7FF3C760BFBDF9F42DE212477C5F8DC8777F0
\noindent 5)
\[
\begin{array}{ccc}
\theta_\eta H_\vartheta^2\text{은 b)에 따라},&
\theta_\eta^2H_\vartheta\text{는 a)에 따라},&
\theta_\eta^2H_\vartheta\text{는 b)에 따라, 이로써 }\theta_\eta^3\text{은 c)에 따라},\\
a_\eta^2(aHb)(bHu)\text{에서};&
a_\eta^2(Hab)(abu)\text{에서};&
a_\eta^3(bHu)(abu)\text{에서}:
\end{array}
\]
\[
\theta_\eta H_\vartheta^2\sim-\frac32\theta_\eta^2H_\vartheta-\frac14s_\vartheta(\theta su)^2+\frac16s_\nu-\frac49ia_\nu
\sim-\theta_\varrho-\frac16s_\nu-\frac19ia_\nu,
\]
\[
\theta_\eta^2H_\vartheta\sim\frac23\theta_\varrho-\frac16s_\vartheta(\theta su)^2+\frac29s_\nu-\frac29ia_\nu,
\]
\[
\theta_\eta^3H_\vartheta\sim\frac23\theta_\varrho-\frac23\theta_\eta^3+\frac5{36}s_\nu\quad\text{이고, 따라서}
\]
\[
\theta_\eta^3\sim\frac14s_\vartheta(\theta su)^2-\frac18s_\nu+\frac13ia_\nu.
\]

% inherited: noether_paper02_ko_translation_001_20260804\targets\Noether_P02_Korean_T13_U74_UNCHECKED.tex
% sha256: 5BC6B1ECFF03E9FEA0632528C1E8BCB2877D2A5BF0A7F0FDC82AA2FAE1EDD6BB
\noindent 6) e)에 따른 $\theta_\eta^2(\theta Hu)^2$:
\[
\A_\nu^2(a\A u)^4=[(a\A u)^4]_{\nu^2}=-\frac{12}{5}\theta_\eta^2(\theta Hu)^2-\frac3{10}\sigma-\frac1{20}(\theta su)^4+\nu\cdot\varrho
\quad(\text{공식 (3.) c에 따라}),
\]
\[
\A_\nu^2(a\A u)^4=[\A_\nu^2]^4_{\widehat{u a}}=\frac35\theta_\eta^2(\theta Hu)^2+\frac15\sigma-\frac3{10}(\theta su)^4+\frac13ij
\quad(\text{공식 (10.) a에 따라}),
\]
\[
\theta_\eta^2(\theta Hu)^2=-\frac16\sigma+\frac1{12}(\theta su)^4-\frac19ij+\frac13\nu\cdot\varrho.
\]

% inherited: noether_paper02_ko_translation_001_20260804\targets\Noether_P02_Korean_T13_U75_UNCHECKED.tex
% sha256: 7A3408920D13A85C4EE275C56E1EC18FDEDA126067526E5A1ABDEE86C587A7E6
\noindent 7) d)와 e)에 따른 $\theta_\eta^2(\vartheta\eta x)^2$. 이로부터 더 높은 형식 $g$를 환원한다.

\S\ 9와 유사하게 계산하면 다음이 성립한다.
\[
\begin{aligned}
\theta_\eta^2(\vartheta\eta x)^2
&=\frac12t\cdot f-\frac12a_\eta^2b_\eta-\frac1{12}g=\frac23t\cdot f-\frac23a_\eta^3b_\eta-\frac16g\\
&=-\frac16g-\frac13(\vartheta\varrho x)^2+\frac4{15}f\cdot a_\varrho^2+\frac8{15}f\cdot t\quad(\text{공식 (11.)에 따라}).
\end{aligned}
\]

% inherited: noether_paper02_ko_translation_001_20260804\targets\Noether_P02_Korean_T13_U76_UNCHECKED.tex
% sha256: 29209ADAE4B21829C4E973EB6BC05A7557EE13C7229508B6A6472CE0FE61747F
위와 같은 방식으로 대응하는 계산을 하면 다음이 성립한다.
\[
(a\A\widehat{\nu}x)^4=[(a\A u)^4]_{\widehat{\nu}x^4}=\frac{21}{10}\theta_\eta^2(\vartheta\eta x)^2+\frac7{10}s_\vartheta^2-\frac3{10}g
\quad(\text{공식 (3.) c에 따라}),
\]
\[
\begin{aligned}
(a\A\widehat{\nu}x)^4&=-\frac13i\A+6[\A_\nu^2]a^2-4[\A_\nu^3]a+\A_\nu^4\cdot f\\
&=-\frac{36}{5}\theta_\eta^2(\vartheta\eta x)^2-\frac95s_\vartheta^2-\frac35g-\frac35(\vartheta\varrho x)^2\\
&\quad+\frac53i\A+\frac{37}{25}f\cdot a_\varrho^2+\frac{74}{25}f\cdot t\quad(\text{공식 (10.)에 따라}).
\end{aligned}
\]
따라서
\[
\begin{aligned}
\frac{93}{10}\theta_\eta^2(\vartheta\eta x)^2
&=-\frac3{10}g-\frac52s_\vartheta^2-\frac35(\vartheta\varrho x)^2\\
&\quad+\frac53i\A+\frac{37}{25}f\cdot a_\varrho^2+\frac{74}{25}f\cdot t,
\end{aligned}
\]
이고, g)에 따라 $\theta_\eta^2(\vartheta\eta x)^2$의 두 값을 비교하면
\[
\tag{13.)\srcfnmark{*)}}
0=g-2s_\vartheta^2+2(\vartheta\varrho x)^2+\frac43i\A-\frac45f\cdot a_\varrho^2-\frac85f\cdot t.
\]
\srcfntext{*)}{이는 서론의 주석에서 언급한 계수차수 6, 변수차수 6인 세 공변식 사이의 관계이다.}

% inherited: noether_paper02_ko_translation_001_20260804\targets\Noether_P02_Korean_T13_U77_UNCHECKED.tex
% sha256: EE0DAB16FBE8C00814E6879CC4044622E9D097180CB28B1D17109D4FC23FE15B
\noindent 8) a)와 b)에 따른 $\theta_\eta^2H_\vartheta(\theta Hu)$, 이로부터 c)에 따른 $\theta_\eta^3(\theta Hu)$(인수 $u_x$를 갖는 형식들은 소멸한다):
\[
\theta_\eta^2H_\vartheta(\theta Hu)=-\frac16s_\vartheta(\theta su)^3-\frac13(\nu\varrho x),
\]
\[
\theta_\eta^2H_\vartheta(\theta Hu)=-\frac13\theta_\eta^3(\theta Hu)-\frac13(\nu\varrho x),
\]
\[
\theta_\eta^3(\theta Hu)=\frac12s_\vartheta(\theta su)^3.
\]

% inherited: noether_paper02_ko_translation_001_20260804\targets\Noether_P02_Korean_T13_U78_UNCHECKED.tex
% sha256: B921CE0D2CAFFABC712401672415563F1660D37D5ED210D144FD559969F3EAD8
\noindent 9) a)와 b)에 따른 $\theta_\eta^2H_\vartheta(\vartheta\eta x)$, 이로부터 c)에 따른 $\theta_\eta^3(\vartheta\eta x)$. d)에 따른 $\theta_\eta^3(\vartheta\eta x)$, 이로부터 g)에 따른 더 높은 형식 $s_\vartheta^2(\theta su)$의 환원(인수 $u_x$를 갖는 형식들은 소멸한다):
\[
\theta_\eta^2H_\vartheta(\vartheta\eta x)=\frac13(atu),
\]
\[
\theta_\eta^2H_\vartheta(\vartheta\eta x)=\frac13(atu)+\frac13\theta_\eta^3(\vartheta\eta x)-\frac16s_\vartheta^2(\theta su),
\]
\[
\theta_\eta^3(\vartheta\eta x)=\frac12s_\vartheta^2(\theta su),
\]
\[
\theta_\eta^3(\vartheta\eta x)=\frac25\theta_\varrho(\vartheta\varrho x)+\frac15(atu):
\]
\srcnumdisplay{(14.)}{%
s_\vartheta^2(\theta su)=\frac45\theta_\varrho(\vartheta\varrho x)+\frac25(atu).
}

% inherited: noether_paper02_ko_translation_001_20260804\targets\Noether_P02_Korean_T13_U79_UNCHECKED.tex
% sha256: E210AC65F0B3AECC05C3C3F4DEC8E01413E68D90ACF0CF64755BB6E7BE41FCB0
\noindent 10) a)에 따른 $\theta_\eta^2H_\vartheta^2$ 및 환원자 $\theta_\nu^2(\vartheta\nu x)^2$을 통한 $\theta_\eta^2H_\vartheta^2$, a)와 c)에 따른 $\theta_\eta^3H_\vartheta$, c)에 따른 $\theta_\eta^4$. 이로부터 g)에 따라 더 높은 형식 $s_\vartheta^2(\theta su)^2$과 $s_\nu^2$을 환원한다.
\[
\begin{aligned}
\text{a)에 따라}\quad
\theta_\eta^2H_\vartheta^2&=[a_\eta^2(aHu)^2]b^2+\frac13[a_\eta^4]_{bu^2}-\frac13H_\nu^2(\nu\eta x)^2\\
&=\frac49ia_\nu^2-\frac16s_\vartheta^2(\theta su)^2-\frac5{18}s_\nu^2+\frac13(atu)^2+\frac1{54}J\cdot u_x^2.
\end{aligned}
\]
\[
\begin{aligned}
\text{환원자 }\theta_\nu^2(\vartheta\nu x)^2\text{을 통해}\quad
\theta_\eta^2H_\vartheta^2
&=[\theta_\nu^2(\vartheta\nu x)^2]_{\nu_1}+\frac13[\theta_\nu^4]_{\nu_1}\widehat{\nu}_1x^2-\frac13s_\vartheta^2(\theta su)^2\\
&=\frac49ia_\nu^2-\frac16s_\nu^2-\frac13s_\vartheta^2(\theta su)^2+\frac13(\nu\varrho x)^2.
\end{aligned}
\]

% inherited: noether_paper02_ko_translation_001_20260804\targets\Noether_P02_Korean_T13_U80_UNCHECKED.tex
% sha256: 40E2F2FA110EA609EB91CE2C5A0C273F4E48BB0FE1A99B45DBC973454F7F3379
\[
\begin{aligned}
\text{a)에 따라}\quad
\theta_\eta^3H_\vartheta&=-[a_\eta^2(aHu)]_{bu}b^2-\frac12[a_\eta^2(aHu)^2]b^2-\frac13[a_\eta^3]b+\frac1{12}[a_\eta^4]_{bu^2}\\
&\quad+\frac12u_x[a_\eta^2(aHu)^2]b^3\\
&=-\frac29ia_\nu^2+\frac1{12}s_\nu^2+\frac4{15}\theta_\varrho^2-\frac7{90}(\nu\varrho x)^2+\frac1{30}(atu)^2+\frac2{27}i^2u_x^2-\frac1{36}J\cdot u_x^2.
\end{aligned}
\]
\[
\begin{aligned}
\text{c)에 따라}\quad
\theta_\eta^3H_\vartheta:\quad a_\eta^3(Hab)(bHu)&=\frac12(atu)^2
=\theta_\eta^2H_\vartheta(\theta Hu)(\vartheta\eta x)+\frac14H_\nu^2(\nu\eta x)^2\\
&\quad+\frac12\theta_\eta^2H_\vartheta^2
=\frac32\theta_\eta^2H_\vartheta^2+\theta_\eta^3H_\vartheta-u_x[a_\eta^2(aHu)^2]b^3\\
&\quad+\frac14H_\nu^2(\nu\eta x)^2,
\end{aligned}
\]
따라서
\[
\begin{aligned}
\theta_\eta^3H_\vartheta&=-\frac23ia_\nu^2+\frac5{12}s_\nu^2+\frac12(\nu\varrho x)^2-\frac12(atu)^2\\
&\quad+\frac4{27}i^2u_x^2-\frac1{12}J\cdot u_x^2.
\end{aligned}
\]

% inherited: noether_paper02_ko_translation_001_20260804\targets\Noether_P02_Korean_T13_U81_UNCHECKED.tex
% sha256: A8EB87301E45F7B888884F8E4F437BBDAA770A574AD696AF3ACA6C0C97A88E14
c)에 따라
\[
\theta_\eta^4=\frac49i\cdot a_\nu^2-\frac16s_\nu^2+\frac{16}{15}\theta_\varrho^2-\frac{14}{45}(\nu\varrho x)^2+\frac2{15}(atu)^2.
\]

% inherited: noether_paper02_ko_translation_001_20260804\targets\Noether_P02_Korean_T13_U82_UNCHECKED.tex
% sha256: 4CCB89873B3AB6C12DD66ECBE009BAB1FD06DC0C938439E0A1C43A4D1A4F51AA
g)에 따라 $\theta_\eta^2H_\vartheta^2$의 두 값에서 다음을 얻는다.
\par\smallskip\noindent
\begin{minipage}[c]{3.6em}
\centering (15.)
\end{minipage}%
\begin{minipage}[c]{\dimexpr\linewidth-3.6em\relax}
\centering
\(\displaystyle
\begin{aligned}
s_\vartheta^2(\theta su)^2
&=\frac23s_\nu^2-2(atu)^2+2(\nu\varrho x)^2-\frac19J\cdot u_x^2\\
&=\frac89i\cdot a_\nu^2+\frac8{15}\theta_\varrho^2-\frac{14}{15}(atu)^2+\frac{38}{45}(\nu\varrho x)^2-\frac4{27}i^2\cdot u_x^2.
\end{aligned}
\)
\end{minipage}
\par\smallskip

% inherited: noether_paper02_ko_translation_001_20260804\targets\Noether_P02_Korean_T13_U83_UNCHECKED.tex
% sha256: 7D97C6BD7CDC9A306398C2DA32A47DE7BAA2690B7686AAE788180900D08E4629
g)에 따라 $\theta_\eta^3H_\vartheta$의 두 값에서 다음을 얻는다.
\par\smallskip\noindent
\begin{minipage}[c]{3.6em}
\centering (16.)
\end{minipage}%
\begin{minipage}[c]{\dimexpr\linewidth-3.6em\relax}
\centering
\(\displaystyle
s_\nu^2=\frac43i\cdot a_\nu^2+\frac45\theta_\varrho^2+\frac85(atu)^2-\frac{26}{15}(\nu\varrho x)^2+\frac16J\cdot u_x^2-\frac29i^2\cdot u_x^2.
\)
\end{minipage}
\par\smallskip

% inherited: noether_paper02_ko_translation_001_20260804\targets\Noether_P02_Korean_T13_U84_UNCHECKED.tex
% sha256: B363B2D92C39234CA8FE8D803D23A03A74384059997D28A38084A186726FBACD
공식 (16.)은 모듈 $(ss'u)^4$의 환원을 위한 기초를, 공식 (13.)은 $s$의 계를 환원하기 위한 기초를 제공한다(\S\ 17.).

\noindent\emph{결론:}

1) $a_\nu(j\nu x)^3$, $\theta_\eta H_\vartheta$, $\theta_\eta^2(\vartheta\eta x)$, $\theta_\eta^2(\theta Hu)^2$은 환원자이며, $a_\nu(j\nu x)^2$, $H_\vartheta$, $H_\vartheta^2$은 분해 가능 형식이다.

2) 계 II의 형식 $j(\nu)=g=(\nu\eta x)^4H_x^2$은 환원 가능하다.

3) 유일한 반변 형식 $g$가 환원 가능해지므로, $\nu$는 같은 계의 형식들과의 거듭제곱이나 곱 안에서 계 I과 전혀 수축되지 않는다.

$\theta$는 반공변식으로 볼 때만 곱이나 거듭제곱 안에서 수축되고, 이에 대응하여 $H$는 공변식으로 볼 때만 그러하다(\S\ 13 B 참조).

% inherited: noether_paper02_ko_translation_001_20260804\targets\Noether_P02_Korean_T14_U87_UNCHECKED.tex
% sha256: 3F1928778872FE6BC49D2E58F91FDE5A48AB9779D18FE8939357215FB462CA02
\begin{center}
\textbf{\S\ 12.}\quad\emph{계수차수 7의 형식들.}
\end{center}
\[
\text{다음의 수축: }\left\{\begin{array}{l}
\theta\cdot K\text{와 }\nu,\\
K,j,\A\text{와 }H,\\
f\text{와 }L,\sigma.
\end{array}\right.
\]

\noindent\emph{환원 불가능 형식들:}
\[
\renewcommand{\arraystretch}{1.35}
\begin{array}{@{}r@{\quad}c|c|c|c@{}}
&\cdot&\cdot&(KHu)^2&\cdot\\
&\cdot&K_\eta&\cdot&K_\eta(KHu)^2\\
\text{B)}&(k\eta x)^2&\cdot&K_\eta^2&\cdot\\
&(k\eta x)^3&H_k(k\eta x)^2,\ K_\eta(k\eta x)^2&\cdot&\cdot\\
&\cdot&K_\eta(k\eta x)^3&\cdot&\cdot
\end{array}\, ,
\]
\[
\renewcommand{\arraystretch}{1.35}
\begin{array}{@{}r@{\quad}c@{\qquad}c@{\qquad}c@{\qquad\qquad}r@{\quad}c@{}}
\text{C)}&(j\eta x)&H_j&\cdot&\text{D)}&(\A Hu)\\
&\cdot&H_j(j\eta x)&H_j^2,&&\A_\eta,
\end{array}
\]
\[
\renewcommand{\arraystretch}{1.35}
\begin{array}{@{}r@{\quad}c@{\qquad}c@{\qquad}c@{\qquad}c@{}}
&\cdot&(aLu)^2&(aLu)^3&\cdot\\
\text{E)}&a_l&\cdot&a_l(aLu)^2&a_l(aLu)^3\\
&\cdot&a_l^2&\cdot&\cdot
\end{array}
\]

% inherited: noether_paper02_ko_translation_001_20260804\targets\P02_T14_U88_r2.tex
% sha256: 3871B99D4C09A63399196AF9056C08464CF3E0B796D62CE8CE39626BB9742518
\noindent\emph{환원:}

\noindent A) $(\vartheta\cdot k,\nu,x)^4$은 분해 가능 형식 $(\vartheta^2\nu x)^4$에 대한 한 번의 전이연산으로서 분해된다.

\noindent B) 형식열 $H_kK_\eta$: 환원자 $H_\vartheta\theta_\eta$.

형식열 $K_\eta^2(KHu)$: 환원자 $a_\eta^2(aHu)$.

형식열 $K_\eta^2(k\eta x)$: 환원자 $\theta_\eta^2(\vartheta\eta x)$.

여기서 생기는 형식열 $K_\eta^3$의 이중 환원은 더 높은 형식열 $(j\eta x)^3$을 환원할 계기를 준다.

$(KHu)$, $(k\eta x)$, $K_\eta(k\eta x)$는 \S\ 3에 따른 분해 가능 형식이다.

$H_k(KHu)$와 $K_\eta(KHu)$는 분해 가능 형식 $K_\nu(k\nu x)$ 및 야코비 행렬식과 한 번 수축하여 생기므로 분해된다.
\[
K_\nu^2\equiv\theta_\nu^2(a\theta u)\equiv a_\nu^2(a\theta u)\modu{\nu^2}.
\]
$K_\nu^2$의 이중 표현에는 분해 가능 형식들 사이의 한 관계가 대응한다. 다음을 얻는다.

% inherited: noether_paper02_ko_translation_001_20260804\targets\Noether_P02_Korean_T14_U89_UNCHECKED.tex
% sha256: 0482F4C37470287AB1AE5F893E29AE2227D7B48146B703AE789388414878654F
\[
H_k(KHu)=2K_{\nu_1}(\nu\nu_1k)=2\theta_{\nu_1}(\nu\nu_1\widehat{a\theta})
=2\theta_\nu^2a_\nu-a_\nu^2\theta_\nu+f\theta_\eta(\theta Hu)^2-f\cdot \nu\cdot\theta_\nu^3\quad(\text{곱 정리})\modu{(\nu^3)},
\]
\[
\begin{aligned}
K_\eta(KHu)&=-K_\nu^2(\nu\nu_1x)=-a_\nu^2(a\theta u)(\nu\nu_1x)=-\theta_\nu^2(a\theta u)(\nu\nu_1x)\\
&=a_\eta(aHu)^2\cdot\theta+a_\nu^2\cdot\theta_\nu
=-\theta_\eta(\theta Hu)^2\cdot f-\theta_\nu^2\cdot a_\nu.
\end{aligned}
\]
\par\smallskip\noindent
\begin{minipage}[c]{4.8em}
\centering (17.)\srcfnmark{*)}
\end{minipage}%
\begin{minipage}[c]{\dimexpr\linewidth-4.8em\relax}
\centering
\(\displaystyle
0=a_\nu^2\cdot\theta_{\nu_1}+\theta_\nu^2\cdot a_{\nu_1}+f\cdot\theta_\eta(\theta Hu)^2+\theta\cdot a_\eta(aHu)^2\modu{(\nu^3)}.
\)
\end{minipage}
\par\smallskip
\srcfntext{*)}{이 경우의 $f\cdot\theta_\nu^3\cdot\nu$처럼 반공변식 하나를 포함하는 형식들의 곱은 분해 가능 형식들 사이의 모든 관계에서 생략한다. 계 $s_{IV}$에 공식을 적용할 때에는 이러한 곱들이 아무 의미도 갖지 않기 때문이다.}
형식 $H_k$, $H_k(k\eta x)$, $H_k^2$, $H_k^2(k\eta x)$는 분해 가능 형식 $H_\vartheta$와 $H_\vartheta^2$와 한 번 수축하여 생기므로 분해된다(\S\ 3. 2), a)와 b) 참조).

% inherited: noether_paper02_ko_translation_001_20260804\targets\Noether_P02_Korean_T14_U90_UNCHECKED.tex
% sha256: 80F531C69E9E6BE779D4EBC074CADAA01D71C6A34738D00BE5A15A1B9FC9735D
다음이 성립한다.
\[
H_k=H_\vartheta(a\theta u)\sim[H_\vartheta]_{\widehat{ua}}-f\cdot H_\vartheta(\theta Hu)\qquad(\text{2) a)에서 }y=\widehat{uv}\text{로 특수화}).
\]
\[
H_k(k\eta x)=H_\vartheta a_\eta-H_\vartheta\theta_\eta f,
\qquad
H_\vartheta a_\eta=\frac54[H_\vartheta]a-\frac14H_\vartheta a_\vartheta
\quad(\S\ 3. 2)b),
\]
\[
H_\vartheta^2(a\theta u)=[H_\vartheta^2]_{\widehat{ua}},
\qquad
H_\vartheta^2a_\eta=[H_\vartheta^2]a=H_k^2(k\eta x)+H_\vartheta^2\theta_\eta f.
\]

% inherited: noether_paper02_ko_translation_001_20260804\targets\Noether_P02_Korean_T14_U91_UNCHECKED.tex
% sha256: D023BC3E4E8DCA42D9868404873FD28EB8DA0EC9CC9DA07FA9729162E509B266
\noindent C) 형식열 $(j\eta x)^3$: B)에 따른 형식열 $K_\eta^3$의 이중 환원으로 환원 가능하다. 다음과 같이 둔다.
\[
\begin{aligned}
0&=a_\eta^3(a\theta u)=\theta_\eta^3(a\theta u)
=\{\theta_\eta+(\widehat{a\theta}\eta x)\}^3(a\theta u)-\theta_\eta^3(a\theta u)\\
&=\frac12(j\eta x)^3\modu{(\nu^3,s,\R,t,i,J)}.\qquad(\text{\S\ 5 끝부분 참조.})
\end{aligned}
\]
유사한 설정은 형식열의 최종형까지 성립한다.
\[
0=H_\vartheta^2(\widehat{a\theta}\eta x)a_\eta^3
=H_\vartheta^2(\widehat{a\theta}\eta x)\theta_\eta^3=\frac12H_j^2(j\eta x)^4\modu{(\nu^3,s,\R,t,i,J)}.
\]
형식 $(j\eta x)^2$: 1) 분해 가능 형식 $H_\vartheta^2$에 대한 관계 ((12.) b)를 두 번 편극하여 분해된다.

2) 곱 $a_\nu\cdot\theta_\nu^3$을 직접 계산하여 분해된다. 이로부터 더 높은 형식 $(\A Hu)^2$이 분해된다. 다음을 얻는다.

% inherited: noether_paper02_ko_translation_001_20260804\targets\Noether_P02_Korean_T14_U92_UNCHECKED.tex
% sha256: 35B73F9D6857A6614F3B0548A45F20D47093503C9E5D9F8B61D3601D4A065E17
\par\smallskip\noindent
\begin{minipage}[c]{4.8em}
\centering (18.)
\end{minipage}%
\begin{minipage}[c]{\dimexpr\linewidth-4.8em\relax}
\centering
\(\displaystyle
\left.\begin{array}{ll}
\text{a)}&\dfrac13(\A Hu)^2+\dfrac13(j\eta x)^2=\theta_\nu^2\cdot a_{\nu_1}^2,\\[0.6em]
\text{b)}&(j\eta x)^2=\dfrac43\theta_\nu^3\cdot a_{\nu_1},
\end{array}\right\}\modu{(\nu^3,s,\R,t,i,J)}
\)
\end{minipage}
\par\smallskip
다음과 같이 두었다.
\[
\begin{aligned}
H_\vartheta^2(a\theta u)^2-\frac13(\A Hu)^2
&=2\theta_\eta^2(a\theta u)(aHu)+\theta_\nu^2\cdot a_{\nu_1}^2\quad(\text{곱 정리})\\
&=(a\theta u)(aHu)\{a_\eta-(\widehat{a\theta}\eta x)\}^2+\theta_\nu^2\cdot a_{\nu_1}^2,
\end{aligned}
\]
\[
\begin{aligned}
a_{\nu_1}\cdot\theta_\nu^3
&=-(a\widehat{\nu\nu_1}u)\theta_\nu^2(\theta\widehat{\nu\nu_1}u)-\frac12(a\widehat{\nu\nu_1}u)\theta_\nu\theta_{\nu_1}(\theta\widehat{\nu\nu_1}u)\\
&=-\frac32\theta_\eta^2(\theta Hu)(aHu)=-\frac32\theta_\eta^2(\theta Hu)(a\theta u)\\
&=-\frac32\{a_\eta-(\widehat{a\theta}\eta x)\}^2\{(aHu)-(a\theta u)\}(a\theta u).
\end{aligned}
\]

% inherited: noether_paper02_ko_translation_001_20260804\targets\P02_T14_U93_r2.tex
% sha256: D188A895261758C00DC4C694217E5C78DE23085A0D7D66F429F77FB7E6C229CD
\noindent D) 형식열 $\A_\eta(\A Hu)$: 환원자 $\A_\nu^2$.

$(\A Hu)^2$은 C)에 따른 분해 가능 형식이다.

\noindent E) 형식열 $a_l^2(aLu)$: 환원자 $a_\eta^2(aHu)$.
형식 $(aLu)$와 $a_l(aLu)$: \S\ 3에 따라 야코비 행렬식에 대한 전이연산으로서 분해된다.

\noindent F) 형식열 $a_\sigma^3$: 환원자 $a_\eta^3$.

형식 $a_\sigma^2$와 $a_\sigma^3$: 대응하는 식
\[
H_\nu a_\nu^3\quad\text{및}\quad H_\nu a_\eta^3,
\qquad
H_\nu a_\nu^3a_\eta\quad\text{및}\quad H_\nu a_\eta^4
\]
을 이중 환원하여 얻는 환원자 $a_\nu^3$과 $a_\eta^3$ (\S\ 10. C)의 $\A_\nu^2$와 $\A_\eta^3$의 환원 참조).

% inherited: noether_paper02_ko_translation_001_20260804\targets\Noether_P02_Korean_T14_U94_UNCHECKED.tex
% sha256: E7CE57F0D663C5716308F1DC217B2044D5F2F470753D3C38C0860DDB57D6541E
이로부터 (16.)을 고려하여 더 높은 형식 $a_\nu s_\nu(asu)^2$, $a_\nu^2s_\nu(asu)^2$ 및 기호 $\R,t(a_\nu,a_\eta^2)$로 된 형식들에 대한 환원을 얻는다.
\par\smallskip\noindent
\begin{minipage}[c]{4.8em}
\centering (19.)
\end{minipage}%
\begin{minipage}[c]{\dimexpr\linewidth-4.8em\relax}
\centering
\(\displaystyle
\left.\begin{array}{rcl}
a_\nu s_\nu(asu)^2&=&\dfrac13 i\theta_\nu^2-\dfrac1{18}iH,\\[0.5em]
a_\nu^2s_\nu(asu)^2&=&0,
\end{array}\right\}\modu{(\R,t)}
\)
\end{minipage}
\par\smallskip
형식 $a_\sigma$: (12.)b를 편극하여 분해되며, 더 짧게는 다음과 같이 두고 곱 $a_\nu^2\cdot\theta_\nu^3$을 직접 전개하여 분해된다.

% inherited: noether_paper02_ko_translation_001_20260804\targets\Noether_P02_Korean_T14_U95_UNCHECKED.tex
% sha256: D6A85D59DA3A52B5CAE50A1D47AE6904818F397F91222D9927F65F5A102CE629
\[
a_\nu^2\cdot\theta_{\nu_1}^3=a_\nu^2\theta_{\nu_1}\{a_{\nu_1}-(\widehat{a\theta}\nu_1x)\}^2
=-2a_\eta(aHu)(a\theta H)(\widehat{a\theta}\eta x)+(\widehat{a\theta}\nu_1x)^2a_\nu^2\theta_{\nu_1}
\]
그리고 $H_j(j\eta x)^2$의 환원(C)을 고려하면 다음을 얻는다.
\par\smallskip\noindent
\begin{minipage}[c]{4.8em}
\centering (20.)
\end{minipage}%
\begin{minipage}[c]{\dimexpr\linewidth-4.8em\relax}
\centering
\(\displaystyle a_\sigma=2a_\nu^2\cdot\theta_{\nu_1}^3\modu{(s,\R,t,i)}.\)
\end{minipage}
\par\smallskip

% inherited: noether_paper02_ko_translation_001_20260804\targets\P02_T14_U96_r2.tex
% sha256: 9193B90F233F0459C89E978D71A283238CA33DE4F7C21703371D64733423B2A3
\noindent\emph{결론:}

1) $(j\eta x)^3$, $\A_\eta(\A Hu)$, $a_\sigma^2$은 환원자이고, $(j\eta x)^2$, $(\A Hu)^2$, $a_\sigma$는 분해 가능 형식이다.

2) $f$는 반변 형식들과의 곱 안에서만 $L$과 수축되며, $f$와 따라서 $K$ 및 $N$은 $\sigma$와 전혀 수축되지 않고, 그러므로 $(\nu\sigma x)$와도 수축되지 않는다. $j$는 반변 형식들과의 곱 안에서만 $H$와 $L$에 수축되며, $\A$는 곱 안에서 $H$나 $L$과 전혀 수축되지 않는다($\A_\eta(\A Hu)$와 $(\A Hu)^2$으로부터 따른다). 마찬가지로 $\sigma$와 따라서 $(\nu\sigma x)$는 계 I의 어떤 형식과도 곱 안에서는 수축되지 않는다. 계 II의 곱들 가운데에는 \S\ 11의 결론들을 고려하면 $H$의 거듭제곱과 $H$와 $L$의 곱만 남는다.

% inherited: noether_paper02_ko_translation_001_20260804\targets\Noether_P02_Korean_T15_U97_UNCHECKED.tex
% sha256: 317EBE67DFA9E121D46E46B8F2D4346468BBEC563D7FF34D32D9332765026D59
\begin{center}
\textbf{\S\ 13.}\quad\emph{계수차수 8의 형식들(계 III).}
\end{center}
\[
\text{다음의 수축: }\left\{\begin{array}{l}
\A\cdot j\text{와 }\nu,\\
\theta^2,
 f\cdot j,
 N\text{과 }H,\\
\theta\text{와 }L,\sigma.
\end{array}\right.
\]


\noindent\emph{환원 불가능 형식들:}
\[
\renewcommand{\arraystretch}{1.28}
\begin{array}{@{}r@{\quad}c@{\qquad\qquad}r@{\quad}c@{\qquad}c@{}}
\text{A)}&\A_\nu(j\nu x),&\text{B)}&(\vartheta^2\eta x)^3&\cdot\\
&&&(\vartheta^2\eta x)^4&H_\vartheta(\vartheta^2\eta x)^3,\ \theta_\eta(\vartheta^2\eta x)^3,
\end{array}
\]
\[
\renewcommand{\arraystretch}{1.28}
\begin{array}{@{}r@{\quad}c@{\qquad}c@{\qquad\qquad}r@{\quad}c@{}}
\text{C)}&(aHu)H_j&\cdot&\text{D)}&H_\eta;\ N_\eta,\\
&\cdot&a_\eta(aHu)H_j&&
\end{array}
\]
\[
\renewcommand{\arraystretch}{1.30}
\begin{array}{@{}r@{\quad}c@{\qquad}c@{\qquad}c@{\qquad}c@{\qquad}c@{}}
&\cdot&\cdot&(\theta Lu)^2&(\theta Lu)^3&\cdot\\
&\cdot&\theta_l&\cdot&L_\vartheta(\theta Lu)^2,\ \theta_l(\theta Lu)^2&\theta_l(\theta Lu)^3\\
\text{E)}&(\vartheta lx)^2&\cdot&\theta_l^2&\cdot&\cdot\\
&\cdot&\theta_l(\vartheta lx)^2&\cdot&\cdot&\cdot
\end{array}
\]

% inherited: noether_paper02_ko_translation_001_20260804\targets\Noether_P02_Korean_T15_U98_UNCHECKED.tex
% sha256: 6614505C55919AE6B9B8536D5324EFEE19F9F68CCB02D549A594268A965E381D
\noindent\emph{환원:}

\noindent A) $\A_\nu(j\nu x)^3$: 환원자 $a_\nu(j\nu x)^3$.

$\A_\nu(j\nu x)^2$은 환원자 $\theta_{\vartheta j}$를 고려할 때 \S\ 3. 2)b에 따라 분해 가능 형식 $a_\nu(j\nu x)^2$에 의해 환원 가능하다. 실제로 \S\ 11 B)에 따라 다음이 성립한다.
\[
\frac3{10}\A_\nu(j\nu x)^2+\frac35a_\nu(j\nu\vartheta)a_\vartheta(j\nu x)
+\frac1{10}a_\nu(j\nu\vartheta)^2
=\frac35\theta_\nu^3\theta_{\vartheta j}'+\frac25\theta_\nu^3\theta_{\vartheta j}\theta_{\vartheta j}'',
\]
(환원자 $a_j$와 $\theta_{\vartheta j}$).

% inherited: noether_paper02_ko_translation_001_20260804\targets\Noether_P02_Korean_T15_U99_UNCHECKED.tex
% sha256: F4116F4A4E030DEE6737215DFD82C719CE4C5C028D02C438CCEBED1D22748E6A
\noindent B) 형식 $H_\vartheta(\vartheta'\eta x)^2$, $H_\vartheta^2(\vartheta'\eta x)$, $H_\vartheta^2(\vartheta'\eta x)^2$은 \S\ 3. 2)b)와 다음 관계에 따라 분해 가능 형식 $H_\vartheta$와 $H_\vartheta^2$, 각각 $H_\vartheta a_\eta$와 $H_\vartheta^2a_\eta$와의 연속적인 한 번의 수축으로 생기므로 분해된다.
\[
H_\vartheta(\vartheta'\eta x)^2=2H_\vartheta a_\eta^2f-2H_\vartheta a_\eta b_\eta.
\]

% inherited: noether_paper02_ko_translation_001_20260804\targets\Noether_P02_Korean_T15_U100_UNCHECKED.tex
% sha256: B4620CBC3499B54F7C03AECAD4E39F2484359FFAC95E452A603C920D499729A7
\noindent C) 형식열 $a_\eta(j\eta x)^2$: 환원자 $a_j$와 $(j\eta x)^3$.

형식 $a_\eta(j\eta x)$은 분해된다. 환원자 $a_j$와 \S\ 3. 2)a에 따른 분해 가능 형식 $(j\eta x)^2$과의 한 번의 수축을 이용한다(특수화 $y=uv$).

형식 $a_\eta H_j$는 다음 두 방식으로 분해된다. 1) 분해 가능 형식 $a_\nu(j\nu x)^2$과 한 번 수축한다.

2) 분해 가능 형식 $(j\eta x)^2$과 특수하게 두 번 수축한다. 이로부터 분해 가능 형식들 사이의 관계가 생긴다.

$a_\nu(j\nu x)^2$으로부터:
\[
0=\theta_\nu'\theta_\nu+2a_\eta H_j\modu{(s,\R,t,i,J)}.
\]
$(j\eta x)^2$으로부터:
\[
0=\theta_\nu'\theta_\nu-\frac32a_\eta H_j\modu{(s,\R,t,i,J)}
\]
(반공변식과의 곱은 생략한다). 다음과 같이 둔다.
\[
0=\frac12a_\nu(abu)^2\theta_\nu^3
+\frac12a_\nu(abu)\theta_{\nu_1}^3(\theta bu)
-\frac34(\widehat{j\eta}bu)^2+\frac34H_j(\widehat{j\eta}bu)
-\frac18fH_j^2
\]
그리고 $a_\nu(abu)(\theta bu)\theta_{\nu_1}^3$에서 기호들을 교환한다.

% inherited: noether_paper02_ko_translation_001_20260804\targets\Noether_P02_Korean_T15_U101_UNCHECKED.tex
% sha256: 115B9F19017AB3481E6803C71B2971FFEDBAEB33147A59116AC78D41AFE8FF61
\noindent D) 형식열 $N_\eta(NHu)$: 환원자 $\A_\nu^2$.

$(NHu)$, $(n\eta x)$, $N_\eta(n\eta x)$, $H_n(NHu)$는 분해 가능 형식이며(\S\ 12 B 참조), $(NHu)^2$은 형식 $(\A Hu)^2$과 한 번 수축함으로써 분해된다.

% inherited: noether_paper02_ko_translation_001_20260804\targets\Noether_P02_Korean_T15_U102_UNCHECKED.tex
% sha256: A683EB0F3FE83D95CB5C3F9A13D2C1E756F781AFD6AEACAF2CF9BB38D30E170A
\noindent E) 형식열 $\theta_lL_\vartheta$, $\theta_l^2(\theta Lu)^2$, $\theta_l^2(\vartheta lx)$:

환원자: $\theta_\eta H_\vartheta$, $\theta_\eta^2(\theta Hu)^2$, $\theta_\eta^2(\vartheta\eta x)$.

형식 $(\theta Lu)$, $(\vartheta lx)$, $L_\vartheta$, $L_\vartheta(\theta Lu)$, $\theta_l(\theta Lu)$, $L_\vartheta(\vartheta lx)$, $\theta_l(\vartheta lx)$, $L_\vartheta^2$, $L_\vartheta^2(\theta Lu)$, $\theta_l^2(\theta Lu)$는 분해된다. (\S\ 12 B의 분해 가능 형식들과 쌍대 관계에 있다.)

% inherited: noether_paper02_ko_translation_001_20260804\targets\Noether_P02_Korean_T15_U103_UNCHECKED.tex
% sha256: 5AFF19C22B815166C202A7A7BB512713EC155A3A5B0B0DDBC55572BEE6F2B3D5
\noindent F) 형식열 $\theta_\sigma(\vartheta\sigma x)$: 환원자 $a_\sigma^2$.

형식 $(\vartheta\sigma x)$와 $(\vartheta\sigma x)^2$은 분해 가능 형식 $a_\sigma$와 한 번 수축하여 생기므로 분해된다. 두 번 수축하여 생긴 $\theta_\sigma$도 다음 형식열 때문에 분해된다.
\[
a_\nu^2(a\theta u)\theta_{\nu_1}^3\equiv H_j^2(j\eta x)\equiv0\;\text{법 }(s,\R,t,i,J):
\]
\[
\theta_\sigma=a_\sigma(abu)^2=a_\nu^2(abu)^2\cdot\theta_{\nu_1}^3.
\]

% inherited: noether_paper02_ko_translation_001_20260804\targets\Noether_P02_Korean_T15_U104_UNCHECKED.tex
% sha256: E6BA78C78DDEA17069CBF0D5F13A03E5219ACFF70CF73C6FE2F4CF70BF3BFF39
\noindent\emph{결론:}

$\theta^3$ 또는 $\theta^2K$는 $H$와 수축되지 않는다. $\theta$는 반공변식으로 볼 때만 거듭제곱이나 곱 안에서 $L$과 수축되며, $\sigma$와 $(\nu\sigma x)$과는 전혀 수축되지 않는다.

$N$은 곱 안에서 계 II의 어떤 형식과도 수축되지 않으며, 계 II 형식들의 거듭제곱이나 곱과도 마찬가지다. 앞 절들의 결론을 고려하면 계 I에 남는 곱은 $f\cdot j$, $\A\cdot j$, $\theta$의 거듭제곱, 그리고 $\theta\cdot K$의 곱뿐이다.

% inherited: noether_paper02_ko_translation_001_20260804\targets\Noether_P02_Korean_T16_U105_UNCHECKED.tex
% sha256: 39054EDB454DC69AE6B6FD249DF192B374AE2D274CB7BB425E68975085B0B23C
\begin{center}
\textbf{\S\ 14.}\quad\emph{계수차수 9의 형식들(계 III).}
\end{center}
\[
\text{다음의 수축: }\left\{\begin{array}{l}
\theta\cdot K\text{와 }H,\\
K,j,\A\text{와 }L,\\
j,\A\text{와 }\sigma,\\
f\text{와 }H^2.
\end{array}\right.
\]


\noindent\emph{환원 불가능 형식들:}
\[
\renewcommand{\arraystretch}{1.35}
\begin{array}{@{}r@{\quad}c@{\qquad}c@{\qquad\qquad}r@{\quad}c@{\qquad}c@{\qquad}c@{\qquad}c@{}}
&\cdot&\cdot&&\cdot&\cdot&(KLu)^3&\cdot\\
\text{A)}&(\vartheta\cdot k,\eta,x)^4&\cdot&\text{B)}&\cdot&\cdot&\cdot&K_l(KLu)^3\\
&\cdot&H_k(\vartheta\cdot k,\eta,x)^4&&\cdot&K_l^2&\cdot&\cdot\\
&&&& (klx)^3&\cdot&\cdot&\cdot\\
&&&&\cdot&K_l(Klx)^3&\cdot&\cdot
\end{array}
\]
\[
\renewcommand{\arraystretch}{1.35}
\begin{array}{@{}r@{\quad}c@{\qquad}c@{\qquad}r@{\quad}c@{\qquad}r@{\quad}c@{\qquad}r@{\quad}c@{\qquad}c@{}}
\text{C)}&L_j&\cdot&\text{D)}&\A_l&\text{E)}&(j\sigma x),&\text{G)}&(aH^2u)^3&(aH^2u)^4\\
&\cdot&L_j^2&&&&&&\cdot&a_\eta(aH^2u)^3
\end{array}
\]

% inherited: noether_paper02_ko_translation_001_20260804\targets\P02_T16_U106_r2.tex
% sha256: 8B3C780E50DB8763FBA238504295CE4EE997BD49FDC38AC5D479A0B2743F8143
\noindent\emph{환원:}

A) 형식 $H_\vartheta(k\eta x)^3$, $H_\vartheta^2(k\eta x)^2$, $H_\vartheta^2(k\eta x)^3$은 분해 가능 형식들과 연속적으로 한 번 수축하여 생기므로 분해된다.

B) 형식열 $K_lL_k$, $K_l^2(KLu)$, $K_l^2(klx)$:

환원자: $\theta_\eta H_\vartheta$, $a_\eta^2(aHu)$, $\theta_\eta^2(\vartheta\eta x)$.

형식 $K_l$, $K_l(KLu)$, $K_l(klx)$, $(KLu)^2$, $(klx)^2$은 야코비 행렬식에 대한 전이연산으로서 분해된다(\S\ 3).

형식 $L_k$, $L_k(KLu)$, $L_k(KLu)^2$, $L_k(klx)$, $L_k(klx)^2$, $L_k^2$, $L_k^2(KLu)$, $L_k^2(klx)$, $L_k^3$, $K_l(KLu)^2$, $K_l(klx)^2$은 다음 분해 가능 형식들과 한 번 수축하여 생기므로 분해된다.
\[
L_\vartheta,
\ H_k(KHu),
\ L_\vartheta(\vartheta lx),
\ H_k^2,
\ L_\vartheta^2,
\ L_\vartheta^2(\theta Lu),
\ K_\eta(KHu),
\ \theta_l(\vartheta lx)
\]
이는 \S\ 3. 2), a)와 b)에 따른다.

% inherited: noether_paper02_ko_translation_001_20260804\targets\Noether_P02_Korean_T16_U107_UNCHECKED.tex
% sha256: F0CCD752F8B532B1955CC2A1973F2303EA85C143DC01F311EF83FAB74F5A3567
C) 형식열 $(jlx)^3$: 환원자 $(j\eta x)^3$.

형식 $(jlx)^2$과 $L_j(jlx)$은 분해 가능 형식 $(j\eta x)^2$과 한 번 수축하여 생긴다.

% inherited: noether_paper02_ko_translation_001_20260804\targets\Noether_P02_Korean_T16_U108_UNCHECKED.tex
% sha256: E5C5ED59FE7138118CD5ADFBE16B862BC9E6E68E2A488A6CFC55204AAFCAC2C2
D) 형식열 $\A_l(\A Lu)$: 환원자 $\A_\nu^2$.

형식 $(\A Lu)^2$, $(\A Lu)^3$: 분해 가능 형식 $(\A Hu)^2$과 한 번 수축한다.

% inherited: noether_paper02_ko_translation_001_20260804\targets\Noether_P02_Korean_T16_U109_UNCHECKED.tex
% sha256: 22A49EF62EDC80863BB4B5BEFF56D1B3FAFD101DAC9A6293D865DF26AC3A8939
E) 형식열 $(j\sigma x)^2$: $j$를 형식열의 더 낮은 형식들로 바꾸어 얻는 환원자 $a_\sigma^2$(\S\ 5):
\[
a_\sigma^2(a\theta u)^2=\frac13(j\sigma x)^2,
\qquad
 a_\sigma^2(a\theta\sigma x)^2=\frac13(j\sigma x)^4.
\]

% inherited: noether_paper02_ko_translation_001_20260804\targets\Noether_P02_Korean_T16_U110_UNCHECKED.tex
% sha256: 5F97B78C1330732F207F8A363122C334FBED3A6F4EFCB8CC96D00419DC9615D0
F) 형식열 $\A_\sigma^2$: 환원자 $a_\sigma^2$.

형식 $\A_\sigma$: $\A$를 형식열의 더 낮은 형식들로 바꾸어 얻는 환원자 $a_\sigma^2$:
\[
a_\vartheta a_\sigma^2=\frac23\A_\sigma\modu{\nu}.
\]

% inherited: noether_paper02_ko_translation_001_20260804\targets\Noether_P02_Korean_T16_U111_UNCHECKED.tex
% sha256: 99E739D6ECA5F3C2507D5104B13D61935FB952FD9D9F55041DD70B04874966D1
\noindent\emph{결론:} 형식 $j$만 $\sigma$와 $(\nu\sigma x)$과 수축된다.

$\A_l$에 대응하는 형식 $N_l$이 분해되므로 $N$은 $L$과 수축되지 않는다.

% inherited: noether_paper02_ko_translation_001_20260804\targets\Noether_P02_Korean_T17_U115_UNCHECKED.tex
% sha256: 8F57672E9BD1B58075A096A6AD0C8ADB4FEF24D7BE72B98CB9C3BA8A4794AAE2
\begin{center}
\textbf{\S\ 15.}\quad\emph{계수차수 10의 형식들(계 III).}
\end{center}
\[
\text{다음의 수축: }\left\{\begin{array}{l}
\theta^2,
\ f\cdot j\text{와 }L,\\
\theta\text{와 }H^2.
\end{array}\right.
\]
\noindent\emph{환원 불가능 형식들:}
\[
\begin{array}{c@{\quad}l@{\qquad}c@{\quad}l}
\text{A)}&(\vartheta^2lx)^4,
\ \theta_l(\vartheta^2lx)^4&
\text{B)}&(aLu)^2L_j,
\ a_l(aLu)^2L_j,\\[0.8em]
\text{C)}&(\theta H^2u)^3,
\ (\theta H^2u)^4,
\ H_\vartheta(\theta H^2u)^3,
\ \theta_\eta(\theta H^2u)^3.
\end{array}
\]

\noindent\emph{환원:}

A)와 B): 나머지 형식들은 분해 가능 형식과 한 번 수축하여 분해된다.

C) 나머지 형식들은 \S\ 13 B)에 따라 쌍대적으로 반대되는 고찰을 통해 분해된다.

\noindent\emph{결론:}

$\theta^2\cdot K$는 $L$과 수축되지 않으며, 이에 대응하여 $H^2L$은 $K$와 수축되지 않는다. $H^3$과 $H^2L$은 $\theta$와 수축되지 않는다. 이로써 \S\ 7의 수축 도식이 증명된다.

% inherited: noether_paper02_ko_translation_001_20260804\targets\Noether_P02_Korean_T18_U116_UNCHECKED.tex
% sha256: 9DE99461050204AC7E9A678959B4FE2DF61C8309E8B0C084693AF6D70B1FB709
\begin{center}
\textbf{\S\ 16.}\quad\emph{계수차수 11, 12, 13, 15의 형식들(계 III).}
\end{center}

\noindent\emph{계수차수 11.}
\[
\text{다음의 수축: }\left\{\begin{array}{l}
\theta\cdot K\text{와 }L,\\
K,j\text{와 }H^2,\\
j\text{와 }(\nu\sigma x),\\
f\text{와 }H\cdot L.
\end{array}\right.
\]
\noindent\emph{환원 불가능 형식들:}
\[
\begin{array}{c@{\quad}l@{\qquad}c@{\quad}l}
\text{A)}&(\vartheta\cdot k,l,x)^5,
\ \theta_l(\vartheta\cdot k,l,x)^5&
\text{B)}&(KH^2u)^4,
\ K_\eta(KH^2u)^4,\\[0.8em]
\text{C)}&(\nu\sigma j),&
\text{D)}&(a_\eta H\cdot L,u)^4.
\end{array}
\]

% inherited: noether_paper02_ko_translation_001_20260804\targets\Noether_P02_Korean_T18_U117_UNCHECKED.tex
% sha256: 622320A768458D226C829CB5027D99217C8FBF539D5199956DAB51A0DC001392
\noindent\emph{계수차수 12.}
\[
\text{다음의 수축: }\left\{\begin{array}{l}
\theta^3\text{와 }L,\\
\theta\text{와 }H\cdot L.
\end{array}\right.
\]
\noindent\emph{환원 불가능 형식들:}
\[
\text{E)}\ (\vartheta^3lx)^6,
\qquad
\text{F)}\ (\theta,H\cdot L,u)^4,
\qquad L_\vartheta(\theta,H\cdot L,u^4).
\]

\noindent\emph{계수차수 13.}
\[
\text{다음의 수축: }K\text{와 }H\cdot L.
\]
\noindent\emph{환원 불가능 형식들:}
\[
\text{G)}\ (K,H\cdot L,u)^5,
\qquad K_\eta(K,H\cdot L,u)^5.
\]

\noindent\emph{계수차수 15.}
\[
\text{다음의 수축: }K\text{와 }H_3.
\]
\noindent\emph{환원 불가능 형식:}
\[
\text{H)}\ (KH^3u)^6.
\]

% inherited: noether_paper02_ko_translation_001_20260804\targets\Noether_P02_Korean_T18_U118_UNCHECKED.tex
% sha256: 25F528A7F973CB6F9DFA17C7F0B5122EAD76BBAD2BA2F3A703DEDCB6FA1D6E17
\noindent\emph{환원:}

형식들 $H_j^3$, $H_j^4$. 다음 관계식에서 얻는 환원자는 $L_j^3$이다:
\[
(\nu lx)L_x^3=-\frac12H^2\modu{(\sigma,s)}.
\]
나머지 형식들은 앞 절들의 고찰에 따라 분해되거나 환원 가능하다.

\noindent\emph{결론:}

\S\ 7의 수축 도식과 \S\S\ 8부터 15까지의 결론들에 따라, 이로써 계 III의 환원 불가능 형식들(117개 형식)이 모두 망라되었다.

\clearpage

% inherited: noether_paper02_ko_translation_001_20260804\targets\Noether_P02_Korean_T19_U119_UNCHECKED.tex
% sha256: 46B1E4497EEF12DC7D0BB536FAECB2AEBD30F4B8DBB29B1FAC81D493FDE021AF
\begin{center}
{\Large\bfseries 제 IV장.\quad $(\varrho,t)$를 법으로 하는 상대 완전계.}
\end{center}

\begin{center}
\textbf{\S\ 17.}\quad\emph{모듈 $(ss'u)^4$와 $s$의 계의 환원.}
\end{center}

% inherited: noether_paper02_ko_translation_001_20260804\targets\Noether_P02_Korean_T19_U120_UNCHECKED.tex
% sha256: 2FD3414F8C57C08F12A66F011580F828334D14E2BFB4840DCCFDE7FDD6FF5E72
바로 다음의 더 높은 상대 완전계를 구성하려면, \S\ 7과 유사하게 다음 모듈에 관하여 $s$의 계를 구성하고
\[
\nu(s)=(ss'u)^4
\]
$s$를 법으로 하는 상대 완전계의 형식들과 수축해야 한다. 더 높은 모듈로서 모듈 $(\varrho,t)$를 도입하면 다음이 성립함을 보이고자 한다.
\[
\begin{array}{ll}
\text{A)}&\text{모듈 }(ss'u)^4\text{가 환원 가능해지고,}\vphantom{\dfrac11}\\
\text{B)}&s\text{의 계는 유일한 형식 }s\text{로 이루어진다.}
\end{array}
\]

% inherited: noether_paper02_ko_translation_001_20260804\targets\Noether_P02_Korean_T19_U121_UNCHECKED.tex
% sha256: 592F2353A6E7D81EAD4C94E6D93541023C4E2D45216421919599D6FCEA7881F2
A) 모듈 $(ss'u)^4$의 환원은 \S\ 11의 공식 (16.)으로 구한 환원자 $s_\nu^2$에서 즉시 얻어진다. 다음 식에서
\[
 s_\nu^2=\frac43 i a_\nu^2+\frac45\theta_\varrho^2+\frac85(atu)^2-
 \frac{26}{15}(\nu\varrho x)^2+\frac16J\cdot u_x^2-
 \frac29i^2\cdot u_x^2
\]
$x$를 $\nu_1$에 따라 편극하면 다음을 얻는다.
\srcnumdisplay{(21.)\srcfnmark{*)}}{%
\begin{aligned}
(ss'u)^4&=-\frac23J\cdot \nu+6s_\nu^2s_{\nu_1}^2\\
&=\frac{24}{5}\theta_\nu^2\theta_\varrho^2+
\frac{48}{5}a_\nu^2(atu)^2-
\frac{52}{5}H_\varrho^2+
\frac43i\cdot(asu)^4+\frac13J\cdot \nu-
\frac49i^2\cdot \nu,
\end{aligned}
}
이로써 모듈 $(ss'u)^4$가 환원된다.

% inherited: noether_paper02_ko_translation_001_20260804\targets\Noether_P02_Korean_T19_U122_UNCHECKED.tex
% sha256: 446BAC481471D75C6132DCF4AC7D38EE67F8F98F8629EB52FDD89BF475C546D7
\srcfntext{*)}{공식 (21.)에서, 공식 (10.)c를 적용하면 서론의 주석에서 언급한 계수차수 9의 세 불변식 사이의 관계가 나온다:
\srcnumdisplay{(22.)}{%
\begin{aligned}
(ass')^4&=\frac{24}{5}\theta_\nu^2\theta_\varrho^2a_\nu^2a_\vartheta^2+
\frac{48}{5}a_\nu^2b_\nu^2(tab)^2-
\frac{52}{5}H_\varrho^2a_\eta^4+
\frac53J\cdot i-\frac49i^3,\\
(ass')^4&=\frac{24}{5}a_\varrho^2a_{\varrho'}^2-
\frac{12}{5}t_\varrho^2+\frac53J\cdot i-\frac49i^3.
\end{aligned}
}
}

% inherited: noether_paper02_ko_translation_001_20260804\targets\Noether_P02_Korean_T19_U123_UNCHECKED.tex
% sha256: 6A9219B4E9ADCD1EA188AE6DB0437275085555F2DF491E459D72209F2CFCF75A
B) 다음 형식의 환원과
\[
Z=(ss'u)^2=\theta(s)
\]
그에 따른 $s$의 계의 환원은, \S\ 7의 공식 (8.)b에 따라 형식 $Z$를 더 낮은 형식 $g_\nu^2$으로 바꾸고 다음 관계를 고려함으로써 얻어진다.
\[
 s_\eta^2=s_\nu^2(\nu\nu_1x)^2-\frac13Z.
\]
형식 $g_\nu^2$은 \S\ 11의 공식 (13.)에 따라 환원자 $s_\nu^2$를 인수로 갖는다. 공식 (8.)과 (16.)에 따라 다음이 성립한다.
\[
 g_\nu^2=-Z+\frac{14}{5}ia_\eta^2+\frac{14}{15}i(asu)^2
 \modu{(\varrho,t)},
\]
그리고 공식 (13.), (14.), (15.), (16.)에 따라 다음이 성립한다.
\[
\begin{aligned}
g_\nu^2&=2[s_\vartheta^2]_{\nu^2}-\frac43 i\A_\nu^2
      =2s_\vartheta^2s_\nu^2-\frac65[s_\vartheta^2(\theta su)^2]\widehat{\nu}x^2
        -\frac43i\A_\nu^2\\
&=\frac45i\cdot a_\eta^2-\frac25i\cdot(asu)^2+\frac13J\cdot\theta
\modu{(\varrho,t)}.
\end{aligned}
\]
따라서 다음을 얻는다.
\srcnumdisplay{(23.)}{%
Z=2i\cdot a_\eta^2+\frac43i\cdot(asu)^2-\frac13J\cdot\theta
 \modu{(\varrho,t)}.
}

% inherited: noether_paper02_ko_translation_001_20260804\targets\Noether_P02_Korean_T19_U124_UNCHECKED.tex
% sha256: 34F4BAC1584957B84E4E5BAE103F78D9BF2B56551C530E115B1B9DA3279B27F8
따라서 $((ss'u)^4,\varrho,t)$를 법으로 하는 상대 완전계는 $(\varrho,t)$를 법으로 하는 상대 완전계로 넘어가며, 이 계를 알려진 두 이차형식의 계와 전이연산하면 절대 완전계가 생긴다. $(\varrho,t)$를 법으로 하는 상대 완전계를 구성하기 위해서는, 공식 (23.)에 따라 $s$의 계가 형식 $s$로 환원되므로, $(s,\varrho,t)$를 법으로 하는 상대 완전계를 $s$ 및 $s$의 거듭제곱들과 전이연산해야 한다. $s$의 거듭제곱들은 계 I과만 수축하게 됨을 보일 것이다.

% inherited: noether_paper02_ko_translation_001_20260804\targets\Noether_P02_Korean_T19_U125_UNCHECKED.tex
% sha256: BEA5609EBE22A2EF99A6569BE8CA2348D3EE949682999E2209AE5B9B8229CFAD
같은 계수차수와 같은 차수를 지닌 형식들 가운데, $(s,\varrho,t)$를 법으로 하는 상대 완전계의 더 높은 형식들을 $s$와 수축하여 생긴 형식들을 \glqq{}더 높은 형식들\grqq{}로 정의한다. 이때 $s$와의 수축은 원래 형식들의 편극으로만 간주한다. 이로써 개별 형식들의 배열은 유일하게 정해진다(\S\ 1, 형식열 2 참조). 예를 들면 다음과 같다.
\[
(\theta_\eta(\theta Hu),s,u)^2\text{는 }s_\eta((\theta Hu)^2,s,u)\text{보다 더 높고}\srcfnmark{*)},
\]
\[
(\theta_\eta(\theta Hu)^2,s,u)\text{는 }(\theta_\eta(\theta Hu),s,u)^2\text{보다 더 높다.}
\]

% inherited: noether_paper02_ko_translation_001_20260804\targets\Noether_P02_Korean_T19_U126_UNCHECKED.tex
% sha256: 2548E3C756F726A999C3D59DB1C92A1CC0B6FF1DD9DC1E09DD7288ACB2A50129
생기는 계를 다음 네 부분계로 나눈다.
\[
\begin{array}{ll}
\text{계 }s_I:&s\text{ 및 }s\text{의 거듭제곱들을 계 I과 수축하여 생긴다,}\\
\text{계 }s_{II}:&s\text{를 계 II와 수축하여 생긴다,}\\
\text{계 }s_{III}:&s\text{를 계 III의 개별 형식들(곱은 제외)과 수축하거나,}\\
&\text{계 I과 II에서 각각 한 형식씩 취한 곱과 수축하여 생긴다,}\\
\text{계 }s_{IV}:&s\text{를 계 I과 III, II와 III, III과 III에서 각각 한 형식씩}\\
&\text{취한 곱과 수축하여 생긴다.}
\end{array}
\]
또한 이제부터 모든 공식은 $(\varrho,t)$를 법으로 이해하고, 공식 (27.)부터는 $(\varrho,t,i,J,\text{ 더 높은 형식들})$을 법으로 이해한다. 이를 매번 명시하지는 않는다.

% inherited: noether_paper02_ko_translation_001_20260804\targets\Noether_P02_Korean_T19_U127_UNCHECKED.tex
% sha256: EEF48AF46C569766166923D92AEE234CBB454BAC2C7D0913D2B2F332E4156528
환원을 위하여 앞서 얻은 공식들을 모아 적는다.
\srcnumdisplay{(24.)}{%
\begin{aligned}
s_\nu^2&=\frac43i\,a_\nu^2+\frac16J\cdot u_x^2-\frac29i^2\cdot u_x^2,\\
s_\nu^3&=\frac13J\cdot u_x,
& s_\nu^4&=J,\\
g&=2s_\vartheta^2-\frac43i\Delta,\\
s_\vartheta^2(\theta su)&=0,\\
s_\vartheta^2(\theta su)^2&=\frac89i\cdot a_\nu^2-\frac4{27}i^2\cdot u_x^2,\\
Z&=2i\cdot a_\eta^2+\frac43i(asu)^2-\frac13J\cdot\theta,\\
g_\nu&=-s_\eta+u_x\left\{2ia_\eta^2-\frac23i(asu)^2+\frac23J\cdot\theta\right\},\\
g_\nu^2&=\frac45ia_\eta^2-\frac25i(asu)^2+\frac13J\cdot\theta,
&g_\nu^3&=0,\quad g_\nu^4=0.
\end{aligned}
}

% inherited: noether_paper02_ko_translation_001_20260804\targets\Noether_P02_Korean_T19_U128_UNCHECKED.tex
% sha256: DDBAE354DAB1CEBA94BC31917DE2F3890B732A766DD8045EC9D5CC53FB592417
\srcfntext{*)}{간결하게 쓰기 위하여, 완전한 전이연산 $[(\theta Hu)^2]_{\widehat{s}u}s$ 대신 전이연산의 항을 계의 형식으로 채택한다. 공식을 세울 때에는 모든 전이연산을 그 선두항으로 바꾼다. 예를 들어 $(\theta_\eta(\theta Hu),s,u)^2$을 $\theta_\eta(\theta Hu)(\theta su)^2$으로 바꾼다.}

% inherited: noether_paper02_ko_translation_001_20260804\targets\Noether_P02_Korean_T19_U129_UNCHECKED.tex
% sha256: 5B287485886137270075B0639220010D7F7650D2DF3DBBE030D4C4D22140A89F
\noindent\emph{결론:}
\[
 s_\nu^2,\quad s_\vartheta^2(\theta su),\quad s_\vartheta^2s_\nu,\quad s_\vartheta^2\theta_\nu
\]
은 환원자이다.

% inherited: noether_paper02_ko_translation_001_20260804\targets\Noether_P02_Korean_T20_U135_UNCHECKED.tex
% sha256: 376131ED89288A5187E55255EBCF4598B96B7FCD69BE75F4B5AA6C2843664D53
\begin{center}
\textbf{\S\ 18.}\quad\emph{계 $s_I$.}
\end{center}

\[
\text{다음의 수축: }\left\{\begin{array}{l}
s\text{와 } f;\ \theta;\ K,j,\Delta;\ f\cdot j,N;\ \Delta\cdot j,\\
s^2\text{와 }\theta,K.
\end{array}\right.
\]

% inherited: noether_paper02_ko_translation_001_20260804\targets\Noether_P02_Korean_T20_U136_UNCHECKED.tex
% sha256: F0C237C3EB7E17C5E996B8A40AB41D21FD61F8E6DF77ED06B9DAFF8C8772AFB6
\noindent\emph{환원 불가능 형식들.}

\noindent 계수차수 5:
\[
(asu);\qquad (asu)^2;\qquad (asu)^3;\qquad (asu)^4.
\]

% inherited: noether_paper02_ko_translation_001_20260804\targets\Noether_P02_Korean_T20_U137_UNCHECKED.tex
% sha256: 6A94BB9A355F4F3BB190A42C020CEAB83A93DBC206D4DE137D3BC8F19163949D
\noindent 계수차수 6:
\[
\begin{array}{cccc}
(\theta su)&(\theta su)^2&(\theta su)^3&(\theta su)^4\\
s_\vartheta&s_\vartheta(\theta su)&s_\vartheta(\theta su)^2&s_\vartheta(\theta su)^3\\
\cdot&s_\vartheta^2&\cdot&\cdot
\end{array}
\]

% inherited: noether_paper02_ko_translation_001_20260804\targets\Noether_P02_Korean_T20_U138_UNCHECKED.tex
% sha256: A3521B2C7D4EDB9D1FC23DE2341A103F51BD209206FAE366D04BB24C79C69651
\noindent 계수차수 7:
\[
\begin{array}{c@{\quad}c@{\quad}c@{\quad}c@{\qquad}l}
\cdot&(Ksu)^2&(Ksu)^3&(Ksu)^4&\\
\text{A)}\ s_k&s_k(Ksu)&s_k(Ksu)^2&s_k(Ksu)^3&
\text{B)}\ s_j,\\
\cdot&s_k^2&\cdot&\cdot&
\text{C)}\ (\Delta su),\ (\Delta su)^2.
\end{array}
\]

% inherited: noether_paper02_ko_translation_001_20260804\targets\Noether_P02_Korean_T20_U139_UNCHECKED.tex
% sha256: 2773F82ADA7406321F5159566936CE52F166EDF98E354B50E0192F65175B1FC6
\noindent 계수차수 8:
\[
\begin{array}{ll}
\text{A)}\ s_j(asu),\ s_j(asu)^2,\ s_j(asu)^3,
&
\text{B)}\ \begin{array}[t]{cc}
\cdot&(Nsu)^2\\
s_n&s_n(Nsu).
\end{array}
\end{array}
\]

% inherited: noether_paper02_ko_translation_001_20260804\targets\Noether_P02_Korean_T20_U140_UNCHECKED.tex
% sha256: DD1FA79FF4BA3E2FD87DAE9277AD16E16B505EF6E6FBE94849AB23AE27EA2807
\noindent 계수차수 10:
\[
\text{A)}\ s_j(\Delta su),\qquad \text{B)}\ s_\vartheta(\theta s'u)^4.
\]

% inherited: noether_paper02_ko_translation_001_20260804\targets\Noether_P02_Korean_T20_U141_UNCHECKED.tex
% sha256: 428C24DCEFC297E62DCB7D57DA646C237D244C2FE8FEDEAE0E2A87D76162F3FE
\noindent 계수차수 11:
\[
\begin{array}{cc}
(Ks^2u)^6&\cdot\\
s_k(Ks^2u)^5&s_k(Ks^2u)^6.
\end{array}
\]

% inherited: noether_paper02_ko_translation_001_20260804\targets\Noether_P02_Korean_T20_U142_UNCHECKED.tex
% sha256: B590E3B0BB11CFE144C2CA86B8F4841CD9E88D3D6F2D79D04094D5F51B173D39
\noindent\emph{환원:}

다음 형식열들은
\[
 s_\vartheta^2(\theta su),\quad s_k^2(Ksu),\quad s_j^2,\quad (\Delta su)^3,\quad (Nsu)^3
\]
환원자 $s_\vartheta^2(\theta su)$를 인수로 가지며, $s_j^2$과 $(\Delta su)^3$은 $j$와 $\Delta$를 형식열의 더 낮은 형식들로 되돌려서 얻는다:
\[
\begin{aligned}
s_\vartheta^2(\theta su)(a\theta u)^3&=-\frac12s_j^2,\\
s_\vartheta^2(\theta su)(a\theta u)^2&=\frac13(\Delta su)^3,\\
s_\vartheta^2(\theta su)^2(a\theta u)^2&=\frac13(\Delta su)^4+\frac13s_j^2.
\end{aligned}
\]

% inherited: noether_paper02_ko_translation_001_20260804\targets\Noether_P02_Korean_T20_U143_UNCHECKED.tex
% sha256: 7225071139AC64C0FF78CDEBC1518360D593AE1A7C46C03517A2C5DCECDED1BE
형식 $s_\vartheta^2s_{\vartheta'}$와 $s_\vartheta^2s_{\vartheta'}^2$: 환원자 $s_\vartheta^2(\theta su)$와 $\theta_{\vartheta'}$:
\[
\begin{aligned}
s_\vartheta^2s_{\vartheta'}&=s_\vartheta^2\theta_{\vartheta'}-s_\vartheta^2(\theta s\widehat{\vartheta'}x),\\
s_\vartheta^2s_{\vartheta'}^2&=s_\vartheta^2s_{\vartheta'}\theta_{\vartheta'}-s_\vartheta^2s_{\vartheta'}(\theta s\widehat{\vartheta'}x).
\end{aligned}
\]
형식열 $s_j(\Delta su)^2$: 환원자 $\Delta_j$와 $(\Delta su)^3$.

% inherited: noether_paper02_ko_translation_001_20260804\targets\Noether_P02_Korean_T20_U144_UNCHECKED.tex
% sha256: 828247B6ECE3D556B9436B6BC04EE54C230C0E6C72CBFC4F4A068E27C95DAD42
\noindent\emph{결론:}

$s$는 거듭제곱으로는 $f,j,\Delta,N$과 수축되지 않는다. $f,j,\Delta,N$은 반변 형식들과의 곱 안에서만 수축되는데, $f$와의 곱에서는 그 반변 형식이 $x$에 관해 여전히 이차일 수 있고, $\Delta$와의 곱에서는 $x$에 관해 여전히 일차일 수 있다. 즉 다음 형식들의 곱을 고찰해야 한다:
\[
 f\cdot(0,\lambda),\quad f\cdot(1,\lambda),\quad f\cdot(2,\lambda),\quad
 j\cdot(\mu,0),\quad N,\quad
 \Delta\cdot(0,\lambda),\quad \Delta\cdot(1,\lambda)\quad(\lambda>0),
\]
여기서 $(\mu,\lambda)$는 $x$에 관해 $\mu$차이고 $u$에 관해 $\lambda$차인 형식을 뜻한다.

% inherited: noether_paper02_ko_translation_001_20260804\targets\P02_T20_U145_r2.tex
% sha256: 74E9293CBBF116833E2C9DD705D3253D61660440A6A52CE302F201EC2A1914A8
$\theta$ 또는 $K$는 임의의 형식과 이루는 거듭제곱이나 곱 안에서 $s$ 또는 $s^2$과 수축되지 않는다. 실제로 야코비 행렬식 $K$와의 곱 $K\cdot\varphi_x$ ($\varphi$는 $f$도 $\theta$도 아님)에 대해서는 분해 가능 형식들 사이의 다음 관계가 성립한다:
\[
K\cdot\varphi_x=(a\theta u)\varphi_x=f\cdot(\varphi\theta u)-\theta\cdot(\varphi a u).
\]
따라서 위에 제시한 계 $s_I$의 수축 도식이 증명되었다.

% inherited: noether_paper02_ko_translation_001_20260804\targets\Noether_P02_Korean_T21_U149_UNCHECKED.tex
% sha256: EC5ED1936725BDEF6AFC4417CEE55F894110EFC02D7FE78845AB61EF1CFAB8B2
\begin{center}
\textbf{\S\ 19.}\quad\emph{계 $s_{II}$.}
\end{center}

$s$와 $\nu;\ H;\ L;\ H^2$의 수축.

\noindent\emph{환원 불가능 형식들:}
\[
\begin{array}{c@{\qquad}c@{\qquad}c@{\qquad}c}
\text{계수차수 6} & \text{계수차수 8} & \text{계수차수 10} & \text{계수차수 12}\\[0.3em]
s_\nu & (Hsu),\ (Hsu)^2 & (Lsu)^2,\ (Lsu)^3 & (H^2su)^3\\
\cdot & s_\eta & s_l & \cdot\\
\cdot & \cdot & \cdot & \cdot\\
s_\nu^4=J & \cdot & \cdot & \cdot
\end{array}
\]

% inherited: noether_paper02_ko_translation_001_20260804\targets\Noether_P02_Korean_T21_U150_UNCHECKED.tex
% sha256: 0EF7C35345D3E4AEADB7B9BC895431BBE2BD0E4550F032C829E08042E2A7951B
\noindent\emph{환원:}

형식열 $s_\eta(Hsu)$와 $s_l(Lsu)$: 환원자 $s_\nu^2$.

형식열 $s_\sigma$: $H$에서 얻는 환원자 $s_\nu^2$,
$s_\nu^2=\dfrac23s_\sigma$.

$(H^2su)^4$은 공식 (7.)a에서 분해된다(\S\ 11. A 참조). 그로부터 $(H\cdot L,s,u)^4$의 분해가 따른다.

% inherited: noether_paper02_ko_translation_001_20260804\targets\P02_T21_U151_r2.tex
% sha256: 195894EEA219FF2EE48D6498C3F7FAE02F44AE59EC35B4791A6E94D94BCC2A63
\noindent\emph{결론:}

$s^2$은 계 II와 수축되지 않는다.

$\nu$는 반변 형식과의 곱에서만 수축된다. $H$는 공변식으로 보았을 때 형식 $(1,\lambda)$, $(2,\lambda)$, $(3,\lambda)$ ($\lambda$는 임의)과의 곱에서 수축된다.

$\sigma$와 $(\nu\sigma x)$은 전혀 수축되지 않으며, 야코비 행렬식으로서의 $L$은 곱 안에서 수축되지 않는다(계 III의 공변식들에 대한 완전 전이연산은 환원 가능해진다).

계 I과 II의 곱으로 남는 것은 다음과 같다:
\[
f\cdot \nu,\qquad \Delta\cdot \nu.
\]

% inherited: noether_paper02_ko_translation_001_20260804\targets\Noether_P02_Korean_T22_U152_UNCHECKED.tex
% sha256: 5396CF17931084BC8D3A279E408E8BE6AE5D64E7AA682080602A0B93A290EA48
\begin{center}
\textbf{\S\ 20.}\quad\emph{계수차수 7의 형식들(계 $s_{III}$).}
\end{center}

$s$와 $f\cdot \nu$, $a_\nu$, $a_\nu^2$의 수축.

\noindent\emph{환원 불가능 형식들:}
\[
\begin{gathered}
s_\nu(asu),\quad a_\nu(asu),\quad s_\nu(asu)^2,\quad a_\nu(asu)^2,\quad
s_\nu(asu)^3,\quad a_\nu(asu)^3,\\
a_\nu s_\nu,\qquad a_\nu s_\nu(asu),\qquad a_\nu^2(asu),\qquad a_\nu^2(asu)^2.
\end{gathered}
\]

% inherited: noether_paper02_ko_translation_001_20260804\targets\P02_T22_U153_r2.tex
% sha256: 69C7BE7574D4881E7A1F0E0382D8154E710BF5A479289EA84AF1AF90F861C0E4
\noindent\emph{환원:}

환원자 $s_\nu^2$과 $s_\eta(Hsu)$와 직접 수축하여 생기는 자명한 환원들은 언급하지 않으며, 야코비 행렬식과의 전이연산이 분해되는 경우도 마찬가지로 언급하지 않는다.

형식 $a_\nu s_\nu(asu)^2$과 $a_\nu^2s_\nu(asu)^2$: 공식 (19.), \S\ 12 참조;
\[
a_\nu s_\nu(asu)^3=a_\nu(a\widehat{\nu_1}\nu_2u)^3(\nu_1\nu_2\nu)=0.
\]

% inherited: noether_paper02_ko_translation_001_20260804\targets\Noether_P02_Korean_T22_U154_UNCHECKED.tex
% sha256: 1CADA308602F60D4A5EBDAE0C819176DF2B0AEC2380BA8CA1A21663435B1C781
형식열 $a_\nu^2s_\nu$: $a_\nu^2$을 더 낮은 형식들로 되돌림으로써 얻는 환원자 $s_\vartheta^2(\theta su)$, 즉 다음 표현들을 이중 환원함으로써 얻는다:
\[
s_\vartheta^2(a\theta s)(a\theta u),\qquad
s_\vartheta^2(a\theta s)(a\theta u)^2
\]
다음과 같이 둔다:
\[
\begin{aligned}
s_\vartheta^2(a\theta s)(a\theta u)
&=[(a\theta u)^2]s^3+\frac15[a_\vartheta(a\theta u)^2]s^2+\frac1{60}[a_\vartheta^2(a\theta u)^2]s\\
&=\frac12a_\nu^2s_\nu-\frac23a_\nu s_\nu^2,\\
s_\vartheta^2(a\theta s)(a\theta u)^2
&=\frac45[a_\vartheta(a\theta u)^2]_{\widehat{us}}s^2+\frac{23}{180}[a_\vartheta^2(a\theta u)^2]_{\widehat{us}}s\\
&=-\frac12a_\nu^2s_\nu(asu).
\end{aligned}
\]
이때 $a_\nu^2b_\nu(abu)=0$이다.

% inherited: noether_paper02_ko_translation_001_20260804\targets\Noether_P02_Korean_T22_U155_UNCHECKED.tex
% sha256: 008B8B7EA65B400BCB73D1B46D23BEAB8C20E359EC1DFF180001760584CB1BD7
다음 공식들을 얻는다:
\srcnumdisplay{(25.)}{%
\begin{array}{rcl@{\qquad}rcl}
\underline{a_\nu^2s_\nu}&=&\dfrac43i\cdot a_\nu^2b_\nu,
& \underline{a_\nu s_\nu(asu)^2}&=&\dfrac13i\cdot\theta_\nu^2-\dfrac1{18}iH,\\[0.7em]
\underline{a_\nu s_\nu(asu)^3}&=&0,
& \underline{a_\nu^2s_\nu(asu)}&=&0,\\[0.7em]
\underline{a_\nu^2s_\nu(asu)^2}&=&0.
\end{array}
}

% inherited: noether_paper02_ko_translation_001_20260804\targets\Noether_P02_Korean_T22_U156_UNCHECKED.tex
% sha256: 6395D2A14C1C6BF863C0C10B948FF9E75ABF6E776C18F4FD445D024E0003DD8F
\noindent\emph{결론:}

1) $a_\nu s_\nu(asu)^2$, $a_\nu^2s_\nu$은 환원자이다.

2) \S\ 19에서 환원 가능하다고 판정된 형식들 $(\Delta su)^3$, $(\Delta su)^4$의 값들을 얻는다. 대응하는 형식열 $(agu)^2$에 대해서도 마찬가지이며, 이 형식열은 $\nu$와 수축할 때 환원자 $g_\nu$에 의해 이중 환원을 일으킨다.
\srcnumdisplay{(26.)}{%
\begin{aligned}
\text{a)}\quad (\Delta su)^3&=-\frac65a_\nu^2(asu)+3a_\nu s_\nu(asu)+2i\theta_\nu(\vartheta \nu x),\\
\text{b)}\quad (\Delta su)^4&=-\frac25a_\nu^2(asu)^2+\frac43i\theta_\nu^2+\frac19iH.
\end{aligned}
}

% inherited: noether_paper02_ko_translation_001_20260804\targets\Noether_P02_Korean_T22_U157_UNCHECKED.tex
% sha256: E5530DDF69564CA4C15FAAD74493D94D4AAC08F86F5525C96FF7628B94CD6732
공식 (24.)와 (26.)에서, $(i,J)$를 법으로 하여(71쪽 참조):
\srcnumdisplay{(27.)}{%
\begin{aligned}
\text{a)}\quad (agu)^2&=\frac23(\Delta su)^2-\frac43a_\nu s_\nu+\frac35s\cdot a_\nu^2,\\
\text{b)}\quad (agu)^3&=2a_\nu s_\nu(asu)-a_\nu^2(asu)^2,\\
\text{c)}\quad (agu)^4&=-a_\nu^2(asu)^2.
\end{aligned}
}

% inherited: noether_paper02_ko_translation_001_20260804\targets\Noether_P02_Korean_T23_U158_UNCHECKED.tex
% sha256: 6A25CA8D10B9FF62381ED31D075E0A4328E05E784A231970E3E57BE4FB264B73
\begin{center}
\textbf{\S\ 21.}\quad\emph{계수차수 8의 형식들(계 $s_{III}$).}
\end{center}

$s$와 계수차수 4의 형식들(계 III)의 수축. (\S\ 9.)

\noindent\emph{환원 불가능 형식들:}
\[
\begin{array}{c@{\quad}c|c|c}
\cdot
& (\vartheta \nu x)s_\nu,\ ((\vartheta \nu x)^2,s,u)
& ((\vartheta \nu x),s,u)^2,\ (\theta_\nu su)
& \bigstar\\[0.45em]
& (\vartheta \nu x)^2s_\nu
& ((\vartheta \nu x)^2,s,u)^2,\ \theta s_\nu
& \\[0.45em]
&& ((\vartheta \nu x)^2,s,u)s_\nu,\ (\theta_\nu(\vartheta \nu x)^2,s,u)
&
\end{array}
\]

% inherited: noether_paper02_ko_translation_001_20260804\targets\Noether_P02_Korean_T23_U159_UNCHECKED.tex
% sha256: E56F4CA025E1C1879B200159D7D9113C586A8BAB332F259C308273FA242D060C
\[
\begin{array}{c|c|c|c}
\bigstar
& ((\vartheta \nu x),s,u)^3,\ (\theta_\nu su)^2
& ((\vartheta \nu x),s,u)^4,\ (\theta_\nu su)^3
& \cdot\\[0.45em]
& ((\vartheta \nu x)^2,s,u)^3,\ (\theta_\nu(\vartheta \nu x),s,u)^2,\ (\theta_\nu^2su)
& (\theta_\nu(\vartheta \nu x),s,u)^3,\ (\theta_\nu^2su)^2
& (\theta_\nu(\vartheta \nu x),s,u)^4\\[0.45em]
& \cdot
& (\theta_\nu^3su)
& \cdot
\end{array}
\]

% inherited: noether_paper02_ko_translation_001_20260804\targets\P02_T23_U160_r2.tex
% sha256: ED1FA9FA9C95277456885B355E7679A8784C9DF10D512A6796AC07BAFB95E06E
\noindent\emph{환원:}

$((\vartheta \nu x),s,u)s_\nu$은 \S\ 3에 따라 야코비 행렬식에 대한 전이연산으로 분해된다.

형식열 $((\vartheta \nu x),s,u)^2s_\nu$: 환원자 $s_\nu^2$과 $s_\vartheta^2\theta_\nu$:
\[
(\widehat{\vartheta \nu su})s_\nu(\theta su)=s_\nu s_\vartheta(\theta su)=-\frac12(\widehat{\vartheta \nu su})^2(\theta su),
\]
\[
\theta_\nu(\widehat{\vartheta \nu su})s_\nu=\theta_\nu s_\nu s_\vartheta=-\frac12\theta_\nu(\widehat{\vartheta \nu su})^2.
\]
형식열 $\theta_\nu s_\nu(\theta su)$: 다음 표현들을 이중 환원하여 얻는 환원자 $a_\nu s_\nu(asu)^2$:
\[
a_\nu s_\nu(asu)^2(abu)\quad\text{및}\quad a_\nu s_\nu(asu)^2(abu)(sbu);
\]
$\theta_\nu s_\nu(\theta su)^3$은 $(agu)^3(abu)(bgu)^3$을 이중 환원하여 얻는다.

% inherited: noether_paper02_ko_translation_001_20260804\targets\P02_T23_U161_r2.tex
% sha256: 8D440D5E6DF7163DA596DF2D1308E3AC6A78C04CC3E0A96467C25AE376673B85
형식열 $\theta_\nu(\vartheta \nu x)s_\nu$: 등식 $\theta_\nu(\vartheta \nu x)s_\nu=a_\nu^2(abu)s_\nu$에서 얻는 환원자 $a_\nu^2s_\nu$.

형식열 $(\theta_\nu(\vartheta \nu x)^2,s,u)$: 형식열 $\theta_\nu(\widehat{\vartheta \nu su})s_\nu$의 형식들 가운데 동시에 형식열 $((\vartheta \nu x)su)^2s_\nu$과 $\theta_\nu(\vartheta \nu x)s_\nu$에 속하는 것들을 이중 환원한다.

형식 $(\theta_\nu(\vartheta \nu x),s,u)$은 야코비 행렬식 $\theta_\nu(\vartheta \nu x)=a_\nu^2(abu)$에 대한 한 번의 전이연산으로 분해된다.

형식 $((\vartheta \nu x)^2,s,u)^4$: 표현 $(a\A u)^2(\A su)^4$을 이중 환원한다. 또는 공식 (27.)에서 나온 $(\theta gu)^4$를 $(j\eta x)^4$의 환원(\S\ 12 C))과 유사하게 이중 환원한다.

% inherited: noether_paper02_ko_translation_001_20260804\targets\Noether_P02_Korean_T23_U162_UNCHECKED.tex
% sha256: 9C24BC3ED7A4EA7F703CFBB4A165520ED55986A9A852B3B3250D43CC2A9078FF
이중 환원을 통하여 다음 공식들을 얻는다:
\srcnumdisplay{(28.)}{%
\begin{aligned}
0&=\theta_\nu s_\nu(\theta su)-\frac16(\widehat{\vartheta \nu su})^2(\theta su)-\frac1{12}(Hsu),\\
0&=\theta_\nu s(\theta su)^2-\frac14(\widehat{\vartheta \nu su})^2(\theta su)^2-\frac1{24}(Hsu)^2,\\
0&=(\widehat{\vartheta \nu su})^2(\theta su)^2+2\theta_\nu(\widehat{\vartheta \nu su})(\theta su)^2
  +3\theta_\nu^2(\theta su)^2-\frac13H(su)^2,\\
0&=\theta_\nu s_\nu(\theta su)^3+\frac12\theta_\nu(\widehat{\vartheta \nu su})(\theta su)^3,\\
0&=\theta_\nu s_\nu s_\vartheta-\frac14s_\eta,
\quad 0=\theta_\nu s_\nu s_\vartheta(\theta su),
\quad 0=\theta_\nu s_\nu s_\vartheta(\theta su)^2.
\end{aligned}
}
마지막 것은 $((\theta_\nu(\vartheta \nu x)^2,s,u)^2,\ldots)$에 대응한다.

% inherited: noether_paper02_ko_translation_001_20260804\targets\Noether_P02_Korean_T23_U163_UNCHECKED.tex
% sha256: 058DC86FD6EFDF3CA75F584824409C3E09E2B90A17259E4D6AF08E136FC6B8A6
\noindent\emph{결론:}

1) $\theta(\vartheta \nu x)s_\nu$, $(\widehat{\vartheta \nu su})(\theta su)s_\nu$, $\theta_\nu(\widehat{\vartheta \nu su})^2$, $(\widehat{\vartheta \nu su})^2(\theta su)^2$은 환원자이다.

2) 계수차수 4의 형식들 $(5,\lambda)$, $(6,\lambda)$ 등은 $s^2$과의 수축에 들어가지 않는다.

% inherited: noether_paper02_ko_translation_001_20260804\targets\Noether_P02_Korean_T23_U164_UNCHECKED.tex
% sha256: 696825DFBF4B1A36AFE22CC7F4F48E189698DEE2FCF21940840E05ED574DC03C
3) 형식열 $(\theta gu)^2$에 대해서는 공식 (27.)과 이 절의 환원들을 고려함으로써 다음 공식들을 얻는다:
\srcnumdisplay{(29.)}{%
\begin{aligned}
\text{a)}\quad (\theta gu)^2&=\frac43\theta_\nu s_\nu+\frac23s_\nu s_\vartheta-\frac13s\cdot\theta_\nu^2-\frac19s\cdot H-\frac13f\cdot a_\nu^2(asu)^2+\frac13a_\nu^2\cdot(asu)^2,\\
\text{b)}\quad (\theta gu)^3&=-\frac13(\widehat{\vartheta \nu su})^2(\theta su)-\theta_\nu(\widehat{\vartheta \nu su})(\theta su)-\theta_\nu^2(\theta su),\\
\text{c)}\quad (\theta gu)^4&=2\theta_\nu^2(\theta su)^2-\frac13(Hsu)^2,\\
g\vartheta(\theta gu)&=(\widehat{\vartheta \nu su})(\vartheta \nu x)s_\nu-\theta_\nu(\vartheta \nu x)(\widehat{\vartheta \nu su}),\\
g\vartheta(\theta gu)^2&=\frac23s\cdot\theta_\nu^3,
\quad g\vartheta(\theta gu)^3=\frac13\theta_\nu^3(\theta su),
\quad g\vartheta(\theta gu)^4=0.
\end{aligned}
}

% inherited: noether_paper02_ko_translation_001_20260804\targets\Noether_P02_Korean_T24_U170_UNCHECKED.tex
% sha256: 3DF506BDB3C267420EE0B03544BBD01708FD31CE57BB5F988FE46E7B0963199D
\begin{center}
\textbf{\S\ 22.}\quad\emph{계수차수 9의 형식들(계 $s_{III}$).}
\end{center}

$s$와 계수차수 5의 형식들(계 III) 및 곱 $\A\cdot \nu$의 수축(\S\ 10).

\noindent\emph{환원 불가능 형식들:}
\[
\begin{array}{c@{\quad}c|c|c|c}
\text{A)}& \cdot & \cdot & (K_\nu su)^2 & \bigstar\\[0.45em]
& \cdot & ((k\nu x)^2,s,u)^2,\ K_\nu s_\nu & ((k\nu x)^2,s,u)^3 & \\
& (k\nu x)^2s_\nu,\ ((k\nu x)^3,s,u) & ((k\nu x)^3,s,u)^2 & \cdot & \cdot\\[0.45em]
& (k\nu x)^3s_\nu & ((k\nu x)^3,s,u)s_\nu,\ (K_\nu(k\nu x)^3,s,u) & \cdot & \cdot\\[0.45em]
& \bigstar\ (K_\nu su)^3 & (K_\nu su)^4 & \cdot & (K_\nu^2su)^4
\end{array}
\]
\[
\begin{array}{c@{\quad}c|c@{\qquad}c@{\quad}c}
\text{B)}& (j\nu x)s_\nu,\ ((j\nu x)^2,s,u) & ((j\nu x)^2,s,u)^2 && \\
& ((j\nu x)^3,s,u) & \cdot && \text{C)}\quad s_\nu(\A su),\ (\A_\nu su),
\end{array}
\]
\[
\begin{array}{c@{\quad}c|c|c}
\text{D)}& \cdot & ((aHu),s,u)^2,\ ((aHu)^2,s,u) & \bigstar\\[0.45em]
& (aHu)s_\eta,\ (a_\eta su) & (aHu)^2s_\eta,\ (a_\eta su)^2 & \\
& \cdot & (a_\eta^2su) & \\
\bigstar& ((aHu),s,u)^3,\ ((aHu)^2,s,u)^2 & ((aHu),s,u)^4 & \\
& (a_\eta su)^3,\ (a_\eta(aHu),s,u)^2,\ (a_\eta(aHu)^2,s,u) & (a_\eta(aHu),s,u)^3 &
\end{array}
\]

% inherited: noether_paper02_ko_translation_001_20260804\targets\Noether_P02_Korean_T24_U171_UNCHECKED.tex
% sha256: 0AD2F19D6D62EDA9A2F24445BDAA408AB705230B4316451CE4F24F5079DAF690
\noindent\emph{환원:}

A) 환원들은 한 번 수축하여 \S\ 21의 환원들에서 직접 따른다. 형식열 $K_\nu s_\nu(Ksu)^2$은 환원자 $a_\nu s_\nu(asu)^2$과 $\theta_\nu s_\nu(\theta su)^2$을 인수로 갖는다. 이로부터 다음과 같이 두면 더 높은 형식 $K_\nu^2(Ksu)^2$, $K_\nu^2(Ksu)^3$이 분해된다:
\[
0=a_\nu s_\nu(asu)^2(a\theta u)=K_\nu s_\nu(Ksu)^2=\theta_\nu s_\nu(\theta su)^2(a\theta u).
\]

% inherited: noether_paper02_ko_translation_001_20260804\targets\Noether_P02_Korean_T24_U172_UNCHECKED.tex
% sha256: 11715A2644F04045375DAF25BEDF51E5AB92B22074FA4BDAFE341A5CE2D8987A
B) 형식열 $(j\nu x)^2s_\nu$: 다음 식에서 얻는 환원자 $a_\nu^2s_\nu$:
\[
(j\nu x)^2s_\nu=3a_\nu^2s_\nu(a\theta u)^2.
\]
$((j\nu x)^3,s,u)^2$: 다음 식에서 얻는 환원자 $a_\nu^2s_\nu$:
\[
((j\nu x)^3,s,u)^2=-6a_\nu^2s_\nu(a\theta u)^2(\theta su).
\]

% inherited: noether_paper02_ko_translation_001_20260804\targets\Noether_P02_Korean_T24_U173_UNCHECKED.tex
% sha256: D3EE7979786937492EE5940121A2079F9851F53B2E461B959CF6824E9EE56263
C) 형식 $\A_\nu(\A su)^2$과 $s_\nu(\A su)^2$: 공식 (27.)a에서 얻는 환원자 $g_\nu$.

형식 $\A_\nu s_\nu$: 1) $a_ga_\nu^2s_\nu=\frac23\A_\nu s_\nu$에서 얻는 환원자 $a_\nu^2s_\nu$;

2) 공식 (27.)a에 따라 표현식 $(\A s\widehat{\nu}x)^2$을 이중 환원함으로써 분해된다.

그로부터 더 높은 형식 $a_\eta s_\eta$의 분해가 따른다.

% inherited: noether_paper02_ko_translation_001_20260804\targets\P02_T24_U174_r2.tex
% sha256: 26E869C10316B6936728DED4A5BCFA3F5E1FBB40D79D0537D3ACCC9FC8AD11C3
D) 형식열 $(aHu)^2s_\eta(asu)$: 형식열 $[a_\nu s_\nu(asu)^2]_{\widehat{\nu_1}x}$을 이중 환원하여 얻는 환원자 $a_\nu s_\nu(asu)^2$; $a_\nu s_\nu(asu)^2$과 $a_\nu s_\nu(asu)^3$에서 $(aHu)^2s_\eta(asu)^2$의 환원을 얻는다. 그로부터 $(a_\eta,s,u)^4$의 환원이 따른다.

형식열 $a_\eta(aHu)s_\eta$: 표현식 $(\A s\widehat{\nu}x)^3$을 이중 환원하여 얻는 환원자 $a_\nu^2s_\nu$, $a_\eta^2s_\eta$. 이는 $a_\eta^2s_\eta$가 형식 $a_\eta(aHu)s_\eta$의 환원 가능 항 $a_\nu^2s_\nu(\nu\nu_1x)$에서 수축으로 생기지 않기 때문이다.

형식열 $(a_\eta^2su)^2$: 다음 식에서 얻는 환원자 $a_\nu^2s_\nu$:
\[
a_\nu^2s_\nu s_{\nu_1}=-\frac12a_\eta^2(Hsu)^2.
\]
형식 $(a_\eta(aHu),s,u)$은 야코비 행렬식에 대한 전이연산으로서 분해된다.

% inherited: noether_paper02_ko_translation_001_20260804\targets\Noether_P02_Korean_T24_U175_UNCHECKED.tex
% sha256: 81508A22D41283D5370C348613FAAB53E5D079F45A78E8853B77369B35F7A5A4
다음 환원 공식들을 얻는다:
\srcnumdisplay{(30.)}{%
\begin{aligned}
\text{a)}\quad 0&=(aHu)^2(asu)s_\eta-a_\eta(asu)(Hsu)^2-a_\eta(aHu)(asu)(Hsu),\\
\text{b)}\quad 0&=(aHu)^2(asu)^2s_\eta-a_\eta(aHu)(asu)^2(Hsu),\\
\text{c)}\quad 0&=a_\eta(asu)^2(Hsu)^2,\\
0&=a_\eta s_\eta+\frac14a_\nu^2\cdot g-\frac14s\cdot a_\eta^2,\\
0&=a_\eta(aHu)s_\eta-\frac12a_\eta^2(Hsu),
\quad 0=a_\eta^2(Hsu)^2,\ldots,
\quad 0=a_\eta^2s_\eta.
\end{aligned}
}

% inherited: noether_paper02_ko_translation_001_20260804\targets\Noether_P02_Korean_T24_U176_UNCHECKED.tex
% sha256: 7A87091496016DDF31C8265B3E1A76FDA7A7E953A4AA9668730ED4058550FEEA
\noindent\emph{결론:}

1) $(j\nu x)^2s_\nu$, $\A_\nu s_\nu$, $\A_\nu(\A su)^2$, 따라서 $N_\nu(Nsu)^2$, $(aHu)^2(asu)s_\eta$, $a_\eta(aHu)s_\eta$, $a_\eta^2(Hsu)^2$은 환원자이다. $a_\eta s_\eta$는 분해 가능 형식이며, 모든 한 번 및 두 번의 수축에서 분해 가능 형식으로 옮겨간다.

2) 계수차수 5의 형식들 $(5,\lambda)$, $(6,\lambda)$ 등은 $s^2$과 수축되지 않는다.

% inherited: noether_paper02_ko_translation_001_20260804\targets\Noether_P02_Korean_T25_U177_UNCHECKED.tex
% sha256: 1112C5D9D67AA85C49036F87C73A67C28C76D86E2E8AB354CB2C18BAE3FD7D96
\begin{center}
\textbf{\S\ 23.}\quad\emph{계수차수 10의 형식들(계 $s_{III}$).}
\end{center}

$s$와 계수차수 6의 형식들(계 III)의 수축(\S\ 11).

\noindent\emph{환원 불가능 형식들:}
\[
\begin{array}{r@{\quad}c}
\text{A)}&
\begin{array}{c|c}
(\vartheta^2\nu x)^3s_\nu&\cdot\\[0.45em]
\cdot&(\vartheta^2\nu x)^3s_\nu s_\vartheta,\ \theta_\nu(\vartheta^2\nu x)^3s_\vartheta
\end{array}
\end{array}
\]

% inherited: noether_paper02_ko_translation_001_20260804\targets\Noether_P02_Korean_T25_U178_UNCHECKED.tex
% sha256: 02F7DE21629492B5D5743D2A74610699A92F9726204FF981488718EF926DD795
\[
\begin{array}{r@{\quad}c}
\text{B)}&
\begin{array}{c|ccc}
\cdot&(a_\nu(j\nu x),s,u)^2&(a_\nu(j\nu x),s,u)^3&(a_\nu(j\nu x),s,u)^4\\[0.45em]
a_\nu(j\nu x)s_\nu&\cdot&(a_\nu^2(j\nu x),s,u)^2&(a_\nu^2(j\nu x),s,u)^3
\end{array}
\end{array}
\]

% inherited: noether_paper02_ko_translation_001_20260804\targets\Noether_P02_Korean_T25_U179_UNCHECKED.tex
% sha256: 668639461A1396A545BD338BD9E7F73DD4BA0369D5BBD81A78DFFE8EEEE90717
\[
\begin{array}{r@{\quad}c}
\text{C)}&
\begin{array}{c|c|c|c}
\cdot&\cdot&((\theta Hu),s,u)^2,\ ((\theta Hu)^2,s,u)&\bigstar\\[0.35em]
\cdot&((\vartheta\eta x),s,u)^2,\ (\theta Hu)s_\eta,\ (\theta_\eta su)&((\vartheta\eta x),s,u)^3,\ (\theta_\eta su)^2&\\[0.35em]
((\vartheta\eta x)^2,s,u)&\cdot&(\theta_\eta^2su)&\\[0.35em]
\cdot&(\theta_\eta(\vartheta\eta x)^2,s,u)&\cdot&
\end{array}
\end{array}
\]
\[
\begin{array}{c@{\quad}|c|c}
\bigstar&((\theta Hu),s,u)^3,\ ((\theta Hu)^2,s,u)^2&((\theta Hu),s,u)^4\\[0.45em]
&(H_\vartheta(\theta Hu),s,u)^2,\ (\theta_\eta(\theta Hu),s,u)^2,\ (\theta_\eta(\theta Hu)^2,s,u)&(\theta_\eta su)^4,\ (\theta_\eta(\theta Hu),s,u)^3\\[0.45em]
&\cdot&(\theta_\eta^2(\theta Hu),s,u)^2
\end{array}
\]

% inherited: noether_paper02_ko_translation_001_20260804\targets\Noether_P02_Korean_T25_U180_UNCHECKED.tex
% sha256: 3A64DA40BC3724E9CF0D466BF07308C85DCDFF12050B1C10ED7E58BA9CEE2BE9
\noindent\emph{환원:}

A) 형식들 $((\vartheta^2\nu x)^3,s,u)^2$, $((\vartheta^2\nu x)^3,s,u)^3$은 분해된다:
\[
(\vartheta \nu x)(\widehat{\vartheta \nu su})(\widehat{\vartheta'\nu su})
=(\vartheta \nu x)s_\vartheta s_{\vartheta'}-(\vartheta \nu x)s_\nu s_\vartheta
=-2\theta(\vartheta \nu x)s_\nu s_\vartheta+(\vartheta \nu x)s_\vartheta^2.
\]
나머지 형식들은 환원자를 인수로 갖거나 분해 가능 형식들과 한 번 수축하여 생긴다.

% inherited: noether_paper02_ko_translation_001_20260804\targets\Noether_P02_Korean_T25_U181_UNCHECKED.tex
% sha256: A5CEF8580A274A1AB97E7D33E179B9167B1EEE2DAA15048A73E27DE63129DE16
B) 환원자는 $a_\nu^2s_\nu$와 $(\widehat{j\nu su})s_\nu$이다.

C) 1. 형식열 $(\theta Hu)^2s_\eta$: a) 환원자 $(aHu)^2(asu)s_\eta$; b) 형식열 $(\theta gu)^3g_\nu$와 $g_\nu(agu)^3(bgu)^2(abu)$의 이중 환원.

이로부터 더 높은 형식들 $(\theta_\eta su)^3$, $(H_\vartheta(\theta Hu),s,u)^3$, $(a_\nu\cdot b_{\nu_1}^2,s,u)^4$의 환원이 따른다[마지막 형식은 $(H_\vartheta su)^4$ 대신이다].

% inherited: noether_paper02_ko_translation_001_20260804\targets\Noether_P02_Korean_T25_U182_UNCHECKED.tex
% sha256: AB5D0CF763BDA46A5BCA67D247BDF83E7AE765C647E915042B43C96CB9946B9E
2. $((\vartheta\eta x),s,u)^4$: 환원자 $(a_\eta su)^4$.

3. $((\vartheta\eta x)^2,s,u)^3$: $(\widehat{\vartheta\eta}su)^2(Hsu)$에서 얻는 환원자 $s_\vartheta^2s_\nu$.
형식들 $(\vartheta\eta x)s_\eta$, $(\vartheta\eta x)^2s_\eta$, $((\vartheta\eta x)^2,s,u)^2$, $\theta_\eta s_\eta$는 분해 가능 형식 $a_\eta s_\eta$와 수축하여 분해된다.

% inherited: noether_paper02_ko_translation_001_20260804\targets\Noether_P02_Korean_T25_U183_UNCHECKED.tex
% sha256: 7567EC0F164D826902283D1239D902F3D8CFCAA14C1C4505DAF576A6092B148F
4. 형식열 $H_\vartheta(\theta Hu)s_\eta$, $\theta_\eta(\theta Hu)s_\eta$, $H_\vartheta(\vartheta\eta x)s_\eta$, $\theta_\eta(\vartheta\eta x)s_\eta$:
환원자 $a_\eta(aHu)s_\eta$와 $a_\eta^2s_\eta$.

형식열 $(\theta_\eta(\vartheta\eta x),s,u)^2$: 환원자 $a_\eta^2(Hsu)^2$ [단, $(\theta_\eta^2(\theta Hu),s,u)^2$은 예외이며, 이는 항 $a_\eta^2(bsu)(Hsu)$에서 수축하여 생긴다].

% inherited: noether_paper02_ko_translation_001_20260804\targets\Noether_P02_Korean_T25_U184_UNCHECKED.tex
% sha256: 90E720DBC3547091FF0F53FDD3788512A1751F89523EF27805EC910257A2C556
5. $(H_\vartheta(\vartheta\eta x),s,u)$, $(H_\vartheta(\vartheta\eta x),s,u)^2$: $a_\eta(aHu)s_\eta(abu)$와 $g_\vartheta(\theta gu)g_\nu$의 이중 환원.

형식열 $(H_\vartheta(\vartheta\eta x),s,u)^3$: 환원자 $((\vartheta \nu x)^2,s,u)^2s_\nu$.

% inherited: noether_paper02_ko_translation_001_20260804\targets\Noether_P02_Korean_T25_U185_UNCHECKED.tex
% sha256: EDCB0BE3314F23D47DEC44FE252DC8109295340452FDCD73C694C5DF0EC7D2B9
6. $(H_\vartheta(\theta Hu),s,u)$는 \S\ 3. 2), c)에 따라 분해 가능 형식 $(\theta_\nu(\vartheta \nu x),s,u)$과 한 번 수축하여 분해된다. 즉 더 낮은 분해 가능 형식 $(H_\vartheta su)^2$을 이중 환원한다:\srcfn{*)}{기호 $\equiv$는 더 낮은 차수의 형식들의 곱을 생략했음을 뜻한다.}
\[
\begin{aligned}
0&\equiv \theta_\nu(\vartheta\nu\nu_1)(\theta su)+\frac23\theta_\nu^2(\vartheta \nu_1 x)(\theta su)+\theta_\nu(\vartheta \nu x)(\theta su)s_{\nu_1}\\
&\equiv -\frac12H_\vartheta(\theta Hu)(\theta su)+\theta_\eta(\theta su)(Hsu)+\frac12H_\vartheta(\theta su)(Hsu)\\
&\equiv -\frac12H_\vartheta(\theta Hu)(\theta su)+\theta_\eta(\theta su)(Hsu)+\frac12[H_\vartheta]_{\widehat{su}^2}-\frac12\cdot\frac35H_\vartheta(\theta Hu)(\theta su)\\
&\equiv -\frac35H_\vartheta(\theta Hu)(\theta su).
\end{aligned}
\]

% inherited: noether_paper02_ko_translation_001_20260804\targets\Noether_P02_Korean_T25_U186_UNCHECKED.tex
% sha256: 4DD9F8F99DA5CB8005CDBAE67340CDD95786B59D19244487A1FBC434C430FB48
1과 5에 따라 이중 환원하면 다음 공식들이 생긴다:
\[
\begin{array}{@{}r@{\quad}l@{}}
&\displaystyle
\begin{aligned}
0={}&\theta_\eta(\theta su)(Hsu)^2+\frac14H_\vartheta(\theta Hu)(\theta su)^2-\frac34\theta_\eta(\theta Hu)(\theta su)(Hsu)\\
&-\frac54\theta_\eta(\theta Hu)^2(\theta su)-(asu)^3\cdot b_\eta(bHu)^2-a_\nu(asu)^3\cdot b_{\nu_1}^2,
\end{aligned}\\[0.55em]
(31.)&\displaystyle 0=H_\vartheta(\theta Hu)(\theta su)^3-7\theta_\eta(\theta Hu)(\theta su)^2(Hsu),\\[0.55em]
&\displaystyle 0=a_\nu(asu)^2b_{\nu_1}^{2}(bsu)^2-\theta_\eta(\theta su)^2(Hsu)^2+6\theta_\eta(\theta Hu)(\theta su)^2(Hsu),\\[0.55em]
&\displaystyle 0\equiv (H_\vartheta(\vartheta\eta x),s,u),\qquad
0\equiv H_\vartheta(\widehat{\vartheta\eta su})(Hsu)-4\theta_\eta^2(Hsu).
\end{array}
\]

% inherited: noether_paper02_ko_translation_001_20260804\targets\Noether_P02_Korean_T25_U187_UNCHECKED.tex
% sha256: ABE6F52DEE0BB4A1C77E2FE6B98119417630DD00BAE09FC90322C8DAB6480001
\noindent\emph{결론:}

1) $(\theta Hu)^2s_\eta$, $H_\vartheta(\theta Hu)s_\eta$, $\theta_\eta(\theta Hu)s_\eta$, $H_\vartheta(\vartheta\eta x)s_\eta$, $\theta_\eta(\vartheta\eta x)s_\eta$, $(H_\vartheta(\vartheta\eta x),s,u)$, $(\theta_\eta(\vartheta\eta x),s,u)^2$은 환원자이다.

2) 계수차수 6의 형식들 $(5,\lambda)$ 등은 $s^2$과의 수축에 들어가지 않는다.

% inherited: noether_paper02_ko_translation_001_20260804\targets\Noether_P02_Korean_T26_U191_UNCHECKED.tex
% sha256: 9E1AD3438B540B94202CA21269A9B98220F8D47A3EB13EC9FD02F1B3130E5C96
\begin{center}\textbf{\S\ 24.}\quad\emph{계수차수 11의 형식들(계 $s_{III}$).}\end{center}
$s$와 계수차수 7의 형식들(계 III)의 수축(\S\ 12).

\noindent\emph{환원 불가능 형식들:}
\[
\text{A)}\quad
\begin{array}{c|c|c|c|c}
\cdot&\cdot&\cdot&((KHu)^2su)^2&\bigstar\\[0.35em]
\cdot&\cdot&(K_\eta su)^2&\cdot&\\[0.35em]
\cdot&\cdot&\cdot&\cdot&\\[0.35em]
((k\eta x)^3,s,u)&\cdot&\cdot&\cdot&\\[0.35em]
\cdot&(K_\eta(k\eta x)^3,s,u)&\cdot&\cdot&
\end{array}
\]
\[
\begin{array}{c@{\quad}|c|c}
\bigstar&((KHu)^2,s,u)^3&((KHu)^2,s,u)^4\\[0.45em]
&\cdot&(K_\eta(KHu)^2,s,u)^3
\end{array}
\]
\[
\text{B)}\quad ((j\eta x),s,u)^2,\ (H_jsu)\ \big|\ ((j\eta x),s,u)^3,
\qquad \text{C)}\quad (\D_\eta su),
\]
\[
\text{D)}\quad
\begin{array}{c|c|c}
\cdot&((aLu)^2,s,u)^2,\ ((aLu)^3,s,u)&((aLu)^2,s,u)^3\\[0.5em]
(aLu)^2s_\eta,\ (a_lsu)^2&(a_lsu)^3&(a_l(aLu)^2,s,u)^2
\end{array}
\]

% inherited: noether_paper02_ko_translation_001_20260804\targets\Noether_P02_Korean_T26_U192_UNCHECKED.tex
% sha256: E3076E2A3FEFF85B44314E375119484D3677F10AE38A258EF19AC18D79551835
\noindent\emph{환원:}

A) 기호 $\theta\cdot H$로 나타낸 대응 형식들과 한 번 수축하여 환원되거나 분해된다.

$(K_\eta(KHu)^2,s,u)$와 $(K_\eta(KHu)^2,s,u)^2$: $K_\nu^2(Ksu)$, $K_\nu^2(Ksu)^2$과 수축하여 \S\ 3. 2), a)에 따라 분해된다.

$(K_\eta(KHu),s,u)^4\equiv(a_\nu^2\cdot\theta_{\nu_1},s,u)^4$: 분해 가능 형식 $K_\nu^2(Ksu)^3$으로부터 \S\ 3. 2), b)에 따라 분해된다.

% inherited: noether_paper02_ko_translation_001_20260804\targets\Noether_P02_Korean_T26_U193_UNCHECKED.tex
% sha256: 2AE10D4F36EB642FF0F8ABBDCDDB59884EE0222315CB2D181F7B2BD5F55784CD
B) 형식열 $(j\eta x)s_\eta$: 환원자 $(j\nu x)^2s_\nu$.

형식 $H_j(\widehat{j\eta su})(Hsu)$: 환원자 $(j\nu x)(\widehat{j\nu su})^2$.

형식 $H_j(j\eta x)(Hsu)$: 환원자 $(j\nu x)^2s_\nu$를 사용하여 분해 가능 형식 $((j\eta x)^2,s,u)^2$을 환원함으로써 분해된다(공식 (18.)a):
\[
0=-2(j\nu x)^2s_\nu s_{\nu_1}=[(j\eta x)^2]_{\widehat{su}^2}+H_j(j\eta x)(Hsu).
\]

% inherited: noether_paper02_ko_translation_001_20260804\targets\Noether_P02_Korean_T26_U194_UNCHECKED.tex
% sha256: D5C602A8C1F697DB2CC1FC789A25594B541C76321D17BC7AE2ED1F2EBC257161
C) 환원자 $\D_\nu s_\nu$와 $\D_\nu(\D su)^2$.

D) 기호 $f\cdot H$로 나타낸 대응 형식들과 한 번 수축하여 환원되거나 분해된다.

% inherited: noether_paper02_ko_translation_001_20260804\targets\Noether_P02_Korean_T26_U195_UNCHECKED.tex
% sha256: 45D0ADE9EA4A60CE6E7EA36F95D3E0B663A3E73D9020ABFB3068491ACE25AD9E
E) (30.)c에서 분해 가능 형식열 $(a_\sigma su)^2$의 환원을 얻는다:
\[
\begin{aligned}
0&=a_\eta(Hsu)^2(as\widehat{\nu}x)^2
=a_\eta(Hsu)^2(a\widehat{\nu}\eta u)^2+2a_\eta(a\widehat{\nu}\eta u)(\widehat{\nu}\eta su)(Hsu)^2
=\frac13a_\sigma(asu)^2,\\
0&=a_\eta a_\nu(Hsu)^2(as\widehat{\nu}x)(asu)
=a_\eta(a\widehat{\nu}\eta u)^2(Hsu)^2(asu)+a_\eta(a\widehat{\nu}\eta u)(\widehat{\nu}\eta su)(Hsu)^2(asu)
=\frac13a_\sigma(asu)^3.
\end{aligned}
\]

% inherited: noether_paper02_ko_translation_001_20260804\targets\Noether_P02_Korean_T26_U196_UNCHECKED.tex
% sha256: 68ABBE71AC67E644641E4AD2212FB9A7EC9F37765FE3DF8C85EE4079C7B7F4D0
\noindent\emph{결론:} $s^2$은 계수차수 7의 형식들과 수축되지 않는다.

% inherited: noether_paper02_ko_translation_001_20260804\targets\Noether_P02_Korean_T27_U197_UNCHECKED.tex
% sha256: 1586327601BCA3CF79A65E252DCBF7563BC2D374EDE6F90991186F65336F84D0
\begin{center}\textbf{\S\ 25.}\quad\emph{계수차수 12, 13, 14, 15의 형식들(계 $s_{III}$).}\end{center}
$s$와 계수차수 8, 9, 10, 11의 형식들(계 III)의 수축(\S\S\ 13--16).

\noindent\emph{환원 불가능 형식들:}

% inherited: noether_paper02_ko_translation_001_20260804\targets\Noether_P02_Korean_T27_U198_UNCHECKED.tex
% sha256: FBA951FF393E28BFE565013D02C5FFDB65A89CB9551D8619E5D9485AF4A9F162
\noindent 계수차수 12.
\[
\begin{adjustbox}{max width=0.98\linewidth}
$\displaystyle
\begin{array}{c@{\qquad}c@{\qquad}c}
\text{A)}&
\begin{array}{c|c}
((\vartheta^2\eta x)^4,s,u)&\cdot\\[0.35em]
\cdot&\theta_\eta(\vartheta^2\eta x)^3s_\vartheta
\end{array}
&
\text{B)}\quad ((aHu)H_j,s,u)^2\;\big|\;((aHu)H_j,s,u)^3
\\[1.0em]
\text{C)}&
\begin{array}{c|c|c}
\cdot&((\theta Lu)^2,s,u)^2,\ ((\theta Lu)^3,s,u)&((\theta Lu)^2,s,u)^3\\[0.35em]
(\theta_lsu)^2&\cdot&(\theta_l(\theta Lu)^2,s,u)^2
\end{array}
&
\end{array}$
\end{adjustbox}
\]

% inherited: noether_paper02_ko_translation_001_20260804\targets\P02_T27_U199_r2.tex
% sha256: B04B5220D204C0705979CB65D317947BCD6761416126E96907B358DBA18C3E8A
\noindent\emph{환원:} 환원자 또는 분해 가능 형식과 직접 수축하여 생기는 환원은 앞으로 더 이상 언급하지 않는다.

B) $((aHu)H_j,s,u)$와 $(a_\eta(aHu)H_j,s,u)$는 야코비 행렬식에 대한 전이연산으로서 분해된다.

% inherited: noether_paper02_ko_translation_001_20260804\targets\Noether_P02_Korean_T27_U200_UNCHECKED.tex
% sha256: 3E14F6E613092197BA51CC87DD8746D56810C455A33A81A1D551F4FCC5436CCA
$(a_\eta(aHu)H_j,s,u)^2$의 분해: 분해 가능 형식 $[\theta_\nu\cdot\theta_{\nu_1}^{\prime 3}]_{\widehat{su^3}}$의 이중 환원(\S\ 13. C)
\begin{align*}
0&\equiv[\theta_{\nu_1}^{\prime 3}\cdot\theta_\nu]_{\widehat{su^3}}
\equiv-\frac34\theta_\nu(\theta su)^2(\theta\theta'u)\theta_{\nu_1}^{\prime 3}\\
&\equiv-\frac98(\theta\theta'\eta su)(\theta\theta'u)(\theta Hu)(\theta'Hu)\theta_\eta'(\theta su)
\equiv-\frac38a_\eta(aHu)H_j(asu)^2.
\end{align*}

% inherited: noether_paper02_ko_translation_001_20260804\targets\Noether_P02_Korean_T27_U201_UNCHECKED.tex
% sha256: 440EF289193636FA8E036DD51E8EE5E111582492704704154A639217FAAC999A
C) 공식 (31.)로부터 분해 가능 형식
\[
 [L_\vartheta(\theta Lu)]_{\widehat{su^3}}\quad\hbox{및}\quad [\theta_l(\theta Lu)]_{\widehat{su^3}}
\]
의 환원이 얻어진다. 즉 곱 $[\theta_\nu^2\cdot H]_{\widehat{su^3}}$와 $[a_\nu^2\cdot(bHu)^2]_{\widehat{su^3}}$에 대해 다음을 얻는다:
\[
\begin{array}{@{}r@{\quad}l@{}}
&\displaystyle 0\equiv\theta_\nu^2(\theta su)^2(Hsu),\qquad
0\equiv [L_\vartheta(\theta Lu)]_{\widehat{su^4}}-7[\theta_l(\theta Lu)]_{\widehat{su^4}},\\[0.55em]
&\displaystyle 0\equiv a_\nu^2(asu)^2(bHu)^2(bsu)+6\theta_l(\theta Lu)^2(\theta su)(Lsu),\\[0.55em]
(32.)&\displaystyle 0\equiv\theta_\nu^3(\theta su)(Hsu)^2
\quad\hbox{은 }0\equiv\theta_l^2(\theta Lu)(\theta su)(Lsu)^2\quad\hbox{에서 얻는다}\\[-0.05em]
&\displaystyle\hfill(\hbox{환원자 }a_\eta^2(Hsu)^2).
\end{array}
\]

% inherited: noether_paper02_ko_translation_001_20260804\targets\Noether_P02_Korean_T27_U202_UNCHECKED.tex
% sha256: 9E922A14F005C092137CF4E83616E05E9E96F6262EC9B9D3504FF7C3A9FA699C
\noindent\emph{환원 불가능 형식들.}

\noindent 계수차수 13.
\[
\begin{array}{rl}
\text{A)}&((KLu)^3,s,u)^2;\ ((KLu)^3,s,u)^3,\\[0.25em]
\text{B)}&(L_jsu)^2,\\[0.25em]
\text{C)}&((aH^2u)^3,s,u)^2.
\end{array}
\]

% inherited: noether_paper02_ko_translation_001_20260804\targets\Noether_P02_Korean_T27_U203_UNCHECKED.tex
% sha256: 3F117865DD27EBB3EC386473A23E6BB718D5CC9B998BACE08360946B75F16F3F
\noindent 계수차수 14.
\[
\begin{array}{rl}
\text{A)}&((aLu)^2L_j,s,u)^2,\\[0.25em]
\text{B)}&((\theta H^2u)^3,s,u)^2.
\end{array}
\]
\noindent 계수차수 15.
\[
 ((KH^2u)^4,s,u)^2.
\]

% inherited: noether_paper02_ko_translation_001_20260804\targets\Noether_P02_Korean_T27_U204_UNCHECKED.tex
% sha256: EE9BDD3E359CCF31E9F124DA6F222F0193A46A650F6525D3928FAABFC4BC30E4
\noindent\emph{환원:} $(K_l(KLu)^3,s,u)^2$, $(H_\vartheta(\theta H^2u)^3,s,u)$, $((K,H\cdot L,u)^5,s,u)$는 \S\ 3. 2), c)에 따라 분해된다. 즉 분해 가능 형식 $(K_l(KLu)^2,s,u)^3$, $(H_\vartheta(\theta H'u)^2,s,u)^2$, $((K,H\cdot L,u)^4,s,u)^2$을 동시에 고찰한다.

\noindent\emph{결론:} $s^2$은 계 III의 개별 형식들과 수축되지 않는다.

% inherited: noether_paper02_ko_translation_001_20260804\targets\Noether_P02_Korean_T28_U211_UNCHECKED.tex
% sha256: 2D15A93A3FD1E99E47CE651D0BD0C9D7D93FA096815F9B725471E154CAA09959
\begin{center}\textbf{\S\ 26.}\quad\emph{계 $s_{IV}$.}\end{center}
\noindent 1) \emph{$s$와 계 I$\cdot$III의 수축.}

\noindent\emph{환원:} 앞 절들의 환원들 $(a_\nu^2s_\nu,\ (aHu)^2(asu)s_\eta$ 등)에 따라 형식 $(0,\lambda)$, $(1,\lambda)$, $(2,\lambda)$는 수축 I과 III을 동시에 허용하지 않는다. 따라서 \S\ 18에 따라 계 $s_{IV}$를 만들 때 $f,\D,N$은 곱 안에서 수축되지 않는다.

$j$는 수축되지 않는다. 같은 환원들에 따라 $j$에 반변인 다음 수축들
\[
(K_\nu(k\nu x)^3,s,u),\qquad ((\vartheta^2\eta x)^4,s,u)\quad\hbox{등}
\]
에 $j$와 동변인 다음 수축들이 대응하기 때문이다:
\[
K_\nu(k\nu x)^2s_k,
\qquad (\vartheta^2\eta x)^3s_\vartheta\quad\hbox{등}.
\]

% inherited: noether_paper02_ko_translation_001_20260804\targets\Noether_P02_Korean_T28_U212_UNCHECKED.tex
% sha256: A33075585EF068E579626FDDA3FEBF753298B1CF4788AEBA2D06180364E168D0
\noindent 2) \emph{$s$와 계 II$\cdot$III의 수축},
즉 $s$와 $H\cdot(1,\lambda)$, $H\cdot(2,\lambda)$, $H\cdot(3,\lambda)$의 수축.

\noindent\emph{환원 불가능 형식들:}
\[
\begin{array}{rl}
\text{A)}&(a_\nu\cdot H,s,u)^4,
((aHu)^2\cdot H,s,u)^3,
((aHu)^2\cdot H,s,u)^4,\\[0.25em]
&((aLu)^2\cdot H,s,u)^4,
((aLu)^3\cdot H,s,u)^3,
((aH^2u)^3\cdot H,s,u)^4,\\[0.45em]
\text{B)}&(a_\nu^2\cdot H,s,u)^3,
(a_\eta(aHu)^2\cdot H,s,u)^3,
(a_l(aLu)^2\cdot H,s,u)^4,\\[0.45em]
\text{C)}&(\theta_\nu\cdot H,s,u)^4,
((\theta Hu)^2\cdot H,s,u)^3,
((\theta Hu)^2\cdot H,s,u)^4,\\[0.25em]
&((\theta Lu)^2\cdot H,s,u)^4,
((\theta Lu)^3\cdot H,s,u)^3,
((\theta H^2u)^3\cdot H,s,u)^4,\\[0.45em]
\text{D)}&(\theta_\eta(\theta Hu)^2\cdot H,s,u)^3,
(\theta_l(\theta Lu)^2\cdot H,s,u)^4,\\[0.45em]
\text{E)}&((KLu)^3\cdot H,s,u)^4,
((KH^2u)^4\cdot H,s,u)^4,\\[0.45em]
\text{F)}&((j\nu x)^2\cdot H,s,u)^3,
((j\nu x)^2\cdot H,s,u)^4,
((j\eta x)\cdot H,s,u)^4,\\[0.25em]
&(H_j\cdot H,s,u)^3,
(L_j\cdot H,s,u)^4.
\end{array}
\]

% inherited: noether_paper02_ko_translation_001_20260804\targets\Noether_P02_Korean_T28_U213_UNCHECKED.tex
% sha256: 02ACA5823AFC9832327D262272834F98542B0E098AD25907098895CD642F5585
\noindent\emph{환원:} $\nu$는 $j$와 마찬가지로 수축되지 않는다.

B)와 D). $(a_\nu^2\cdot H,s,u)^4$과 $(\theta_\nu^2\cdot H,s,u)^4$은 $(gHu)^2=-\frac13s\cdot\sigma$를 고려하면 공식 (27.)c)와 (29.)c)에서 분해된다:
\[
(a_\nu^2\cdot H,s,u)^4=\frac13(asu)^4\cdot\sigma,
\qquad
(\theta_\nu^2\cdot H,s,u)^4=-\frac16(\theta su)^4\cdot\sigma;
\]
그로부터 $(a_\eta(aHu)\cdot H,s,u)^4$, $(\theta_\eta(\theta Hu)\cdot H,s,u)^4$의 분해가 따른다.

% inherited: noether_paper02_ko_translation_001_20260804\targets\Noether_P02_Korean_T28_U214_UNCHECKED.tex
% sha256: 470C23713BB4C4890A419D19B186A9EBEE30B18D29C6C4FD8A1633A333028D1C
$(\theta_\nu^2\cdot H,s,u)^3$, $(\theta_\nu^3\cdot H,s,u)^3$, 그리고 그로부터 $(\theta_\eta^2(\theta Hu)\cdot H,s,u)^4$은 \S\ 25의 계수차수 12의 형식 C)에 따라 분해된다.

% inherited: noether_paper02_ko_translation_001_20260804\targets\Noether_P02_Korean_T29_U215_UNCHECKED.tex
% sha256: D8CB6DC9CEC436B73271B15EFD46D094EB796DE3C882F53CBF3128AEB557C5F8
\noindent 3) \emph{$s$와 계 III$\cdot$III의 수축}, 즉 $s$와 계 III의 다음 형식들의 수축:
\[
(3,\lambda)(3,\lambda),\quad (3,\lambda)(2,\lambda),\quad (3,\lambda)(1,\lambda),\quad
(2,\lambda)(2,\lambda),\quad (2,\lambda)(1,\lambda),\quad
(1,\lambda)(1,\lambda).
\]
\noindent\emph{환원 불가능 형식들:}
\[
\begin{array}{rl}
\text{A)}&(a_\nu^2\cdot a_\nu^2,s,u)^3,
(a_\nu^2\cdot a_\nu^2,s,u)^4,
(a_\nu^2\cdot a_\eta(aHu)^2,s,u)^3,\\[0.45em]
\text{B)}&(a_\nu\cdot H_j,s,u)^4,
((aHu)^2\cdot H_j,s,u)^3,
((aLu)^2\cdot H_j,s,u)^4,\\[0.25em]
&((aLu)^3\cdot H_j,s,u)^2,
((aH^2u)^3\cdot H_j,s,u)^3.
\end{array}
\]

% inherited: noether_paper02_ko_translation_001_20260804\targets\Noether_P02_Korean_T29_U216_UNCHECKED.tex
% sha256: 0DB1EC5BF7666430331914EB8B9C66A9E8DE1912D5CC1F86C9DCB399A9181E5C
\noindent\emph{환원:}

기호 $\nu$와 $f$에서 계수차수가 같을 때, 곱 II$\cdot$III은 곱 III$\cdot$III보다 더 높은 것으로 간주한다.

환원들은 일부는 곱들 사이의 관계식(시지지)에서 생기고, 일부는 곱들을 대응하는 분해 가능 형식들로 바꾸고 이 형식들과 $s$의 수축을 환원함으로써 생긴다. (반공변식을 포함하는 곱들은 곱으로 넘어가므로 항상 생략한다.)

A) $f$에 관하여 이차이고, 기호 $\nu$는 임의의 계수차수이다. \S\ 11의 공식 (12.)와 유사한 계산에 따라 다음이 성립한다:
\begin{align*}
1)\quad 0&=a_\nu\cdot b_{\nu_1}+\frac12f\cdot(aHu)^2-\frac14\theta\cdot H+\frac12u_xH_\vartheta-\frac14u_x^2\cdot H_\vartheta^2,\\
2)\quad 0&=a_\nu\cdot b_{\nu_1}^2+f\cdot a_\eta(aHu)^2-\theta_\eta-\frac12H_\vartheta,\\
3)\quad 0&=a_\nu^2\cdot b_{\nu_1}^2-H_\vartheta^2+2\theta_\eta^2.
\end{align*}

% inherited: noether_paper02_ko_translation_001_20260804\targets\Noether_P02_Korean_T29_U217_UNCHECKED.tex
% sha256: 52801A4DA6412BD2E4173C053EB4BB14076A54CBBDDD6C1A5C888E772100B00F
그로부터 다음이 따른다:
\begin{align*}
4)\quad 0&=a_\nu^2\cdot b_{\nu_1}^2(j\nu_1x)\quad(3)\text{에서}),\\
5)\quad L_\vartheta(\theta Lu)&=-a_\nu\cdot b_\eta(bHu)^2+\frac34a_\nu^2\cdot(bHu)^2+\frac34\theta_\nu^2\cdot H\quad(2)\text{에서}),\\
6)\quad L_\vartheta(\theta Lu)&=-6a_\nu\cdot b_\eta(bHu)^2+2a_\nu^2\cdot(bHu)^2-\frac12\theta_\nu^2\cdot H\quad(1)\text{에서}),\\
7)\quad 0&=a_\nu\cdot(bHu)^2+f\cdot(aLu)^3+\theta_\nu\cdot H\quad(1)\text{에서}),
\end{align*}
\[
\left.
\begin{aligned}
8)\quad 0&=a_\nu\cdot(bLu)^3+\frac34(\theta Hu)^2\cdot H\\
9)\quad 0&=(aHu)^2\cdot(bHu)^2+2(\theta Hu)^2\cdot H
\end{aligned}
\right\}\quad(7)\text{에서}),
\]
\begin{align*}
10)\quad 0&=a_\eta(aHu)^2\cdot b_\eta(bHu)^2\quad(5)\text{와 }6)\text{에서}),\\
11)\quad 0&\equiv [a_\nu^2\cdot b_\eta(bHu)]_{\widehat{su^4}}\quad(\text{공식 (30.)b에서}).
\end{align*}

% inherited: noether_paper02_ko_translation_001_20260804\targets\Noether_P02_Korean_T29_U218_UNCHECKED.tex
% sha256: EEBF7062875818798B077AF1EE0842B23D80A33C3AF687D7728CE15FFBFF80CB
다음 형식들의 환원
\[
(H_\vartheta su)^4,
\qquad (L_\vartheta(\theta Lu),s,u)^3,
\qquad (L_\vartheta(\theta Lu),s,u)^4
\]
(공식 (31.)과 (32.))은 공식 1)부터 11)까지와 함께, 기호 $f$에 관하여 이차이고 $\nu$에 관하여 임의인 곱들과 $s$의 모든 수축을 환원한다.

% inherited: noether_paper02_ko_translation_001_20260804\targets\Noether_P02_Korean_T30_U219_UNCHECKED.tex
% sha256: 20CF7671D6F85071BDD81893260B028549915C48E34375DD2FA962984E30E52D
B) 기호 $f,\theta,j$ 중 각각 두 개의 곱이며, $\nu$는 임의의 계수차수이다. 공식 (17.), (18.), (20.)과 직접 계산에 따라 다음 관계식들이 성립한다:
\[
\left.
\begin{aligned}
1)\quad 0&=a_\nu\cdot\theta_{\nu_1}+\frac14\theta\cdot(aHu)^2+\frac14f\cdot(\theta Hu)^2,\\
2)\quad 0&=\theta_\nu\cdot\theta_{\nu_1}+\frac12\theta\cdot(\theta Hu)^2
\end{aligned}
\right\}\quad\text{(A 1)과 유사한 직접 계산).}
\]
그로부터 다음이 따른다:
\[
\left.
\begin{aligned}
3)\quad 0&=3a_\nu\cdot(\theta Hu)^2+\theta\cdot(aLu)^3+2f\cdot(\theta Lu)^3\\
4)\quad 0&=3\theta_\nu\cdot(aHu)^2+f\cdot(\theta Lu)^3+2\theta\cdot(aLu)^3
\end{aligned}
\right\}\quad\text{(1)과 A 6)에서),}
\]

% inherited: noether_paper02_ko_translation_001_20260804\targets\Noether_P02_Korean_T30_U220_UNCHECKED.tex
% sha256: 731A65C24FC10C164B15E21E7EC926D3CC43EB8BCC950B5D3D864B6EAC6386C6
\begin{align*}
5)\quad 0&=a_\nu\cdot(\theta Lu)^3,\quad 0=\theta_\nu\cdot(aLu)^3,\quad 0=(aHu)^2\cdot(\theta Hu)^2
\quad\text{(3), 4), A 8)에서),}\\
6)\quad 0&=a_\nu\cdot\theta_\nu^2+\theta_\nu\cdot a_\nu^2+f\cdot\theta_\eta(\theta Hu)^2+\theta\cdot a_\eta(aHu)^2\quad(\text{공식 (17.)}),\\
7)\quad 0&=\theta_\nu\cdot\theta_\nu^2+\theta\cdot\theta_\eta(\theta Hu)^2+\frac16f\cdot H_j\\
&\qquad [6)\text{을 }\theta_\nu^2(\theta\theta'\nu_1x)\text{에서 유도한 것과 유사한 유도}],\\
8)\quad 0&=a_\nu\cdot(j\nu x)^2+f\cdot H_j\quad(6)\text{에서}),\\
9)\quad 0&=(aHu)^2\cdot(j\nu x)^2-2a_\nu\cdot H_j\quad(8)\text{에서}),\\
10)\quad 0&=\theta_\nu\cdot(j\nu x)^2,\quad 0=\theta\cdot H_j\quad(7)\text{과 }8)\text{에서}),
\end{align*}
\[
\left.
\begin{aligned}
11)\quad 0&=a_\nu\cdot\theta_\nu^3-\frac34(j\eta x)^2,\\
12)\quad 0&=a_\nu^2\cdot\theta_\nu^2-\frac13(j\eta x)^2-\frac13(\D Hu)^2
\end{aligned}
\right\}\quad\text{(공식 (18.)),}
\]
\begin{align*}
13)\quad 0&=a_\nu^2\cdot\theta_\nu^3-\frac12a_\sigma\quad(\text{공식 (20.)}),\\
14)\quad 0&=\theta_\nu\cdot\theta_\nu^3,
&0&=a_\eta H_j\quad(\S\ 13.\ C)).
\end{align*}

% inherited: noether_paper02_ko_translation_001_20260804\targets\Noether_P02_Korean_T30_U221_UNCHECKED.tex
% sha256: A219A9DB8A450D0262B874C17D79A43C48922E53E84B522364A3D2E4218B4044
공식들은 곱들이 편극 가능해질 때까지 전개하였다. 즉, 다음 이유로 수축을 통해 단 \emph{하나}의 새로운 곱만 생기는 때까지이다:
\begin{center}
\begin{tabular}{rl}
a)&곱 자체와의 수축이 환원 가능해지고,\\
b)&수축으로 생기는 나머지 곱들이 환원 가능해지거나
\end{tabular}
\end{center}
반공변식으로 이어진다(\S\ 3. 2), a)와 b) 참조).

공식 1)부터 5)까지는 $a_\nu\cdot\theta_\nu$에서 수축으로 생기는 곱들과 $s$의 모든 수축을 환원하고,

공식 6)부터 10)까지는 $a_\nu\cdot\theta_\nu^2$에서 수축으로 생기는 곱들을 환원한다.

\S\ 24 A)에 따르면, 공식 6)부터 10)까지를 고려할 때 $a_\nu^2\cdot\theta_\nu$에서 생기는 곱들과 $s$의 수축들은 분해된다.

% inherited: noether_paper02_ko_translation_001_20260804\targets\Noether_P02_Korean_T30_U222_UNCHECKED.tex
% sha256: 6ACC6BEE9E52602AEF747470C12F3B665A588CEB9AA74D438B6B01DE147C0256
다음을 얻는다:
\begin{align*}
0&\equiv a_\nu^2(asu)^2\theta_{\nu_1}(\theta su)^2,
&0&\equiv a_\nu^2(asu)\theta_{\nu_1}(\theta su)^3-K_\eta(KHu)^2(Ksu)^3,\\
0&\equiv a_\nu^2(asu)^2\theta_{\nu_1}(\theta su),
&0&\equiv a_\nu^2(asu)\theta_{\nu_1}(\theta su)^2,\\
0&\equiv a_\nu^2(asu)\theta_{\nu_1}(\theta su).
\end{align*}
환원자들
\[
((j\eta x)^2,s,u)^3
\ [\hbox{$(\vartheta\eta\widehat{su})^2(Hsu)$에서. \S\ 23.\ C\ 3)}],
\]
\[
(H_j(j\eta x),s,u)^2\ [\S\ 24,\ B],
\quad ((\D Hu)^2,s,u)^2\ [\S\ 24,\ C],
\]
\[
(a_\sigma su)^2\ [\S\ 24,\ E]
\]
은 공식 11)부터 14)까지와 함께 $a_\nu\cdot\theta_{\nu_1}^3$, $a_\nu^2\cdot\theta_{\nu_1}^2$, $a_\nu^2\cdot\theta_{\nu_1}^3$에서 생기는 모든 곱을 환원한다.

% inherited: noether_paper02_ko_translation_001_20260804\targets\Noether_P02_Korean_T31_U223_UNCHECKED.tex
% sha256: 6D80D7796CE2265D0754DD93645F235162A6B0E64BC5CCF2FA8A94E08F25EF6D
지금까지의 계산은 다음 \emph{결론}으로 요약된다:

상대 완전계 $\modu{(\R,t)}$는 다음 331개 형식과 부분계로 이루어지며, 이어서 이를 표로 배열하여 제시한다:
\[
\left.
\begin{array}{lr}
u_x&1\ \text{개 형식}\\
\text{계 I}&6\ \text{개 형식}\\
\text{계 II}&5\ \text{개 형식}\\
\text{계 III}&117\ \text{개 형식}
\end{array}
\right\}
\quad
\begin{gathered}
(s,\varrho,t)\text{를 법으로 하는 상대 완전계}\\
[\text{표 I 참조}].
\end{gathered}
\]

% inherited: noether_paper02_ko_translation_001_20260804\targets\Noether_P02_Korean_T31_U224_UNCHECKED.tex
% sha256: 3CDA607A9148DF06261B46ACCCB1C9C6227F13E993A5C504E0AAEFCFB94CFAD9
\[
\left.
\begin{array}{lr}
s&1\ \text{개 형식}\\
\text{계 }s_I&35\ \text{개 형식}\\
\text{계 }s_{II}&9\ \text{개 형식}\\
\text{계 }s_{III}&125\ \text{개 형식}\\
\text{계 }s_{IV}&32\ \text{개 형식}
\end{array}
\right\}
\quad [\text{표 II 참조}].
\]

% inherited: noether_paper02_ko_translation_001_20260804\targets\Noether_P02_Korean_T32_U225_UNCHECKED.tex
% sha256: 183E52729DDCFCC9AF95D25349126822DB69869A7C3E0AD64D33B3E896777364
\clearpage
\begin{landscape}
\thispagestyle{plain}
\begin{center}\bfseries 표 I (90쪽 참조).\end{center}
\vspace{-0.6em}
\begingroup
\setlength{\tabcolsep}{2pt}
\renewcommand{\arraystretch}{1.35}
\tiny
\begin{center}
\begin{adjustbox}{max width=0.98\linewidth,max totalheight=0.72\textheight,center}
\begin{tabular}{|c|c|c|c|c|c|c|c|c|c|c|c|c|c|c|c|c|c|c|c|c|c|c|}
\hline
{} & 0 & 1 & 2 & 3 & 4 & 5 & 6 & 7 & 8 & 9 & 10 & 11 & 12 & 13 & 14 & 15 & 16 & 17 & 18 & 19 & 20 & 21 \\ \hline
0 & \mcell{i} &  &  &  & \mcell{f} &  & \mcell{\A} &  & \mcell{K_{\nu}(k\nu x)^3} &  & \mcell{K_{\eta}(k\eta x)^3} &  &  &  & \mcell{(\vartheta^3\eta x)^4} & \mcell{H_k(k\cdot\vartheta,\eta x)^4} &  & \mcell{\theta_l(\vartheta\cdot k,lx)^5} &  &  &  & \mcell{(\vartheta^2lx)^6} \\ \hline
1 &  & \mcell{u_x} &  &  &  & \mcell{\theta_\nu(\vartheta\nu x)^2} &  & \mcell{\theta_\eta(\vartheta\eta x)^2} & \mcell{N} & \mcell{(k\nu x)^3} & \mcell{\theta_\nu(\vartheta^2\nu x)^3} & \mcell{(k\eta x)^3} & \mcell{H_{\vartheta}(\vartheta^2\eta x)^3\\ \theta_\eta(\vartheta^2\eta x)^3} &  & \mcell{\theta_l(\vartheta^2lx)^4} &  & \mcell{(\vartheta k,\eta,x)^4} &  & \mcell{(\vartheta k,l,x)^5} &  &  &  \\ \hline
2 &  &  & \mcell{a_\nu^2} &  & \mcell{\theta\\ a_\eta^2} &  & \mcell{(\vartheta\nu x)^2} & \mcell{K_\nu(k\nu x)^2} & \mcell{(\vartheta\eta x)^2} & \mcell{H_k(k\eta x)^2\\K_\eta(k\eta x)^2} &  & \mcell{(\vartheta^2\nu x)^3\\ K_l(klx)^3} &  & \mcell{(\vartheta^2\eta x)^3} &  & \mcell{(\vartheta^2lx)^4} &  &  &  &  &  &  \\ \hline
3 &  & \mcell{\theta_\nu^3} &  & \mcell{a_\nu} & \mcell{\theta_\nu(\vartheta\nu x)} & \mcell{\A_\nu\\a_\eta} & \mcell{K\\H_{\vartheta}(\vartheta\eta x)\\\theta_\eta(\vartheta\eta x)} & \mcell{\A_\eta} & \mcell{(k\nu x)^2\\\theta_l(\vartheta lx)^2} &  & \mcell{(k\eta x)^2} &  & \mcell{(klx)^3} &  &  &  &  &  &  &  &  &  \\ \hline
4 & \mcell{\nu} &  & \mcell{H\\\theta_\nu^2} & \mcell{(j\nu x)^3\\a_\eta(aHu)} & \mcell{\theta_\eta^2} & \mcell{(\vartheta\nu x)\\a_l^2} &  & \mcell{N_\nu\\(\vartheta\eta x)} &  & \mcell{H_\eta\\N_\eta\\(\vartheta lx)^2} &  &  &  &  &  &  &  &  &  &  &  &  \\ \hline
5 &  & \mcell{a_\eta(aHu)^2} & \mcell{\theta_\eta^2(\theta Hu)} & \mcell{\theta_\nu} & \mcell{\frac{K_\nu^2}{(aHu)}} & \mcell{\theta_\eta} & \mcell{K_\eta^2\\(\A Hu)\\a_l} &  & \mcell{\A_l} &  &  &  &  &  &  &  &  &  &  &  &  &  \\ \hline
6 & \mcell{j\\\sigma} &  & \mcell{(j\nu x)^2\\(aHu)^2} & \mcell{L\\a_\nu^2(j\nu x)\\H_{\vartheta}(\theta Hu)\\\theta_\eta(\theta Hu)} & \mcell{\frac{K_\nu}{\theta_l^2}} &  &  & \mcell{K_\eta} &  &  &  &  &  &  &  &  &  &  &  &  &  &  \\ \hline
7 &  & \mcell{\theta_\eta(\theta Hu)^2} & \mcell{H_j(j\eta x)\\a_l(aLu)^2} &  & \mcell{a_\nu(j\nu x)\\(\theta Hu)} &  & \mcell{\A_\nu(j\nu x)\\\theta_l} & \mcell{K_l^2} &  &  &  &  &  &  &  &  &  &  &  &  &  &  \\ \hline
8 & \mcell{H_j^2\\a_l(aLu)^3} & \mcell{(\nu\sigma x)\\(j\nu x)} & \mcell{(\theta Hu)^2} & \mcell{K_\eta(KHu)^2\\(j\eta x)\\(aLu)^2} &  &  &  &  &  &  &  &  &  &  &  &  &  &  &  &  &  &  \\ \hline
9 &  & \mcell{H_j\\(aLu)^3} & \mcell{a_\eta(aHu)H_j\\L_{\vartheta}(\theta Lu)^2\\\theta_l(\theta Lu)^2} &  & \mcell{(KHu)} &  &  &  &  &  &  &  &  &  &  &  &  &  &  &  &  &  \\ \hline
10 & \mcell{\theta_l(\theta Lu)^3} & \mcell{L_j\\(j\sigma x)\\a_\eta(aH^2u)^3} &  & \mcell{H_j(aHu)\\(\theta Lu)^2} &  &  &  &  &  &  &  &  &  &  &  &  &  &  &  &  &  &  \\ \hline
11 &  & \mcell{(\theta Lu)^3} & \mcell{K_l(KLu)^3\\L_j\\(aH^3u)^3} &  &  &  &  &  &  &  &  &  &  &  &  &  &  &  &  &  &  &  \\ \hline
12 & \mcell{(aH^2u)^4} & \mcell{a_l(aLu)^2L_j\\H_{\vartheta}(\theta H^2u)^3\\\theta_\eta(\theta H^2u)^3} &  & \mcell{(KLu)^3} &  &  &  &  &  &  &  &  &  &  &  &  &  &  &  &  &  &  \\ \hline
13 & \mcell{(\nu\sigma j)} &  & \mcell{(aLu)^2L_j\\(\theta H^2u)^3} &  &  &  &  &  &  &  &  &  &  &  &  &  &  &  &  &  &  &  \\ \hline
14 & \mcell{(\theta H^2u)^4} & \mcell{K_\eta(KH^2u)^4\\(a,H\!\cdot L,u)^4} &  &  &  &  &  &  &  &  &  &  &  &  &  &  &  &  &  &  &  &  \\ \hline
15 & \mcell{L_\vartheta(\theta,H\!\cdot L,u)^4} &  & \mcell{(KH^2u)^4} &  &  &  &  &  &  &  &  &  &  &  &  &  &  &  &  &  &  &  \\ \hline
16 & \mcell{(\theta,H\!\cdot L,u)^4} &  &  &  &  &  &  &  &  &  &  &  &  &  &  &  &  &  &  &  &  &  \\ \hline
17 & \mcell{K_\eta(K,H\!\cdot L,u)^5} &  &  &  &  &  &  &  &  &  &  &  &  &  &  &  &  &  &  &  &  &  \\ \hline
18 &  & \mcell{(K,H\!\cdot L,u)^5} &  &  &  &  &  &  &  &  &  &  &  &  &  &  &  &  &  &  &  &  \\ \hline
19 &  &  &  &  &  &  &  &  &  &  &  &  &  &  &  &  &  &  &  &  &  &  \\ \hline
20 &  &  &  &  &  &  &  &  &  &  &  &  &  &  &  &  &  &  &  &  &  &  \\ \hline
21 & \mcell{(KH^3u)^6} &  &  &  &  &  &  &  &  &  &  &  &  &  &  &  &  &  &  &  &  &  \\ \hline
\end{tabular}
\end{adjustbox}
\end{center}
\vspace{-0.2em}
\begin{center}\footnotesize (맨 위 가로행의 수들은 $x$에 관한 차수를, 첫 번째 세로행의 수들은 $u$에 관한 차수를 나타낸다.)\end{center}
\endgroup
\end{landscape}

% inherited: noether_paper02_ko_translation_001_20260804\targets\Noether_P02_Korean_T32_U226_UNCHECKED.tex
% sha256: CBBB3A2769D90F15D00A6A22C490178F6D1FE8B5442679D3262BA9DB96FE1D14
\clearpage
\begin{landscape}
\thispagestyle{plain}
\begin{center}\bfseries 표 II (90쪽 참조).\end{center}
\vspace{-0.6em}
\begingroup
\setlength{\tabcolsep}{2pt}
\renewcommand{\arraystretch}{1.35}
\tiny
\begin{center}
\begin{adjustbox}{max width=0.98\linewidth,max totalheight=0.72\textheight,center}
\begin{tabular}{|c|c|c|c|c|c|c|c|c|c|c|c|c|c|c|c|c|c|}
\hline
{} & 0 & 1 & 2 & 3 & 4 & 5 & 6 & 7 & 8 & 9 & 10 & 11 & 12 & 13 & 14 & 15 & 16 \\ \hline
0 & \mcell{J} &  &  &  & \mcell{s} &  & \mcell{s_\vartheta^2} &  &  &  &  & \mcell{s_\nu} & \mcell{(k\nu x)^3s_\nu} & \mcell{(\vartheta^3\nu x)^3s_\nu s_\vartheta\\\theta_\nu(\vartheta^2\nu x)^3s_\vartheta} &  & \mcell{\theta_\eta(\vartheta^2\eta x)^3s_\vartheta} &  \\ \hline
1 &  &  &  &  &  &  & \mcell{(asu)} & \mcell{s_\vartheta} & \mcell{s_k^2\\(\Delta su)} & \mcell{s_\nu(Nsu)\\(\vartheta\nu x)^2s_\nu} & \mcell{((k\nu x)^3,s,u)s_\nu\\(K_\nu(k\nu x)^3,s,u)} &  & \mcell{(K_\eta(k\eta x)^3,s,u)} &  & \mcell{(\vartheta^2\nu x)^3s_\nu} &  & \mcell{((\vartheta^2\eta x)^4,s,u)} \\ \hline
2 &  &  &  &  & \mcell{(asu)^2} & \mcell{s_\vartheta(\theta su)} & \mcell{(\Delta su)^2\\a_\nu s_\nu} & \mcell{((\vartheta\nu x)^2,s,u)s_\nu\\(\theta_\nu(\vartheta\nu x)^2,s,u)} &  & \mcell{s_k\\(\theta_\eta(\vartheta\eta x)^2,s,u)} &  & \mcell{(k\nu x)^2s_\nu\\((k\nu x)^3,s,u)} &  & \mcell{((k\eta x)^3,s,u)} &  &  &  \\ \hline
3 &  &  & \mcell{(asu)^3} & \mcell{s_\vartheta(\theta su)^2\\s_\nu} & \mcell{a_\nu s_\nu(asu)\\a_\nu^2(asu)} & \mcell{s_\eta} & \mcell{(\theta su)\\(a_\eta^2su)} & \mcell{s_k(Ksu)} & \mcell{(Nsu)^2\\(\vartheta\nu x)s_\nu\\((\vartheta\nu x)^2,s,u)} & \mcell{((k\nu x)^3,s,u)^2} & \mcell{((\vartheta\eta x)^2,s,u)} &  &  &  &  &  &  \\ \hline
4 & \mcell{(asu)^4} & \mcell{s_\vartheta(\theta su)^3} & \mcell{a_\nu^2(asu)^2} & \mcell{(\theta_\nu^3su)} & \mcell{(\theta su)^2} & \mcell{s_k(Ksu)^2\\a_\nu(asu)\\s_\nu(asu)} & \mcell{((\vartheta\nu x)^2,s,u)^2\\\theta_\nu s_\nu} & \mcell{s_\nu(\Delta su)\\(\Delta_\nu su)\\(aHu)s_\eta\\(a_\eta su)} &  & \mcell{(\Delta_\eta su)} &  &  &  &  &  &  &  \\ \hline
5 &  &  & \mcell{(\theta su)^3} & \mcell{s_k(Ksu)^3,s_j\\s_\vartheta(\theta s^\prime u)^4\\s_\nu(asu)^2\\a_\nu(asu)^2} & \mcell{(Hsu)\\((\vartheta\nu x)^2,s,u)^3\\(\theta_\nu(\vartheta\nu x),s,u)^2\\(\theta_\nu^2su)} & \mcell{((j\nu x)^3,s,u)\\(aHu)^2s_\eta\\(a_\eta su)^2} & \mcell{(Ksu)^2\\s_l\\(\theta_\eta^2su)} &  & \mcell{((k\nu x)^2,s,u)^2\\K_\nu s_\nu} &  &  &  &  &  &  &  &  \\ \hline
6 & \mcell{(\theta su)^4} & \mcell{s_\nu(asu)^3\\a_\nu(asu)^3} & \mcell{(Hsu)^2\\(\theta_\nu(\vartheta\nu x),s,u)^3\\(\theta_\nu^2su)^2} & \mcell{(a_\eta(aHu),s,u)^2\\(a_\eta(aHu)^2,s,u)} & \mcell{(Ksu)^3} & \mcell{s_j(asu)\\((\vartheta\nu x),s,u)^3\\(\theta_\nu su)} & \mcell{((k\nu x)^2,s,u)^3} & \mcell{((\vartheta\eta x),s,u)\\(\theta Hu)s_\eta\\(\theta_\eta su)} &  &  &  &  &  &  &  &  &  \\ \hline
7 & \mcell{(\theta_\nu(\vartheta\nu x),s,u)^4} & \mcell{(a_\eta(aHu),s,u)^3} & \mcell{(Ksu)^4\\(\theta_\eta^2(\theta Hu),s,u)^2\\(a_\nu^2\cdot a_{\nu_1}^2,s,u)^3} & \mcell{s_j(asu)^2\\s_k(Ks^2u)^3\\((j\nu x),s,u)^3\\(\theta_\nu su)^2} & \mcell{(j\nu x)s_\nu\\((j\nu x)^2,s,u)\\((aHu),s,u)^2\\((aHu)^2,s,u)} & \mcell{((\vartheta\eta x),s,u)^3\\(\theta_\eta su)^2} & \mcell{(aLu)^3s_l\\(asu)^3} &  &  &  &  &  &  &  &  &  &  \\ \hline

8 & \mcell{(a_\nu^2\cdot a_\nu^2,s,u)^4} & \mcell{s_j(asu)^3\\s_k(Ks^2u)^6\\((\vartheta\nu x),s,u)^4\\(\theta_\nu su)^3} & \mcell{((j\nu x)^2,s,u)^2\\((aHu),s,u)^3\\((aHu)^2,s,u)^3} & \mcell{(Lsu)^2\\(a_\nu^2(j\nu x),s,u)^2\\(H_\vartheta(\theta Hu),s,u)^2\\(\theta_\eta(\theta Hu),s,u)^2\\(\theta_\eta(\theta Hu)^2,s,u)} & \mcell{(a_lsu)^3} & \mcell{(K_\nu su)^2} &  & \mcell{(K_\eta su)^2} &  &  &  &  &  &  &  &  &  \\ \hline
9 & \mcell{(K_\nu^2su)^4\\((aHu),s,u)^4} & \mcell{(Lsu)^2\\(a_\nu^2(j\nu x),s,u)^3\\(\theta_\eta su)^4\\(\theta_\eta(\theta Hu),s,u)^3} & \mcell{(Ks^2u)^6\\(a_l(aLu)^2,s,u)^2\\(a_\nu^2\cdot H,s,u)^3} & \mcell{(K_\nu su)^3} & \mcell{(a_\nu(j\nu x),s,u)^2\\((\theta Hu),s,u)^2\\((\theta Hu)^2,s,u)} &  & \mcell{(\theta_lsu)^3} &  &  &  &  &  &  &  &  &  &  \\ \hline
10 &  & \mcell{(K_\nu su)^4\\(a_\nu^2\cdot a_\eta(aHu)^2,s,u)^3} & \mcell{(a_\nu(j\nu x),s,u)^3\\((\theta Hu),s,u)^3\\((\theta Hu)^2,s,u)^2} & \mcell{((j\eta x),s,u)^2\\(H_jsu)\\((aLu)^2,s,u)^2\\((aLu)^3,s,u)} &  &  &  &  &  &  &  &  &  &  &  &  &  \\ \hline
11 & \mcell{(a_\nu(j\nu x),s,u)^4\\((\theta Hu),s,u)^4} & \mcell{(K_\eta(KHu)^2,s,u)^3\\((j\eta x),s,u)^3\\((aLu)^2,s,u)^4\\(a_\nu\cdot H,s,u)^4} & \mcell{(H^2su)^3\\(\theta_l(\theta Lu)^2,s,u)^2} &  & \mcell{((KHu)^2,s,u)^2} &  &  &  &  &  &  &  &  &  &  &  &  \\ \hline
12 &  & \mcell{(a_\eta(aHu)^2\cdot H,s,u)^3} & \mcell{((KHu)^2,s,u)^3} & \mcell{((aHu)H_j,s,u)^2\\((\theta Lu)^2,s,u)^2\\((\theta Lu)^3,s,u)} &  &  &  &  &  &  &  &  &  &  &  &  &  \\ \hline
13 & \mcell{((KHu)^2,s,u)^4} & \mcell{((aHu)H_j,s,u)^3\\((\theta Lu)^2,s,u)^3\\(\theta_\nu\cdot H,s,u)^4} & \mcell{(L_jsu)^2\\((aH^2u)^3,s,u)^2\\((j\nu x)^2\cdot H,s,u)^3\\((aHu)^2\cdot H,s,u)^3} &  &  &  &  &  &  &  &  &  &  &  &  &  &  \\ \hline
14 & \mcell{((j\nu x)^2\cdot H,s,u)^4\\((aHu)^2\cdot H,s,u)^4} & \mcell{(\theta_\eta(\theta Hu)^2\cdot H,s,u)^3} &  & \mcell{((KLu)^3,s,u)^2} &  &  &  &  &  &  &  &  &  &  &  &  &  \\ \hline
15 & \mcell{(a_l(aLu)^2\cdot H,s,u)^4} & \mcell{((KLu)^3,s,u)^3} & \mcell{((aLu)^2L_j,s,u)^2\\((\theta H^2u)^3,s,u)^2\\((\theta Hu)^2\cdot H,s,u)^3} &  &  &  &  &  &  &  &  &  &  &  &  &  &  \\ \hline
16 & \mcell{((\theta Hu)^2\cdot H,s,u)^4\\(a_\nu\cdot H_j,s,u)^4} & \mcell{((j\eta x)\cdot H,s,u)^4\\(H_j\cdot H,s,u)^3\\((aLu)^2\cdot H,s,u)^4\\((aLu)^3\cdot H,s,u)^3} &  &  &  &  &  &  &  &  &  &  &  &  &  &  &  \\ \hline
17 & \mcell{(\theta_l(\theta Lu)^2\cdot H,s,u)^4} &  & \mcell{((KH^3u)^4,s,u)^3} &  &  &  &  &  &  &  &  &  &  &  &  &  &  \\ \hline
18 &  & \mcell{((\theta Lu)^2\cdot H,s,u)^4\\((\theta Lu)^3\cdot H,s,u)^3\\((aHu)^2\cdot H_j,s,u)^3} &  &  &  &  &  &  &  &  &  &  &  &  &  &  &  \\ \hline
19 & \mcell{((aH^3u)^3\cdot H,s,u)^4\\(L_j\cdot H,s,u)^4} &  &  &  &  &  &  &  &  &  &  &  &  &  &  &  &  \\ \hline
20 &  & \mcell{((KLu)^3\cdot H,s,u)^4} & \mcell{((aLu)^3\cdot H_j,s,u)^2} &  &  &  &  &  &  &  &  &  &  &  &  &  &  \\ \hline
21 & \mcell{((\theta H^2u)^3\cdot H,s,u)^4\\((aLu)^2\cdot H_j,s,u)^4} &  &  &  &  &  &  &  &  &  &  &  &  &  &  &  &  \\ \hline
22 &  &  &  &  &  &  &  &  &  &  &  &  &  &  &  &  &  \\ \hline
23 & \mcell{((KH^2u)^4\cdot H,s,u)^4} & \mcell{((aH^2u)^3\cdot H,s,u)^3} &  &  &  &  &  &  &  &  &  &  &  &  &  &  &  \\ \hline
\end{tabular}
\end{adjustbox}
\end{center}
\vspace{-0.2em}
\begin{center}\footnotesize (맨 위 가로행의 수들은 $x$에 관한 차수를, 첫 번째 세로행의 수들은 $u$에 관한 차수를 나타낸다.)\end{center}
\endgroup
\end{landscape}
